\documentclass[10pt,a4paper]{article}
\newcommand{\ReportShortTitle}{Bayesian approximation algebra}
\usepackage[T1]{fontenc}
\usepackage[utf8]{inputenc}
\usepackage{lmodern}
\usepackage{microtype}
\usepackage[a4paper,margin=24mm,headheight=14pt]{geometry}
\usepackage{amsmath,amssymb,mathtools}
\usepackage{booktabs,tabularx,array}
\usepackage{enumitem}
\usepackage{xcolor}
\usepackage{fancyhdr}
\usepackage{titlesec}
\usepackage[hidelinks]{hyperref}

\definecolor{ink}{HTML}{18212B}
\definecolor{accent}{HTML}{135F73}
\definecolor{muted}{HTML}{5E6A73}
\definecolor{pale}{HTML}{EEF5F7}
\color{ink}
\hypersetup{colorlinks=true,linkcolor=accent,urlcolor=accent}
\setlength{\parindent}{0pt}
\setlength{\parskip}{0.55em}
\setlist{nosep,leftmargin=1.55em}
\setcounter{secnumdepth}{2}
\titleformat{\section}{\Large\bfseries\color{accent}}{\thesection}{0.6em}{}
\titleformat{\subsection}{\large\bfseries\color{ink}}{\thesubsection}{0.55em}{}
\pagestyle{fancy}
\fancyhf{}
\fancyhead[L]{\small ODMA--URA research record}
\fancyhead[R]{\small\color{muted}\ReportShortTitle}
\fancyfoot[C]{\thepage}
\renewcommand{\headrulewidth}{0.3pt}
\newcolumntype{Y}{>{\raggedright\arraybackslash}X}
\newcommand{\R}{\mathbb{R}}
\newcommand{\C}{\mathbb{C}}
\newcommand{\Z}{\mathbb{Z}}
\newcommand{\F}{\mathbb{F}}
\newcommand{\one}{\mathbf{1}}
\newcommand{\given}{\,|\,}
\newcommand{\CN}{\mathcal{CN}}
\DeclareMathOperator{\diag}{diag}
\DeclareMathOperator{\Cov}{Cov}
\DeclareMathOperator{\Var}{Var}
\DeclareMathOperator{\supp}{supp}
\DeclareMathOperator*{\argminop}{arg\,min}
\DeclareMathOperator*{\argmaxop}{arg\,max}
\newcommand{\evidence}[1]{\textcolor{accent}{\textbf{Evidence status:}} #1}
\providecommand{\ReportUpdated}{30 August 2026}
\newcommand{\ProjectReport}[4]{%
  \pagenumbering{roman}
  \begin{titlepage}
    \vspace*{18mm}
    {\large\color{accent}\textbf{ODMA--URA research record}}\\[12mm]
    {\huge\bfseries #1}\\[5mm]
    {\Large\color{muted}#2}\\[18mm]
    \begin{tabularx}{\textwidth}{@{}p{0.16\textwidth}Y@{}}
      Stage & #3\\
      Updated & \ReportUpdated\\
      Status & #4
    \end{tabularx}
    \vfill
    \textcolor{muted}{This report is a chronological research record. Measured results, interpretations, and open
    hypotheses are labelled separately. Detailed numerical ledgers remain in \texttt{results/}.}
  \end{titlepage}
  \pagenumbering{arabic}
}

\newcommand{\E}{\mathbb E}
\newcommand{\ind}{\mathbf 1}
\newcommand{\TV}{\operatorname{TV}}
\begin{document}
{\LARGE\bfseries Polar Bayesian decoding: approximation algebra}\par
{\large Exact bottlenecks, explicit assumptions and decision-error control}\par
{\color{muted}19 October 2026 --- Report 9; theoretical completion audit before implementation}

\textbf{Research direction.} Retain the Bayesian decoder objective from Report 7 and investigate the expensive
operations in its different exact evaluations. An assumption must follow from the algebra of an identified
operation, explain what computation it could remove, and expose the information and accuracy it would lose.
Exact discoveries are resources for this investigation; they do not commit us to one receiver architecture.

\textbf{Scope.} We fix the target, map the exact computational bottlenecks and derive shared error-control
statements. The entries develop certified restriction, continuation quantization, sparse regional interactions,
positive features, connected-component truncation, residual moment closure, marked Fourier approximations,
weighted sampling, prefix queries, cyclic feature contraction, weighted correction bounds, residual-guided
proposals and bounded multiplicity reconstruction. Further entries develop tilted residual laws, positive Gaussian
auxiliary fields, consistent moment bounds, certified polar inner queries, variational references and signed
compression. A final audit closes the finite constructions and separates their conditions from unproved
useful scaling. The existing decoder compositions remain options rather than a selected receiver. Each entry has
a specified mathematical replacement, a conditional accuracy statement and a cost account.
None is established as a scalable decoder for this encoder, and previous experimental rankings are not carried forward. There are
no implementations, simulation panels or empirical PUPE results in this report. Report 8 supplies the fixed
encoder; Report 7 retains exact identities, reductions, obstructions and certified decisions.

\section{Keep the posterior and Bayes decision fixed}
Let $\mathcal A$ be the legal message-count space. The receiver knows $K$, the real AWGN variance $\sigma^2$
and the encoding rule, but neither messages nor patterns. With the fixed balanced polar/CRC encoder,
\[
 \phi(p,v)=S_p\diag(\sqrt{e_p})(-1)^{Lv},\qquad
 y=\Phi a+z,\quad z\sim\mathcal N(0,\sigma^2I_n),\quad \|\phi(w)\|^2=1.
\]
Here $P=2^{B_p}$ patterns, $V=2^{B_c}$ payloads per pattern and $M=PV=2^B$ complete messages.
$L$ is the $N_c$-by-$B_c$ binary CRC/polar generator. The matrix $\Phi$ is mathematical notation, not a
requirement to materialize the full codebook. The IID source is primary:
\[
 \mathcal A=\{a\in\mathbb Z_+^M:\one^Ta=K\},\qquad
 \pi(a)=\frac{K!}{M^K\prod_w a_w!}.
\]
The distinct-message control restricts $a_w$ to $0/1$ with a uniform legal prior. It retains same-pattern
collisions between different payloads. An inference simplification must not silently change this source law.
Define
\begin{equation}
 \omega(a)=\pi(a)e^{-\|y-\Phi a\|^2/(2\sigma^2)},\quad
 Z=\sum_{a\in\mathcal A}\omega(a),\quad P_y(a)=\omega(a)/Z.
 \label{eq:target}
\end{equation}
The loss and its required statistics are
\begin{equation}
 \ell(a,d)=1-\frac1K\sum_w\min(a_w,d_w),\qquad
 T_{w,j}=P_y(a_w\geq j),\quad \one^Td\leq K.
 \label{eq:loss-target}
\end{equation}
The exact Bayes action allocates $K$ occurrences by the largest available next count tails. Approximation
changes how these quantities are obtained. A mode, a small signal residual or marginal bit confidence alone
does not replace the Bayes objective. No MAP/least-squares alternative is an active route in this study.

\section{Map each exact calculation to its expensive object}
The following are representations of the same posterior and decision, not different statistical targets.
Let $H(L,h)$ denote the number of legal count configurations at occupancy $h$ on an alphabet of size $L$:
\[
 H_{\rm IID}(L,h)=\binom{L+h-1}{h},\qquad H_{\rm distinct}(L,h)=\binom Lh.
\]
For fixed $K$ and large $M$, both scale as $M^K/K!$. Polynomial cost in $M$ is still exponential in $B$.

\begin{center}\small
\begin{tabularx}{\textwidth}{@{}>{\raggedright\arraybackslash}p{.22\textwidth}YY@{}}\toprule
Exact evaluation & Expensive object & Exact resources to examine before simplifying\\\midrule
Count posterior / residual recursion & $H(M,K)$ configurations or $M^K$ ordered branches; each residual needs continuation evidence & Factorial source weights, conditional polar sums, symmetry and admissible completion bounds\\
Inside/outside / compatible frontier & Generated transitions and reachable continuation states, not just final table size & Character keys, future projection, finalized likelihood directions and marked backward queries\\
Polar pair / three-user contractions & Explicit payload/pattern maps and tables; joint interfaces at several neighbours & XOR convolution, overlap-parity quotients, positive K2/K3 contractions and exact reference changes\\
Bit/pattern factorization / elimination & Joint intermediate separator factors and complete-message count observables & Exact XOR/CRC factors, shared pattern variables and conditional separation\\
Finite Fourier evaluation & $(2K+1)^d$ modes for $d$ integer-signature dimensions, plus marked coefficient access & Convolution powers, source-specific coefficient identities and omitted-mode bounds\\
Connected-user expansion & Connected coefficients through order $K$, their marked contractions and signed cancellation & Pair/triangle subcalculations, exact reference changes, rooted-tree remainder bounds\\
Certified Bayes search & Search nodes, likelihood envelopes and unresolved decision contrasts & Support duals, cached directions and adverse-event prior masses\\\bottomrule
\end{tabularx}
\end{center}
A proposed reduction must identify its place in this map. Removing one expensive object can expose another.
For example, a compact interaction table does not make the sum over unknown patterns free, and a small
separator dimension does not imply a small number of states or small exponential-kernel rank.

\subsection{Parameterize cost before invoking a favourable regime}
Let $S,T$ count constructed states and transitions, and let $E$ count intersecting pattern pairs. Define
\[
 J_{\rm pair}=\sum_{p<q,\,\mathrm{overlap}}2^{r_{pq}+r_{qp}}
\]
as the number of cross-pattern parity-table entries. Existing exact schedules include
\[
 C_{\rm score}=O(PN_c+M\log V),\qquad
 C_{\rm pair}=O(PV\log V+EV+J_{\rm pair}),
\]
where endpoint maps and their construction remain charged. The complete K3 contraction uses $M$ pair
applications plus row extraction and scalar work, with small additional streaming memory. It is not a
constant number of pair applications for arbitrary $K$.

Inside/frontier cost follows $S,T$ and the cost of constructing/projecting keys. A scalar reverse pass can
return all marked counts from the recorded arithmetic without $M$ complete forward reruns, but cannot
remove forward state growth. Exact elimination follows joint domain products in intermediate bags.
Connected remainder bounds can be cheaper than connected coefficients while still being too loose to
resolve a decision. Every proposed favourable parameter regime must address these distinctions.

\section{How an assumption enters this report}
Each retained entry is assessed against the following requirements:
\begin{enumerate}
\item \textbf{Exact origin.} State the exact equation and costly object being replaced, including its variables,
      legal support and any conditioning. Cite the relevant exact result rather than restarting the model.
\item \textbf{Algebraic motivation.} Remove available exact redundancy and identify a precise property that could
      make the remaining object simpler. Name the regime and parameter in which that property is assumed.
\item \textbf{Replacement.} Define the new weights, factor, state, expectation or contraction. Say how normalizers,
      source exclusions, cardinality and full-message/count identity are retained.
\item \textbf{Accuracy statement.} Derive a bound or exact error identity, state the assumptions needed for it,
      and connect it to count tails or competing Bayes allocations. If no useful bound is available, state
      exactly what remains unproved. Include a limiting case or algebraic counterexample.
\item \textbf{Complete computational effect.} Count construction, pattern search, inference, readout, memory,
      precision and error-control work as $B_p,B_c,K,n,N_c$ grow. A small retained object is insufficient if
      obtaining it requires the discarded exponential calculation.
\item \textbf{Dependencies.} State what can be composed, which constraints must survive and whether another
      assumption is needed. A result conditional on a supplied good list or true pattern is not a complete receiver.
\end{enumerate}
A condition can be mathematically meaningful before we know whether it is affordable to establish. Keep its
logical implication separate from the claim that the actual polar family satisfies it. An unavailable
certificate is an open computational problem, not permission to substitute truth or exact posterior quantities.

No decoder engineering, tuning or new testing is part of the present stage. An entry enters the main text
when its algebra and limitations are clear; the working log records unresolved derivations. Many assumptions
may eventually be investigated, but a long menu of named algorithms is not the unit of this record.

\section{Positive weight errors give posterior-error bounds}
\label{sec:positive-errors}
This is a reusable statement about finite distributions, not an assumption that the physical posterior has
any simplifying structure. Let $\widetilde\omega(a)>0$ on the same legal support and $Q(a)=\widetilde\omega(a)/\widetilde Z$.
Suppose the log-ratio range is bounded:
\[
 \log\frac{\omega(a)}{\widetilde\omega(a)}\in[c,c+\Omega]\quad\text{for all legal }a.
\]
Then
\begin{equation}
 \TV(P_y,Q)\leq\tanh(\Omega/4).
 \label{eq:tv-range}
\end{equation}
A common multiplicative error changes no normalized probability; the range matters rather than an arbitrary
absolute offset. To prove the bound, put $R=\omega/\widetilde\omega\in[u,v]$ and $\mu=\E_QR$. Convexity of
$|R-\mu|$ bounds it by its endpoint chord, giving
\[
 \TV(P_y,Q)\leq\frac{(v-\mu)(\mu-u)}{(v-u)\mu}
 \leq\frac{\sqrt v-\sqrt u}{\sqrt v+\sqrt u}
 =\tanh\left(\frac14\log\frac vu\right).
\]
The maximum occurs at $\mu=\sqrt{uv}$ and is attainable with endpoint-supported ratios. If $u=v$, the
posterior is unchanged and the bound follows by continuity.

For example, a uniform absolute log-weight error at most $e$ has range at most $2e$, so
$\TV(P_y,Q)\leq\tanh(e/2)$. Positivity and common legal support are conditions of this result.
Hard source constraints must remain hard; signed truncated coefficients do not define $Q$ automatically.

\subsection{Compatible local modifications}
If the same configuration weight is changed through factors with log-ratio ranges $\Omega_1,\ldots,\Omega_m$,
then the total range is at most their sum:
\[
 \TV(P_y,Q)\leq\tanh\left(\frac14\sum_{i=1}^m\Omega_i\right).
\]
This is a deterministic bound and needs no independence of errors. The factors must belong to one compatible
representation of the original configurations. Applying the full Gaussian afresh at several cuts would
reuse the observation and would not meet this condition.

\clearpage
\section{Support restriction has a separate probability cost}
\label{sec:restriction-foundation}
For a retained event $E$ with $P_y(E)>0$, put $\eta=P_y(E^c)$ and $P_E=P_y(\cdot\mid E)$.
A candidate-word restriction means every transmitted copy lies in the proposed alphabet; a configuration
restriction means the entire count vector belongs to a retained set. These events are different.
Exact conditioning gives
\begin{equation}
 \TV(P_y,P_E)=\eta,\qquad D(P_E\|P_y)=-\log(1-\eta).
 \label{eq:restriction}
\end{equation}
More generally, any $Q$ supported on $E$ satisfies
\[
 D(Q\|P_y)=D(Q\|P_E)-\log P_y(E).
\]
Thus exact conditioning is the reverse-KL optimum within that support, although its joint inference can
remain expensive. If $\TV(P_E,Q)\leq\tau$, then
\begin{equation}
 \TV(P_y,Q)\leq\eta+(1-\eta)\tau.
 \label{eq:restriction-plus}
\end{equation}
Indeed, $P_y=(1-\eta)P_E+\eta P_{E^c}$, and convexity gives the stated bound. For a candidate alphabet $C$,
any action supported on $C$ also has unavoidable posterior expected missed-copy contribution
\[
 \rho_C=\frac1K\sum_{w\notin C}\E_{P_y}a_w\leq\eta.
\]
This is an absolute risk contribution, not by itself a lower bound on excess risk relative to Bayes.

Posterior omitted mass, prior event probability and empirical truth recall are distinct quantities. A
restriction must supply a way to control $\eta$, or explicitly leave that control as an assumption. Computing
$\eta$ with the full exact posterior would not establish a tractable restriction procedure.

\section{Transfer accuracy to the Bayes allocation}
\label{sec:decision-foundation}
For any action $d$ with total $K$, its posterior expected matching gain is
\[
 G_T(d)=\sum_w\sum_{j=1}^{d_w}T_{w,j},\qquad R_T(d)=1-G_T(d)/K.
\]
Suppose approximate scores $\widetilde T_{w,j}$ have uniform error at most $e$, and $\widetilde d$ maximizes
$G_{\widetilde T}$ over feasible actions. With $d^*$ the true Bayes action,
\begin{align}
 K\big[R_T(\widetilde d)-R_T(d^*)\big]
 &=G_T(d^*)-G_T(\widetilde d)\\
 &\leq |G_T(d^*)-G_{\widetilde T}(d^*)|
       +|G_{\widetilde T}(\widetilde d)-G_T(\widetilde d)|
 \leq 2Ke.
\end{align}
Therefore
\begin{equation}
 0\leq R_T(\widetilde d)-R_T(d^*)\leq\min(1,2e).
 \label{eq:decision-error}
\end{equation}
If the scores are count tails of a positive approximate law with TV error $\tau$, then $e\leq\tau$.
Combining with \eqref{eq:restriction-plus} gives the bound
$\min\{1,2[\eta+(1-\eta)\tau]\}$ for that law's Bayes action. A direct tail bound can be more informative
than a uniform full-posterior bound. If approximate scores are not monotone count tails, simple greedy
allocation need not maximize their feasible objective; the optimization condition above must still hold.
One repair is to sort each word's $K$ estimated tails in nonincreasing order after clipping to $[0,1]$.
Sorting cannot increase their uniform distance from the true ordered tails: coordinatewise errors at
most $e$ imply the same bound for every order statistic. The sorted scores therefore retain the error
bound and permit greedy allocation. They define per-word count laws, without asserting a consistent
joint posterior. This costs $O(K\log K)$ per explicitly queried word.

\clearpage
\subsection{Exact decisions can need less than an accurate full posterior}
For monotone approximate tails and their greedy action, let $g$ be the smallest last-selected versus
next-unselected tail gap across different words. If $g>2e$, every required ordering remains valid under the
true tails, certifying a true Bayes action. Same-word prefix order is retained by the exact count-tail law.
A tie gives no such strict certificate.

More directly, let $H=\ind\{a_u\geq j\}-\ind\{a_v\geq l\}$. On common legal support, Report 7's exact
reweighting identity gives
\[
 \E_{P_y}H-\E_QH=\frac{\Cov_Q(H,R)}{\E_QR},\qquad R=\omega/\widetilde\omega.
\]
A comparison-specific bound can exploit cancellations that separate tail bounds lose. However, its expectation
may be as hard to obtain as the original inference. Multiple corrections use updated references; their
initial covariance effects are not generally additive. A useful assumption must address both accuracy and cost.

\section{A geometric error resource for continuation calculations}
\label{sec:width-foundation}
At one exact alphabet/cardinality cut, let $s$ be a partial signal and $t\in\mathcal T$ a legal future sum.
Compare it with $\bar s$, write $\delta=s-\bar s$, and set $g_s(t)=e^{-\|y-s-t\|^2/(2\sigma^2)}$. Then
\[
 \log\frac{g_s(t)}{g_{\bar s}(t)}
 =c_s-\frac{\delta^Tt}{\sigma^2},\qquad
 c_s=\frac{\delta^T(y-\bar s)-\|\delta\|^2/2}{\sigma^2}.
\]
Define support $h_{\mathcal T}(d)=\max_{t\in\mathcal T}d^Tt$, directional width
$w_{\mathcal T}(d)=h_{\mathcal T}(d)+h_{\mathcal T}(-d)$ and midpoint
$m_\delta=[h_{\mathcal T}(\delta)-h_{\mathcal T}(-\delta)]/2$. Absorbing the scalar
$e^{c_s-m_\delta/\sigma^2}$ gives
\begin{equation}
 \left|\log\frac{g_s(t)}{e^{c_s-m_\delta/\sigma^2}g_{\bar s}(t)}\right|
 \leq\frac{w_{\mathcal T}(\delta)}{2\sigma^2}.
 \label{eq:width}
\end{equation}
The midpoint minimizes the worst absolute log error among scalar rescalings. Width zero recovers Report 7's
lossless proportional-row merging criterion. This is the algebraic reason future dependence, rather than
Euclidean signal distance alone, is a relevant approximation quantity.

For $h$ remaining IID messages, the width is $h$ times the largest-minus-smallest allowed-word score
$\delta^T\phi(w)$. For distinct messages it is top-$h$ minus bottom-$h$ sums. Upper support bounds in both
directions supply a safe wider interval; incomplete candidate maxima do not. States being compared must
have compatible legal futures, and their marked count weights must be retained. Multiple cuts require the
compatible likelihood schedule of Report 7.

This is an error-control resource, not an assumption that a small number of useful groups exists. A future
entry must separately derive how to construct them, their widths, their generated-state cost and any pruned
posterior mass. No grouping algorithm or favourable width on the research bank is asserted here.

\section{Signed calculations need numerator and normalizer control}
\label{sec:signed-foundation}
Fourier and connected representations can be exact while partial sums are signed. Let $N$ be an exact
count-tail numerator and $Z$ its positive evidence, with real estimates $\widetilde N,\widetilde Z$ satisfying
$|N-\widetilde N|\leq\delta_N$ and $|Z-\widetilde Z|\leq\delta_Z$. If $\widetilde Z>\delta_Z$, then
\begin{equation}
 \left|\frac NZ-\frac{\widetilde N}{\widetilde Z}\right|
 \leq\frac{\delta_N+|\widetilde N/\widetilde Z|\delta_Z}
             {\widetilde Z-\delta_Z}.
 \label{eq:ratio-error}
\end{equation}
This follows by adding and subtracting $\widetilde N/Z$ and using $Z\geq\widetilde Z-\delta_Z$.
It shows why small absolute evidence error need not give useful posterior accuracy when evidence is small.
The numerator/remainder bounds must also be computable without the full discarded calculation.

A signed approximation is not automatically a positive approximate posterior. If rigorous intervals enclose
the true tails, they can support exact decision comparisons without treating the partial sum as a probability
law. Clipping coefficients into a posterior would be an additional transformation requiring its own argument.
Source legality, normalization and arithmetic cancellation are separate issues.

\section{One invariant links two apparent approximation problems}
\label{sec:common-interaction}
Before choosing a simplification, remove dependence that is only a choice of reference. For the positive
cross factor $b(s,t)=e^{-s^Tt/\sigma^2}$, Report 7 gives the four-point identity
\begin{equation}
 \log\frac{b(s,t)b(s',t')}{b(s,t')b(s',t)}
 =-\frac{(s-s')^T(t-t')}{\sigma^2}.
 \label{eq:four-point}
\end{equation}
Multiplying rows or columns by positive scalars changes none of these values. The same identity applies to
Gaussian continuation rows, because their remaining terms depend on only $s$ or only $t$. Thus the part of
past-state variation that matters is precisely the part that interacts with future variation. For regions,
it is the variation in one region that interacts with variation in the other.

Here is a useful obstruction with a narrow claim. On a feasible rectangular domain $\mathcal S\times\mathcal T$,
let $\chi$ be the largest absolute value in \eqref{eq:four-point}. Suppose we replace $b(s,t)$ by
$u(s)v(t)>0$ and the residual log error $e(s,t)=\log b(s,t)-\log u(s)-\log v(t)$ has range length $\Omega$.
Every four-point expression is a sum of two errors minus two others, so
\begin{equation}
 \Omega\geq\chi/2,\qquad \sup_{s,t}|e(s,t)|\geq\chi/4.
 \label{eq:rank-one-obstruction}
\end{equation}
No unary rescaling can make a strongly varying rectangular interaction uniformly separable. Conversely,
choosing anchors $s_0,t_0$ and the positive replacement
\[
 \widetilde b(s,t)=\frac{b(s,t_0)b(s_0,t)}{b(s_0,t_0)}
\]
leaves exactly $e(s,t)=-(s-s_0)^T(t-t_0)/\sigma^2$. Bounds on this bilinear remainder give sufficient error
conditions directly. If $\chi=0$, the replacement is exact.

This does not prove that posterior or decision error must be large when $\chi$ is large. The troublesome
states may have little mass, or their errors may cancel in a particular decision contrast. It also does not
exclude a sum of several separable factors. Hard source zeros and nonrectangular feasible domains must be
retained separately. In particular, cardinality identities can cancel errors across factors even when one
factor alone is not separable. The following constructions investigate several ways to reduce computation
around this structure, including positive sums of separable factors.

\section{Entry A: restrict hypotheses with an evidence bound}
\label{sec:entry-restriction}
\subsection{The exact operation and proposed assumption}
Count enumeration and residual search sum positive weights over $H(M,K)$ legal states. Let $E$ retain some
of those states. The proposed structural assumption is that a set with large posterior mass can be
\emph{found and certified cheaply}. Small retained size alone is not this assumption.
If the retained weights are evaluated exactly, the replacement is the conditional law $P_E$, with
$\TV(P_y,P_E)=\eta=P_y(E^c)$ and Bayes excess loss at most $\min(1,2\eta)$.

There is a non-oracle route to an upper bound on $\eta$. Cover all omitted states by disjoint search subtrees
$b$, each with a valid unnormalized evidence upper bound $U_b$. Let $Z_E$ be the retained evidence and
$U_D=\sum_b U_b$. Then
\begin{equation}
 \eta=\frac{Z_D}{Z_E+Z_D}\leq\frac{U_D}{Z_E+U_D}.
 \label{eq:discard-mass}
\end{equation}
A certified positive lower bound $L_E\leq Z_E$ can replace $Z_E$ in the right side. All weights must use the
same prior normalization. An overlapping frontier must first be made disjoint or treated with a valid upper
bound, rather than silently counted as an exact partition.

Report 7 supplies bounds without knowing the transmitted words. For a count prefix with fixed counts
$c_i$, residual $r$, $h$ remaining users and $m$ allowed remaining words, the unnormalized factorial-prior
mass is $m^h/(h!\prod_i c_i!)$ under IID, or $\binom mh$ under distinct sampling. Multiplying it by
$\exp[-\max(0,\|r\|-h)^2/(2\sigma^2)]$ bounds its evidence. Directional support bounds can tighten the
distance term. Finding useful directions and bounding their maxima remain charged computation.

\subsection{Why a likelihood gap alone is insufficient}
Consider one retained explanation with weight $1$ and $N$ omitted explanations each with weight $e^{-\Delta}$.
Their total posterior mass is
\[
 \eta=\frac{N e^{-\Delta}}{1+N e^{-\Delta}}.
\]
To make this at most $\varepsilon$ requires
$\Delta\geq\log N+\log[(1-\varepsilon)/\varepsilon]$. Individually poor explanations can collectively
dominate. A residual threshold or a large best-versus-second-best gap does not control the omitted sum.
The factorial source weights, pattern counts and occupancy multiplicities belong in that sum.

\subsection{Computational effect and precise boundary}
Let $F$ be the number of retained complete count states, $J$ the number of constructed search nodes and
$C_{\rm node}$ the per-node construction/bounding cost. With sparse count lists, a direct calculation costs
\[
 O\!\left(JC_{\rm node}+F(n+K)+FK\log(2+FK)\right)
\]
including likelihood evaluation from maintained residuals, accumulation of occupied count-tail events and sorting the scores.
The union of those events has at most $FK$ entries. This is a conservative arithmetic count, with search
storage additional to the $O(FK)$ retained sparse record. Waveform construction and incremental residual
updates are included in $C_{\rm node}$; constructing each leaf sum independently would instead cost $O(nK)$ per leaf.
The full $H(M,K)$ enumeration disappears only when $J$, $F$ and $C_{\rm node}$ are themselves affordable.
Generating each candidate by scanning all $M$ words is still exponential in $B$.

This entry applies to count/residual search, pruned inside/outside paths and certified Bayes search. A
candidate-word set $C$ induces a different retained event, namely $\supp(a)\subseteq C$; one must bound
all excluded configurations, not just missing individual words. If approximate inference is then used on
$E$, Section~\ref{sec:restriction-foundation} combines its error with the omitted mass. No new physical
independence assumption is needed, but cheap evidence certification remains unproved for the target bank.

\clearpage
\section{Entry B: quantize only the dependence on future messages}
\label{sec:entry-continuation}
\subsection{One cut: define a positive replacement before discussing its size}
At a disjoint alphabet cut and fixed occupancy, write an exact configuration weight as
$\alpha(a_A)\beta(a_B)g_s(t)$, with $s=\Phi_Aa_A$, $t=\Phi_Ba_B$ and the exact source factors in
$\alpha,\beta$. Group partial histories with the same legal future set. For a representative $\bar s_g$,
put $\delta_s=s-\bar s_g$ and define
\[
 \kappa_s=\exp(c_s-m_{\delta_s}/\sigma^2),\qquad
 \widetilde\omega(a_A,a_B)=\alpha(a_A)\beta(a_B)\kappa_s g_{\bar s_g}(t).
\]
Section~\ref{sec:width-foundation} proves a pointwise log error at most
$\varepsilon_g$ whenever $w_{\mathcal T}(\delta_s)/(2\sigma^2)\leq\varepsilon_g$ for all histories in
group $g$. The grouped weight is $\sum_{a_A\in g}\alpha(a_A)\kappa_s$; its marked version must retain each
past count event. All legal count configurations still exist in the approximate law. The representative
waveform need not be a legal transmitted sum: it represents a likelihood row, not a proposed message identity.

For one cut with $S$ past states and $T$ future states, $R$ representatives replace the dense $ST$ row
contraction by $RT$, after constructing and marking the groups. Enumerating all $S$ states first retains
that cost. The substantive assumption is a small, constructible cover in the seminorm
\begin{equation}
 d_{\mathcal T}(s,s')=\frac{w_{\mathcal T}(s-s')}{2\sigma^2}.
 \label{eq:future-seminorm}
\end{equation}
Its zero directions are exactly the lossless projection directions already found in Report 7.

\subsection{A constructive covering bound exposes the remaining exponential}
Remove those zero directions with an orthonormal basis $Q$. Write $u=Q^Ts$ and $z=Q^Tt$ in dimension $d$.
Suppose the retained past domain has coordinate ranges $A_j$ and the future domain has coordinate ranges
$D_j$. For any past displacement,
\[
 w_{\mathcal T}(s-\bar s)\leq\sum_{j=1}^d D_j|u_j-\bar u_j|.
\]
Grid cells of width $h_j$ with midpoint representatives therefore have error at most
$\sum_jD_jh_j/(4\sigma^2)$. With all $D_j>0$, choosing
$h_j=4\sigma^2\varepsilon/(dD_j)$ gives the sufficient cover size
\begin{equation}
 R\leq\prod_{j=1}^d\max\{1,\lceil A_j/h_j\rceil\}
 \leq\prod_{j=1}^d\left(1+\frac{dA_jD_j}{4\sigma^2\varepsilon}\right).
 \label{eq:cover-size}
\end{equation}
Coordinates with zero range need no subdivision. A nonrectangular attainable domain can use fewer occupied
cells, but finding them must also be affordable. Support upper bounds give safe coordinate ranges without
enumerating completions; sharp ranges are not assumed free. Maximizing and minimizing the linear score over
this future-coordinate box also gives a valid interval for constructing $\kappa_s$, with the same stated bound.

This is an upper bound for a specified construction, not a theorem that this many states are necessary.
It makes the difficulty visible: a fixed small $d$ can turn state growth into a controlled power of the
accuracy requirement, while large $d$ can retain exponential cost. Low physical noise also demands finer
cells. Exact projection alone provides no small-cover guarantee for the selected polar bank.

\subsection{A complete sequential construction, with telescoping errors}
One-cut accuracy is insufficient if a recursive algorithm repeatedly merges states. The following explicit
construction shows what a compatible composition must mean. Order the complete alphabet as $1,\ldots,M$.
At step $i$, each legal path chooses its exact count $a_i$; preserve the running total $k$ and source factor
$1/a_i!$ (or $a_i\in\{0,1\}$ for distinct sampling). Let $\bar s_0=0$ and set
\[
 x_i=\bar s_{i-1}+a_i\phi_i,\qquad
 \bar s_i=\text{representative of }x_i\text{ at cut }(i,k).
\]
The remaining legal sum set depends only on $(i,K-k)$, so merged paths have compatible futures. Apply the
one-cut rule to $x_i,\bar s_i$, obtaining a positive scalar $\kappa_i$ and an error bound $\varepsilon_i$.
Cover ranges must include the incoming $x_i$ of this approximate recurrence, not only exact partial sums;
previous quantization can enlarge those ranges. No out-of-range transition may be silently discarded.
The resulting complete-path weight is
\begin{equation}
 \widetilde\omega(a)=\beta(a)\left(\prod_{i=1}^M\kappa_i\right)g_{\bar s_M}(0),
 \qquad \beta(a)=\prod_i1/a_i!\quad\text{(IID)}.
 \label{eq:quantized-path}
\end{equation}
The omitted source normalization is common to all paths. No path is deleted, and every legal count vector
has one path. Aggregating paths into the same state therefore computes a positive law on the original support.

To check the error, fix a path and let $t_i=\sum_{j>i}a_j\phi_j$. The identities
$g_{\bar s_{i-1}}(a_i\phi_i+t_i)=g_{x_i}(t_i)$ make the ratios telescope:
\begin{equation}
 \frac{g_0(\sum_i a_i\phi_i)}{(\prod_i\kappa_i)g_{\bar s_M}(0)}
 =\prod_i\frac{g_{x_i}(t_i)}{\kappa_i g_{\bar s_i}(t_i)}.
 \label{eq:telescoping-quantization}
\end{equation}
Thus if $\sum_i\varepsilon_i(a)\leq E_*$ for every legal path, then
\begin{equation}
 \TV(P_y,Q)\leq\tanh(E_*/2),\qquad
 R(\widetilde d)-R(d^*)\leq\min\{1,2\tanh(E_*/2)\}.
 \label{eq:quantization-error}
\end{equation}
This consumes the Gaussian exactly once through a telescoping identity; it does not multiply a new
independent observation likelihood at every step. Taking the largest error at each cut and summing gives
a conservative $E_*$. Path-specific bounds can be retained instead, at additional bookkeeping cost.

\subsection{What state merging saves, including full count readout}
If at most $R$ representatives are constructed per cut and occupancy, there are $O(MKR)$ states and
$O(MK^2R)$ count transitions. Let $C_{\rm step}$ include signal/key updates, grid assignment and the
certified support interval needed for $\kappa_i$. Forward arithmetic costs
$O(MK^2R C_{\rm step})$, in addition to constructing the encoding and cover bounds.
One always available interval is $\delta_i^Tt_i\in[-h\|\delta_i\|,h\|\delta_i\|]$ for $h$ remaining
unit-energy messages. Its midpoint is zero, and it gives local error at most $h\|\delta_i\|/\sigma^2$
with $O(n)$ arithmetic. This cheap bound may need many more representatives; obtaining sharper polar
support intervals can itself require explicit-alphabet work. Neither accuracy nor certificate cost is free.

All count probabilities can be recovered on the recorded graph: attach a formal multiplier $z_{i,c}$ to
every transition choosing count $c$ at step $i$. At $z=1$,
\[
 \frac{1}{\widetilde Z}\frac{\partial\widetilde Z}{\partial\log z_{i,c}}
 =Q(a_i=c).
\]
Each complete path chooses precisely one such transition at that step. A reverse pass over the same scalar
sum/product graph obtains all these derivatives with arithmetic linear in its recorded size; count tails
then cost $O(MK)$. Hold the constructed groups and scalars fixed when differentiating. This does not
differentiate through a clustering rule or replace multiplicity by a slot marginal.

Forward working memory can be $O(nKR)$ numbers for full signal records; the simple all-query reverse calculation
stores $O(MK^2R)$ transitions and $O(nMKR)$ signal entries unless recomputation or projection is used. With $R$ small, the
combinatorial $H(M,K)$ dependence has been replaced by explicit $M$ times a controlled state count.
It is still exponential in $B$. Furthermore, an equal error budget over $M$ cuts requires
$\varepsilon\simeq E_*/M$, worsening \eqref{eq:cover-size}; a small one-cut cover does not establish a cheap
whole decoder. Prefix-level construction could avoid visiting all $M$ leaves only after its own legal
transitions, covers and full-message readout are derived. No such sub-exhaustive construction is claimed here.

\subsection{A stronger schedule: charge only occupied-word transitions}
\label{sec:occupied-quantization}
The factor $M$ in the preceding \emph{error budget} is avoidable. A legal path visits $M$ alphabet positions
but chooses a positive count at at most $K$ of them. Make every zero-count transition the identity:
\[
 a_i=0:\quad \bar s_i=\bar s_{i-1},\qquad \kappa_i=1.
\]
Quantize only when $a_i>0$. Use a fixed representative grid at each occupancy, independent of the alphabet
position, so an identity transition remains a valid state. All histories at the same cut and occupancy still
have the same legal future. The telescoping proof \eqref{eq:telescoping-quantization} is unchanged, but only
occupied-word transitions contribute error. This is a new schedule for B, not an assumption that zero counts
have small error.

Here is a complete sufficient cover, including its drift. Choose a common center $c$, replace $\phi_w$ by
$\psi_w=\phi_w-c$ and $y$ by $y-Kc$, and take an orthonormal basis $Q$ spanning all $\psi_w$. In its $d$
coordinates put
\[
 b_j=\max_w|(Q^T\psi_w)_j|,\quad
 D_j=\max_w(Q^T\psi_w)_j-\min_w(Q^T\psi_w)_j.
\]
Safe upper bounds can replace these maxima. A centered lattice with coordinate spacings $h_j$ rounds
each incoming coordinate by at most $h_j/2$. After any positive-count transition ending at occupancy $k$,
every reachable representative satisfies
\[
 |(Q^T\bar s)_j|\leq k(b_j+h_j/2).
\]
Indeed, its true increments have total count $k$ and there have been at most $k$ rounding operations.
The same inequality includes every incoming transition after its next rounding. Thus it suffices to store,
per occupancy, the lattice points in these boxes, with the conservative bound
\begin{equation}
 R\leq\prod_{j=1}^d\left(1+\frac{2Kb_j}{h_j}+K\right).
 \label{eq:occupied-cover}
\end{equation}
Actual reachable lattice points may be fewer. This construction does not assume a cover of exact signals
also covers approximate signals.

With $K-k$ users left, the future width of a rounding displacement is at most
$(K-k)\sum_jD_jh_j/2$. Let $\theta=\sum_jD_jh_j/(4\sigma^2)$. The positive occupancies along a path are
strictly increasing integers $k_1<\cdots<k_L\leq K$, so
\begin{equation}
 E(a)\leq\theta\sum_{\ell=1}^L(K-k_\ell)
 \leq\theta\,\frac{K(K-1)}2.
 \label{eq:occupied-error}
\end{equation}
The scalar correction uses a certified interval over the still-allowed future words; the global coordinate
box above supplies one without a search over completions. Distinct sampling has a smaller future set and
obeys the same conservative bound. At occupancy $K$ the future is zero and the final merge is exact after
scalar absorption. For $K=1$ every such merge is exact, although all candidate likelihoods must still be summed.

This replaces a possible $M$-fold accuracy penalty by an occupancy-controlled one. For $K>1$, one sufficient
choice is $\theta\leq2E_*/[K(K-1)]$, followed by \eqref{eq:quantization-error}. It does \emph{not} remove the
$O(MK^2R)$ transition cost or the $M$ candidate queries. The global grid can also be larger than a finely
adapted cut-specific grid; its dimension and inverse-noise dependence remain in \eqref{eq:occupied-cover}.
Directions, bounds, scalar factors and the reverse marked pass have the same accounting obligations as
before. The useful deduction is specific: traversing many absent words need not itself accumulate
approximation error.

\clearpage
\section{Entry C: retain exact regional counts and simplify only cross interactions}
\label{sec:entry-regions}
\subsection{Start with disjoint alphabets so the source remains exact}
Partition the complete message alphabet into $G$ disjoint regions $\mathcal C_r$ of size $L_r$. A region
could contain one pattern, a prefix or a chosen group of full messages. Its variable $a_r$ is an entire local
count configuration, with occupancy $k_r$ and source weight $\beta_r(a_r)$. Retain $\sum_r k_r=K$ exactly.
Distinct-source exclusions are internal to these disjoint alphabets; they need not be approximated by
independent labeled users. A region has
\[
 D_r=\sum_{k=0}^K H(L_r,k)
\]
possible local states, with infeasible distinct terms zero. Large $D_r$ is one of this construction's costs.

Choose one common center $c\in\R^n$, and set
$\psi_w=\phi_w-c$, $x_r=\sum_{w\in\mathcal C_r}a_w\psi_w$, $y'=y-Kc$. This is an exact change of reference:
$y-\Phi a=y'-\sum_r x_r$. After discarding an observation-only constant, the weight is
\begin{equation}
 \omega(a)\ \propto\ \ind\{\textstyle\sum_r k_r=K\}
 \prod_r \underbrace{\beta_r(a_r)
 e^{(y'^Tx_r-\|x_r\|^2/2)/\sigma^2}}_{f_r(a_r)}
 \prod_{r<q}e^{-x_r^Tx_q/\sigma^2}.
 \label{eq:regional-exact}
\end{equation}
All within-region interference, same-word multiplicity and pattern/payload association remain in $a_r$.
The local factors are positive on their legal domains. This is an exact count representation before any
edge is removed, not an assumption that regions are independent.

\subsection{Delete residual edges, and bound exactly what was deleted}
Retain a regional edge set $\mathcal E$ and replace every other cross factor by $1$. Denote the resulting
positive, cardinality-constrained law by $Q$. Define
\[
 \gamma_{rq}=\max_{u\in\mathcal C_r,\ v\in\mathcal C_q}
          \frac{|\psi_u^T\psi_v|}{\sigma^2},\qquad
 E(k)=\sum_{(r,q)\notin\mathcal E,\ r<q}\gamma_{rq}k_rk_q.
\]
By expanding the two sums in $x_r^Tx_q$, the log-weight error obeys
\begin{equation}
 \left|\log\frac{\omega(a)}{\widetilde\omega(a)}\right|\leq E(k),\qquad
 E_*:=\max_{\substack{k\text{ feasible}\\\sum k_r=K}}E(k)
 \leq\binom K2\gamma_{\max}.
 \label{eq:edge-error}
\end{equation}
Here the same common constant is used for both unnormalized weights. The last inequality counts only
cross-region pairs; the maximum over an empty omitted edge set is zero. Consequently the bounds in
\eqref{eq:quantization-error} hold with this $E_*$. For a target small bound at growing $K$, a worst-case
argument can require residual interaction of order $K^{-2}$, not merely a small fixed correlation.

The four-point invariant explains why the reference should be chosen before judging these couplings.
Nevertheless, a large $\gamma_{\max}$ is only a failed sufficient bound, not proof that deletion is poor.
Tighter signed sums can have smaller range, and occupancy-weighted envelopes can replace the maximum.
An individual edge's error also cannot be judged solely by its covariance with a decision under the original
reference: multiple reweightings change that reference, as Report 7's counterexample shows.

\subsection{A forest gives an explicit tractable contraction conditional on local sizes}
Suppose $\mathcal E$ is a tree (a forest is combined by an additional count convolution). For a directed
edge $r\to q$, keep the parent's full regional state $a_q$ and a polynomial whose exponent records the
number of users on the $r$ side. The exact message for this surrogate is
\begin{equation}
 m_{r\to q}(a_q,z)=\sum_{a_r} f_r(a_r)e^{-x_r^Tx_q/\sigma^2}z^{k_r}
       \prod_{v\in\partial r\setminus q}m_{v\to r}(a_r,z),
 \label{eq:regional-message}
\end{equation}
truncated above degree $K$. Extracting degree $K$ at the root enforces the known total. A reverse pass gives
the local state marginal $Q(a_r)$ and hence every $Q(a_w\geq j)$ inside that region. These are exact
probabilities for the specified approximate law, not locally normalized independent answers.

With direct polynomial products, prefix/suffix reuse around each node and stored local states, both passes
have the conservative arithmetic cost
\begin{equation}
 O\!\left(K\sum_{(r,q)\in\mathcal E}D_rD_q
       +K^2\sum_r(1+\deg r)D_r\right).
 \label{eq:region-cost}
\end{equation}
For each $a_r$, multiply the incoming degree-$K$ polynomials; then sum over $a_r$ for every $a_q$ and
coefficient. This proves the two terms. Readout adds $O(K\sum_r D_r)$ sparse count-tail accumulation and
score selection. Message storage is $O(K\sum_r(1+\deg r)D_r)$; storing all edge tables adds
$O(\sum_{(r,q)\in\mathcal E}D_rD_q)$. Direct construction of those tables adds
$O(n\sum_{(r,q)\in\mathcal E}D_rD_q)$ work, besides local-state construction and omitted-edge certification.
Direct local signal construction costs $O(nK\sum_rD_r)$; directly certifying every $\gamma_{rq}$ costs
$O(n\sum_{(r,q)\notin\mathcal E,\ r<q}L_rL_q)$, potentially $O(nM^2)$ even for a sparse retained graph.
Bounds from the encoding geometry could reduce that cost, with their own accuracy and construction account.
Structured polar/parity contractions may replace a dense edge multiplication, but their required joint
boundary must still be checked. A general retained graph can instead have exponential separator-table cost.

Thus sparse \emph{retained} interactions and small local state spaces are separate assumptions. The discarded
interactions must also have a useful aggregate error bound. A small region in signal dimension need not have
few count configurations. Capping occupancy at $b<K$ reduces its local space to
$\sum_{k=0}^bH(L_r,k)$, but deletes all configurations exceeding that cap; Entry A must account for their
mass. Unknown patterns are included in the alphabet, not supplied by an oracle. If the entire alphabet is
listed to build the regions, this construction still visits $M=2^B$ messages.

\subsection{Why independently centering each region can reintroduce a hard sum}
One might obtain much smaller local variation using a different center $c_r$ in each region. Put
$x_r=s_r-k_rc_r$. Now the residual is
\[
 y-\sum_r s_r=\left(y-\sum_r k_rc_r\right)-\sum_r x_r.
\]
For \emph{fixed} occupancy vector $k$, this gives a valid centered regional calculation. But its effective
observation depends on every $k_r$. Across unknown occupancies, those supposedly unary factors need not
remain local. Even if every within-region deviation $x_r$ vanished, the term
$\|y-\sum_r k_rc_r\|^2$ would still couple the occupancies.

An exact remedy is to sum the unnormalized sector answers over all feasible $k$, retaining their source
weights. With unrestricted capacities there are $\binom{G+K-1}{K}$ sectors; fixing occupancies removes the
count-polynomial work inside each sector but can add this large outer sum. A different remedy would retain
the center/occupancy interactions as explicit factors and analyze their separators. Neither remedy is free.
The common-center construction avoids this issue because $\sum_r k_rc=Kc$ is known. This is a concrete
case where a promising local approximation needs additional global algebra before it defines a decoder.

\clearpage
\section{Entry D: positive features for cross interactions}
\label{sec:positive-features}
\subsection{First separate small geometric dimension from small interaction rank}
Entry B replaces several rows by a rescaled representative. Entry C can replace an edge by one. Neither
exhausts what a positive boundary can do. Consider a retained regional factor
$b(x_r,x_q)=\exp(-x_r^Tx_q/\sigma^2)$. Suppose $x_r=U_r z_r$ and $x_q=U_q z_q$ in orthonormal local bases.
The exact interaction is $z_r^TC_{rq}z_q$, where
\[
 C_{rq}=U_r^TU_q/\sigma^2.
\]
Take a singular-value decomposition $C_{rq}=\sum_j\lambda_j u_jv_j^T$. Keeping $d$ singular directions gives
\[
 \alpha_j=\sqrt{\lambda_j}\,u_j^Tz_r,\qquad
 \beta_j=\sqrt{\lambda_j}\,v_j^Tz_q,\qquad
 z_r^TC_{rq}z_q=\sum_{j=1}^d\alpha_j\beta_j+r_{rq}.
\]
If the state domains satisfy $\|z_r\|\leq A_r$, $\|z_q\|\leq A_q$, then
\begin{equation}
 |r_{rq}|\leq\epsilon_{rq}:=\lambda_{d+1}A_rA_q.
 \label{eq:geometric-rank}
\end{equation}
Taking all nonzero singular directions gives $\epsilon_{rq}=0$. Bounds conditional on local occupancies
can be sharper. Bases, singular directions and safe domain bounds have construction costs; small $d$ is a
condition to investigate, not a consequence of calling the code polar.

The exponential of a rank-$d$ bilinear form need not have rank $d$ as a table. The next derivation supplies
the missing relation between interaction strength, accuracy and a separable feature count.

\subsection{Shift the coordinates so the expansion is nonnegative}
Choose coordinate bounds
$\alpha_j\in[l_j,l_j+A_j]$ and $\beta_j\in[u_j-D_j,u_j]$, with $A_j,D_j\geq0$.
Define $X_j=\alpha_j-l_j\geq0$ and $Y_j=u_j-\beta_j\geq0$. Then
\[
 -\alpha^T\beta=-l^T\beta-X^Tu+X^TY.
\]
The first two terms are unary. Thus the retained-direction factor has the exact form
\begin{equation}
 b_d(x_r,x_q)=a(x_r)c(x_q)e^\xi,\qquad
 a=e^{-X^Tu},\quad c=e^{-l^T\beta},\quad
 \xi=X^TY\in[0,\Lambda],\quad \Lambda=\sum_jA_jD_j.
 \label{eq:positive-feature-origin}
\end{equation}
For an integer $L\geq0$, replace $e^\xi$ by its degree-$L$ Taylor polynomial. The multinomial identity gives
\begin{equation}
 \widetilde b_{d,L}
 =a(x_r)c(x_q)\sum_{|\gamma|\leq L}\frac{X^\gamma Y^\gamma}{\gamma!},
 \qquad R_{d,L}=\binom{d+L}{L}.
 \label{eq:positive-features}
\end{equation}
Here $\gamma$ is a vector of nonnegative integers. Each summand is a product of one nonnegative feature
of $x_r$ and one of $x_q$. The constant feature makes the whole factor strictly positive even when some
coordinates vanish. No likelihood state or legal source assignment has been removed.
For one retained direction and $L=2$, the three features are simply
$1$, $XY$ and $X^2Y^2/2$, multiplied by the unary factors. The endpoint states remain different; their
interaction is communicated through these three numbers rather than through a table entry for every pair.

Define
\[
 p_L(\Lambda)=e^{-\Lambda}\sum_{m=0}^L\frac{\Lambda^m}{m!}>0,\qquad p_L(0)=1.
\]
Since $e^{-\xi}\sum_{m=0}^L\xi^m/m!$ decreases with $\xi\geq0$,
\begin{equation}
 p_L(\Lambda)\leq\frac{\widetilde b_{d,L}}{b_d}\leq1.
 \label{eq:feature-ratio}
\end{equation}
This is also the probability that a Poisson variable of mean $\Lambda$ is at most $L$; that interpretation
does not introduce a random approximation. Combining the geometric remainder and the polynomial error,
\[
 \log\frac{b}{\widetilde b_{d,L}}
 \in[-\epsilon_{rq},\,\epsilon_{rq}-\log p_L(\Lambda)].
\]
For several replaced factors on the same legal configurations, their log-ratio ranges add. Consequently
\begin{equation}
 \TV(P_y,Q)\leq
 \tanh\!\left(\frac14\sum_{\text{replaced }(r,q)}
       [2\epsilon_{rq}-\log p_{L_{rq}}(\Lambda_{rq})]\right).
 \label{eq:feature-tv}
\end{equation}
Other retained-support modifications add their own ranges; actual support restriction still costs Entry A's
omitted mass. A small local singular-value tail and a short positive expansion are two separate conditions.

\subsection{What controls the required number of features?}
With exact geometric directions, increasing $L$ gives nested lower approximations to every weight and
evidence. A uniform upper weight is the lower weight divided by the product of the $p_L$ factors. Applying
the positive event/complement envelopes below can give more useful tail intervals than one global TV bound.
The four-point obstruction to \emph{one} feature is respected: at $L=0$, the interaction error can be as
large as $\Lambda$.

For an explicit sufficient degree, let $m=L+1>\Lambda>0$. A Poisson variable $N$ has
$\E e^{tN}=\exp[\Lambda(e^t-1)]$. Markov's inequality with $e^t=m/\Lambda$ gives
\[
 1-p_L(\Lambda)=\Pr(N\geq m)
 \leq e^{-\Lambda}\left(\frac{e\Lambda}{m}\right)^m.
\]
Thus a desired relative error specifies a sufficient $L$ without evaluating the original table. But even
at $d=1$, increasing interaction strength can require increasing $L$; at larger $d$ the count
$\binom{d+L}{L}$ can be prohibitive. A small geometric interface does not imply a cheap exponential kernel.
The coordinate-box value $\Lambda$ is a sufficient bound and can be loose when coordinate extrema are
incompatible. Establishing a smaller attainable maximum is additional work.

\subsection{Use the features in a complete cardinality-constrained contraction}
On Entry C's retained forest, write a replaced edge as
$\sum_{\ell=1}^{R_{rq}}u_{rq,\ell}(a_r)v_{rq,\ell}(a_q)$.
For each polynomial coefficient in \eqref{eq:regional-message}, first sum the weighted $r$-states against
each $u_{rq,\ell}$, then expand in the parent's $v_{rq,\ell}$. The edge multiplication now costs
$O(KR_{rq}(D_r+D_q))$, instead of $O(KD_rD_q)$. Both passes and all count readout therefore have the bound
\begin{equation}
 O\!\left(K\sum_{(r,q)\in\mathcal E}R_{rq}(D_r+D_q)
       +K^2\sum_r(1+\deg r)D_r+K\sum_rD_r\right).
 \label{eq:feature-cost}
\end{equation}
Local-state enumeration and the exact factorial/distinct priors remain. Feature construction costs, for
example, $O(\sum_{(r,q)}d_{rq}R_{rq}(D_r+D_q))$ by direct monomial evaluation, in addition to computing the
coordinates and bounds. Storing the features costs $O(\sum_{(r,q)}R_{rq}(D_r+D_q))$; node polynomial
storage is as in C. Edge tables of size $D_rD_q$ need not be materialized.

For a graph with cycles, introducing one feature index per edge does not make the entire contraction a
forest. Joint indices can multiply in an elimination separator. Likewise, a cheap isolated boundary
replacement is not yet a complete compressed polar recursion. These are genuine cost boundaries, not
source-independence assumptions. On an already chosen forest, however, \eqref{eq:feature-cost} is a complete
positive refinement of C with retained cardinality and full count tails.

Large unary rescalings and monomials may need scaled arithmetic. The bounds above are in exact real
arithmetic; rounding enclosures and the precision needed for the final margin would remain implementation
costs. The algebra does not establish a useful $d,L$ for the fixed bank.

\clearpage
\section{Entry E: make connected-component truncation a positive restriction}
\label{sec:positive-clusters}
\subsection{An exact rescaling creates a different expansion}
Report 7 expands pair factors as $b=1+f$, where $f=b-1$ may be signed. Keeping a few connected orders
in that expansion need not define probabilities. There is another exact reference for the IID source.
Write its labeled-slot posterior as
\[
 P_y(W)\propto\prod_{i=1}^K\nu(w_i)\prod_{i<j}b(w_i,w_j),\qquad
 \nu(w)=\frac{e^{y^T\phi_w/\sigma^2}}{\sum_u e^{y^T\phi_u/\sigma^2}},\qquad
 b(u,v)=e^{-\phi_u^T\phi_v/\sigma^2}.
\]
Choose a certified scalar $0<b_0\leq\min_{u,v}b(u,v)$, and put
\[
 f_+(u,v)=b(u,v)/b_0-1\geq0.
\]
The factor $b_0^{\binom K2}$ is constant, so replacing every $b$ by $1+f_+$ changes no posterior.
Expand the product using a graph $G$ on the $K$ labeled slots:
\begin{equation}
 \Omega(W,G)=\prod_i\nu(w_i)\prod_{(i,j)\in E(G)}f_+(w_i,w_j).
 \label{eq:positive-graph}
\end{equation}
Summing $G$ restores the exact rescaled weight of $W$. All joint weights are nonnegative; the empty graph
has strictly positive weight for every $W$. The graphs are algebraic auxiliary variables. They are not a
claim that physical resource overlaps are independently random.

Let $E_q$ be the event that every connected component of $G$ has at most $q$ vertices, with $1\leq q\leq K$.
Restrict this \emph{joint} law to $E_q$, then marginalize $G$, obtaining $Q_q(W)$. Every legal IID tuple
still has positive approximate weight. Nevertheless, the restriction theorem applies on the larger space:
\begin{equation}
 \TV(P_y,Q_q)\leq\eta_q:=\Pr(E_q^c\mid y),\qquad
 R(d_q)-R(d^*)\leq\min(1,2\eta_q).
 \label{eq:cluster-restriction}
\end{equation}
Marginalization cannot increase TV, since any event in $W$ is also a joint event. This is different from
restricting the transmitted message list: it discards auxiliary explanations of each tuple. At $q=K$ it
is exact; at $q=1$ it is the product reference $\nu^{\otimes K}$.

\subsection{Assemble all small components, including their multiplicities}
For a queried word $w$, let $J_w(W)=\sum_i\ind\{w_i=w\}$ and define the nonnegative polynomial
\[
 C_{m,w}(z)=\E_{\nu^{\otimes m}}\!\left[
 z^{J_w(W)}\sum_{\substack{G\text{ connected}\\\text{on }m\text{ labeled slots}}}
       \prod_{(i,j)\in E(G)}f_+(w_i,w_j)\right].
\]
For $m=1$ the edge product is one. Grouping graphs by the component containing slot 1 gives the finite
recurrence
\begin{equation}
 B^{(q)}_{0,w}(z)=1,\qquad
 B^{(q)}_{k,w}(z)=\sum_{m=1}^{\min(k,q)}\binom{k-1}{m-1}
                 C_{m,w}(z)B^{(q)}_{k-m,w}(z).
 \label{eq:positive-cluster-assembly}
\end{equation}
It includes arbitrarily many small components, not just one group of $q$ users. In particular, $q=2$
keeps disjoint interacting pairs and singletons; $q=3$ additionally keeps the three chains and the triangle
within each three-slot component. It does not discard all terms above degree $q$ in an interaction strength.
The full count law is
\[
 Q_q(a_w=j)=\frac{[z^j]B^{(q)}_{K,w}(z)}{B^{(q)}_{K,w}(1)}.
\]
Every query uses the same denominator. Exchangeability retains the exact count multiplicity
$K!/\prod_w a_w!$. These are count tails of one common positive law, with repetitions retained.
For example, with $K=4$ and unmarked $C_1=1$,
\[
 B_4^{(2)}(1)=1+6C_2+3C_2^2,\qquad
 B_4^{(3)}(1)=1+6C_2+3C_2^2+4C_3.
\]
The first formula already includes two simultaneous interacting pairs. The missing term at $q=3$ is
the connected four-slot coefficient $C_4$, not every possible effect involving four slots.

\subsection{One explicit error bound uses only the chance of an auxiliary edge}
Conditional on $W$, the joint graph has independent edges with probabilities
\[
 p(u,v)=\frac{f_+(u,v)}{1+f_+(u,v)}=1-\frac{b_0}{b(u,v)}.
\]
Suppose $p(u,v)\leq p_*$. A component of size at least $q+1$ contains a tree on some $q+1$ of its vertices.
For any specified tree its $q$ edges are all present with conditional probability at most $p_*^q$.
There are $\binom K{q+1}$ vertex subsets and $(q+1)^{q-1}$ labeled trees on each. A union bound,
valid for every $W$ and therefore also after its posterior average, proves
\begin{equation}
 \eta_q\leq
 \min\!\left\{1,\binom K{q+1}(q+1)^{q-1}p_*^q\right\},
 \quad 1\leq q<K;\qquad \eta_K=0.
 \label{eq:cluster-tree-probability}
\end{equation}
This answers a limited version of the question about the first few connected orders: higher components
do have a controlled tail when the \emph{rescaled} edge probabilities meet this bound. It does not follow
from the original log-likelihood being pairwise.

For perspective, $r!\geq(r/e)^r$ gives, with $r=q+1$,
\[
 \binom Kr r^{r-2}p_*^{r-1}
 \leq\frac{Ke}{r^2}(eKp_*)^{r-1}.
\]
Thus if $eKp_*<1$ with a fixed gap, a component cutoff growing logarithmically with $K$ suffices for a
small stated error. This is a sufficient, conservative regime, not a sharp graph threshold and not a
tractability result: obtaining the small-component coefficients still costs work exponential in their size.

\subsection{A weighted certificate reuses the exact polar pair operator}
A uniform $p_*$ can be dominated by irrelevant words. A second bound carries the observation-dependent
reference $\nu$. Put $F=f_+$, $\epsilon=\max F$, and let $Q_m$ be the sum of all weighted labeled trees
averaged under $\nu^{\otimes m}$. Report 7's finite tree argument now bounds the \emph{positive} connected
coefficient $C_m=C_{m,w}(1)$:
\[
 C_m\leq Q_m(1+\epsilon)^{(m-1)(m-2)/2}.
\]
Every connected graph has a spanning tree; summing trees and optional extra edges overcounts it safely.
Use exact $C_m$ through $m=q$ and these upper bounds above $q$ as coefficients $\overline C_m$.
Assemble an upper evidence by
\[
 U_0=1,\qquad U_k=\sum_{m=1}^k\binom{k-1}{m-1}\overline C_m U_{k-m}.
\]
All coefficients are nonnegative, so with $Z_q=B^{(q)}_{K,w}(1)$ and exact rescaled evidence $Z_+$,
\begin{equation}
 Z_q\leq Z_+\leq U_K,\qquad
 \eta_q=1-\frac{Z_q}{Z_+}\leq1-\frac{Z_q}{U_K}.
 \label{eq:positive-cluster-majorant}
\end{equation}
Take the smaller valid bound from \eqref{eq:cluster-tree-probability} and
\eqref{eq:positive-cluster-majorant}. The latter can still be loose: optional-edge overcounting remains.

The tree coefficients obey
\[
 T_u(t)=t\exp\!\left(\sum_v F(u,v)\nu(v)T_v(t)\right),\qquad
 Q_m=(m-1)!\sum_u\nu(u)[t^m]T_u(t).
\]
Through degree $K$, the coefficient recurrence needs $K-1$ applications of $F$, $O(K^2M)$ additional
arithmetic and $O(KM)$ storage. Here there is a useful structural difference from a general entrywise
absolute remainder kernel:
\begin{equation}
 F x=b_0^{-1}B x-\one(\one^Tx),\qquad B_{uv}=b(u,v).
 \label{eq:positive-cluster-operator}
\end{equation}
Thus the polar pair operator applies directly, with one scalar subtraction per output. The positive
rescaling retains its XOR/parity structure. A small computed difference can require controlled precision;
nonnegativity of the exact operator is not a guarantee about floating-point subtraction.

\subsection{Full query cost and the price of positivity}
Let $C_{\rm blocks}(q)$ include construction of all connected coefficients and their count marks.
One direct bound is
$O(\sum_{m=1}^q m^2 M^m2^{\binom m2})$ arithmetic once the pair entries are available.
For each tuple, its connected weight can update the coefficient for each distinct occupied word and its
count, at most $m$ updates. Unmarked coefficients are shared; zero-occurrence coefficients follow by
subtracting the positive-count totals. This avoids rerunning tuple enumeration $M$ times. Exact polar
pair and three-user contractions can reduce particular block computations, but not all higher orders
automatically.

For each word, assemble its degree-$m$ connected polynomials through $m=q$ with
\eqref{eq:positive-cluster-assembly}. Direct polynomial products cost at most $O(K^2q^2)$ per word.
Thus an explicit full-count construction costs
\[
 C_{\rm blocks}(q)+O(MK^2q^2+K C_{\rm pair}+K^2M),
\]
including the weighted tree certificate when used. Storing the marked connected coefficients uses
$O(Mq^2)$ numbers; assembly can stream words using $O(K^2)$ scratch, with $O(MK)$ output tails.
Pair storage, waveform access, reference construction and block workspaces are additional. Even a small
$q$ is still explicit in $M=2^B$; a logarithmic $q$ can leave $M^q$ unaffordable.

The main limitation is visible before any experiment. For unit-energy IID messages, the universal choice
$b_0=e^{-1/\sigma^2}$ is attained on repeated words. Since $b(u,v)\leq e^{1/\sigma^2}$, the resulting
uniform bound is $p_*\leq1-e^{-2/\sigma^2}$. At low physical noise this can be almost one. Even when most
distinct-word correlations are small, the diagonal can force a floor far below their factors, making many
auxiliary edges likely. An exact reference change may help but has no guaranteed favourable range.
The benefit of positivity must be weighed against this inflation.

The same component assembly does \emph{not} extend unchanged to the distinct-message control: its hard
pair zeros admit no positive floor. Keeping those exclusions outside the graph couples different components,
destroying the product-reference factorization used above. Report 7's signed construction still handles
that control. This entry is an IID construction with an explicit error certificate, not a relaxation of
distinctness. Nor does it supersede the signed expansion: a better signed remainder can be preferable
when forcing positivity is too expensive.

\clearpage
\section{Entry F: close the residual distribution through its smoothed likelihood}
\label{sec:residual-closure}
\subsection{Derive the excluded-word queries before choosing moments}
For the IID source, the exact residual evidence is an expectation over a sum of independent prior
waveforms. A moment approximation must apply to that expectation, not silently replace interference by
an extra physical noise term. Fix a queried word $w$, assume $M\geq2$, and put $p=1/M$.
Let $X^{(-w)}$ be a uniform waveform on the other $M-1$ words. Define
\[
 L_h^{(-w)}(r)=
 \E\exp\!\left[-\frac{\|r-\sum_{i=1}^hX_i^{(-w)}\|^2}{2\sigma^2}\right],
 \qquad L_0^{(-w)}(r)=e^{-\|r\|^2/(2\sigma^2)}.
\]
Conditioning the prior on the exact count of $w$ gives its unnormalized count masses:
\begin{equation}
 N_{w,j}=\binom Kj p^j(1-p)^{K-j}
       L_{K-j}^{(-w)}(y-j\phi_w),\qquad j=0,\ldots,K.
 \label{eq:residual-count-mass}
\end{equation}
For every $w$, $\sum_jN_{w,j}=Z$ from \eqref{eq:target}. Excluding $w$ is essential: without the exclusion,
the remaining draws can add further copies and the expression is not an exact-count probability.
The case $M=1$ has deterministic counts and needs no closure.

\subsection{A positive Gaussian closure, and what it does not guarantee}
Let $\mu_w=\E X^{(-w)}$ and $\Sigma_w=\operatorname{Cov}(X^{(-w)})$. Replace the $h$-word sum by a
Gaussian with mean $h\mu_w$ and covariance $h\Sigma_w$. Completing the square in the Gaussian integral gives
\begin{equation}
 \widetilde L_h^{(-w)}(r)
 =\det(I+h\Sigma_w/\sigma^2)^{-1/2}
   \exp\!\left[-\frac12(r-h\mu_w)^T
     (\sigma^2I+h\Sigma_w)^{-1}(r-h\mu_w)\right].
 \label{eq:gaussian-residual}
\end{equation}
It is well defined even for singular $\Sigma_w$ because $\sigma^2>0$. The determinant is part of the
likelihood and cannot be dropped when comparing different $w,j$. At $h=0$ the formula is exact.
All moments come from the fixed prior and encoder; none is fitted from the transmitted messages.

Substitute \eqref{eq:gaussian-residual} in \eqref{eq:residual-count-mass}, then normalize the $K+1$ positive
masses separately for each $w$. This gives genuine per-word count distributions and monotone tails.
It need not give the marginals of a single joint law on $\sum_w a_w=K$: the approximate evidence
$\sum_j\widetilde N_{w,j}$ can differ between words. We therefore use direct tail-error bounds and
Section~\ref{sec:decision-foundation}, not a claimed full-posterior TV bound.
The greedy allocation remains feasible and maximizes these monotone approximate gains.

\subsection{A finite smoothing bound instead of an unquantified Gaussian claim}
Here is a sufficient error bound for any independent waveform law $X$ with mean $\mu$, covariance $\Sigma$
and finite third absolute centered moment. Let $\xi=X-\mu$ and $G\sim\mathcal N(0,\Sigma)$.
For the kernel $f_r(s)=e^{-\|r-s\|^2/(2\sigma^2)}$, its third directional derivative obeys
\[
 |D^3f_r(s)[v,v,v]|\leq \frac{3}{\sigma^3}\|v\|^3.
\]
To see a safe constant, differentiation produces a cubic term and three linear terms in
$(r-s)/\sigma$; their norm is bounded by
$(t^3+3t)e^{-t^2/2}<3$ for $t\geq0$.
Taylor expansion through order two, conditional on all other variables, therefore has remainder at most
$\|v\|^3/(2\sigma^3)$. Replacing one $\xi$ by one $G$ cancels the constant, first and second terms because
their first two moments agree. Replacing all $h$ increments and summing remainders proves
\begin{equation}
 |L_h(r)-\widetilde L_h(r)|
 \leq\delta_h:=
 \min\!\left\{1,\frac{h}{2\sigma^3}
       [\E\|\xi\|^3+\E\|G\|^3]\right\}.
 \label{eq:gaussian-smoothing-coarse}
\end{equation}
The bound is uniform in the queried residual $r$ and uses the actual physical smoothing scale $\sigma$.
The cap one follows because both expectations lie in $[0,1]$. This is a finite replacement proof, not an
appeal to asymptotic normality.

There is a sharper version using smoothing already supplied by the Gaussian increments in the telescoping
sequence. In the step with $i-1$ original increments and $h-i$ Gaussian increments other than the replaced one,
put
\[
 A_i=\sigma^2I+(h-i)\Sigma,\qquad
 c_i=\det(I+(h-i)\Sigma/\sigma^2)^{-1/2}.
\]
Integrating the latter increments first gives a kernel of height $c_i$ and covariance $A_i$. The same
Taylor proof, now in the coordinates $A_i^{-1/2}s$, gives
\begin{equation}
 |L_h(r)-\widetilde L_h(r)|
 \leq\min\!\left\{1,\frac12\sum_{i=1}^h c_i
 \left[\E\|A_i^{-1/2}\xi\|^3+\E\|A_i^{-1/2}G\|^3\right]\right\}.
 \label{eq:gaussian-smoothing-shaped}
\end{equation}
Each term is no larger than the corresponding coarse term. It can expose useful smoothing in directions
with residual variance, but does not assert a useful bound at the intended noise or load.

Simple certificates need not enumerate a third moment for every excluded alphabet. Since the codewords
have norm one, $\|\xi\|\leq2$ and
\[
 \E\|\xi\|^3\leq2\operatorname{tr}\Sigma,\qquad
 \E\|G\|^3\leq3^{3/4}(\operatorname{tr}\Sigma)^{3/2}.
\]
The Gaussian inequality follows from
$\E\|G\|^4=(\operatorname{tr}\Sigma)^2+2\operatorname{tr}(\Sigma^2)$ and the power-mean inequality.
For the shaped bound one can use
\[
 \E\|A_i^{-1/2}\xi\|^3
 \leq 2\|A_i^{-1/2}\|\,\operatorname{tr}(A_i^{-1}\Sigma),\qquad
 \E\|A_i^{-1/2}G\|^3
 \leq3^{3/4}\operatorname{tr}(A_i^{-1}\Sigma)^{3/2}.
\]
Sharper directional moments might help, but constructing them is additional work.

\subsection{Carry the remainder through the count normalization}
Use either bound with $(\mu,\Sigma)=(\mu_w,\Sigma_w)$ and call it $\delta_{w,h}$.
Equation \eqref{eq:residual-count-mass} gives
\[
 |N_{w,j}-\widetilde N_{w,j}|
 \leq d_{w,j}:=\binom Kj p^j(1-p)^{K-j}\delta_{w,K-j}.
\]
Put $L_{w,j}=\max(0,\widetilde N_{w,j}-d_{w,j})$ and
$U_{w,j}=\widetilde N_{w,j}+d_{w,j}$. For a tail at $l$, define $L_A,U_A$ by summing these bounds over
$j\geq l$, and $L_{\bar A},U_{\bar A}$ over $j<l$. Then
\begin{equation}
 \frac{L_A}{L_A+U_{\bar A}}\leq T_{w,l}
 \leq\frac{U_A}{U_A+L_{\bar A}}.
 \label{eq:gaussian-count-envelopes}
\end{equation}
Use $[0,1]$ when necessary. Subtract the normalized approximate tail from the two endpoints to obtain its
maximum absolute error; the largest such error gives the Bayes excess-loss bound. Since each word yields
an interval for the same exact $Z$, their intersections can also tighten denominator bounds.
All statements concern the original IID count law even though the approximate per-word laws may be mutually
inconsistent.

The important assumption is now precise: these \emph{smoothed}, marked residual errors must be small
relative to the evidence relevant to the tail or decision. Merely matching a covariance is insufficient.
For example, let a scalar $X$ be $\pm1$ with equal probability and take $r=1$, $h=1$.
As $\sigma\downarrow0$, $L_1(1)\to1/2$, whereas its moment-matched Gaussian closure gives
\[
 \widetilde L_1(1)=\frac{\sigma}{\sqrt{\sigma^2+1}}
             e^{-1/[2(\sigma^2+1)]}\longrightarrow0.
\]
The first two moments agree exactly but the smoothed likelihood does not. A small absolute error is also
insufficient when the normalizer is still smaller. Neither many users nor a large ambient dimension alone
settles these conditions.

\subsection{Use polar characters for the moments, and charge for all count queries}
Let $\mu,\Sigma$ be the unconditional one-word mean and covariance and $\delta_w=\phi_w-\mu$.
The excluded moments follow exactly from one rank-one correction:
\begin{equation}
 \mu_w=\frac{\mu-p\phi_w}{1-p},\qquad
 \Sigma_w=\frac{\Sigma}{1-p}
          -\frac{p}{(1-p)^2}\delta_w\delta_w^T.
 \label{eq:excluded-moments}
\end{equation}
A direct pass over the alphabet constructs the global moments in $O(Mn^2)$ work. For this encoder, the
character identities can remove that payload enumeration. If $\ell_t$ is row $t$ of the binary generator $L$,
uniform payload bits satisfy
\[
 \E_v(-1)^{\ell_t v}=\ind\{\ell_t=0\},\qquad
 \E_v(-1)^{\ell_t v+\ell_u v}=\ind\{\ell_t=\ell_u\}.
\]
With $D_p=\diag(\sqrt{e_p})$, $z_t=\ind\{\ell_t=0\}$ and
$J_{tu}=\ind\{\ell_t=\ell_u\}$, the exact global moments are
\[
 \mu=\frac1P\sum_pS_pD_pz,\qquad
 \Sigma=\frac1P\sum_pS_pD_pJD_pS_p^T-\mu\mu^T.
\]
Zero or repeated generator rows must be retained; silently assuming all symbol signs independent would
give the wrong covariance. Direct sparse insertion costs $O(PN_c^2+n^2)$ after constructing the row classes
and patterns, with $O(n^2)$ covariance storage. Pattern enumeration remains.

For each $h$, pre-factorize $A_h^0=\sigma^2I+h\Sigma/(1-p)$. Equation \eqref{eq:excluded-moments} shows that
$\sigma^2I+h\Sigma_w$ is its rank-one downdate. Matrix determinant and inverse rank-one formulas therefore
give all Gaussian queries with the conservative dense arithmetic cost
\[
 O(Kn^3+MKn^2)
\]
after the moments, instead of factoring a new $n$-by-$n$ matrix for every $w,j$. The covariance plus
$\sigma^2I$ remains positive definite; a nearly singular downdate formula requires appropriate numerical
precision, not an ad hoc change of variance. Processing one $h$ at a time uses $O(n^2+MK)$ storage for
the matrix work and all count masses. Waveform generation and constructing the reference moments remain
charged separately.

The coarse certificates use traces from \eqref{eq:excluded-moments}, hence only $O(MK)$ further scalar work
once the column data are available. Fully evaluating the shaped certificate for every $w,h$ has extra
matrix/trace costs; it is not included for free in the displayed Gaussian-query bound. These formulas replace
the combinatorial residual sums by moment/matrix calculations, but the complete decoder still visits
$M$ words and can have vacuous error intervals. Distinct-message residuals are sampled without replacement;
the binomial decomposition and independent-increment proof above do not apply to that control unchanged.

\subsection{Retain separated modes and close only the variation inside them}
The scalar counterexample suggests a more selective assumption than making the entire residual sum Gaussian.
Partition the excluded-word alphabet into $G$ known groups with prior masses $\rho_g$, conditional means
$\mu_g$ and covariances $\Sigma_g$. Remove empty groups. For $h$ remaining IID words, their occupancy
vector $k$ has exact prior mass
\[
 \pi_h(k)=\frac{h!}{\prod_g k_g!}\prod_g\rho_g^{k_g},\qquad \sum_gk_g=h.
\]
Conditional on $k$, the sums inside the groups are independent. Gaussian replacement \emph{within} each
group therefore gives the positive mixture
\begin{equation}
 \widetilde L_h(r)=\sum_{\sum k_g=h}\pi_h(k)\,
 \det(I+V_k/\sigma^2)^{-1/2}
 e^{-\frac12(r-m_k)^T(\sigma^2I+V_k)^{-1}(r-m_k)},
 \quad m_k=\sum_gk_g\mu_g,\quad V_k=\sum_gk_g\Sigma_g.
 \label{eq:mixture-residual}
\end{equation}
Different group centers remain separated. Their unknown occupancy is summed with its correct prior,
not replaced by a mean occupancy or supplied by an oracle. For each sector the same replacement proof gives
\[
 |L_{h,k}(r)-\widetilde L_{h,k}(r)|
 \leq\delta(k):=\min\!\left\{1,\frac1{2\sigma^3}
       \sum_g k_g[\E\|X_g-\mu_g\|^3+\E\|G_g\|^3]\right\},
 \quad G_g\sim\mathcal N(0,\Sigma_g).
\]
Hence $|L_h-\widetilde L_h|\leq\sum_k\pi_h(k)\delta(k)$. A cheap further upper bound moves the minimum
outside the average and uses $\E k_g=h\rho_g$. Thus a certificate can be constructed from weighted
within-group third moments without enumerating sectors, although the mixture likelihood itself still
needs their contraction. Count queries continue to use the excluded-word groups and
\eqref{eq:gaussian-count-envelopes}.

There are $\binom{G+h-1}{h}$ sectors, each needing up to $O(n^3)$ matrix work per query before any
demonstrated reuse. A weaker within-group assumption therefore costs an extra occupancy sum.
This is C's regional-centering obstacle in another exact representation.

Point-mass groups give exact closures and preserve the discrete counterexample's modes. But if every
message needs its own group, the mixture restores count enumeration. Useful groupings must control both
within-group variation and occupancy-sector work. Dropping sectors needs an evidence bound:
prior-weighted \emph{likelihood error} and posterior omitted mass are different quantities.

\section{Entry G: truncate Fourier work with marked remainder bounds}
\label{sec:fourier-approximation}
\subsection{The evidence and count queries damp different powers}
Use Report 7's exact integer signatures $\phi_w=Ct_w$, with $t_w\in\{-1,0,1\}^d$.
On the centered grid $\mathcal G=\{-K,\ldots,K\}^d$, let $Q=2K+1$, $N_g=Q^d$ and
$\theta_\ell=2\pi\ell/Q$ for $\ell\in\{0,\ldots,Q-1\}^d$. Define
\[
 g_y(t)=e^{-\|y-Ct\|^2/(2\sigma^2)},\quad
 c_\ell=\frac1{N_g}\sum_{t\in\mathcal G}g_y(t)e^{i\theta_\ell^Tt},\quad
 \zeta_{\ell,w}=e^{-i\theta_\ell^Tt_w},\quad
 \chi_\ell=\frac1M\sum_w\zeta_{\ell,w}.
\]
These conventions give $g_y(t)=\sum_\ell c_\ell e^{-i\theta_\ell^Tt}$. Under the IID prior,
the normalized evidence and exact-count masses are therefore
\begin{align}
 Z&=\sum_\ell c_\ell\chi_\ell^K,\\
 N_{w,j}&=\pi_j\sum_\ell c_\ell\zeta_{\ell,w}^{\,j}
                     \chi_{\ell,-w}^{\,K-j},\qquad
 \pi_j=\binom Kj p^j(1-p)^{K-j},\quad p=1/M,\\
 \chi_{\ell,-w}&=\frac{\chi_\ell-p\zeta_{\ell,w}}{1-p}.
 \label{eq:fourier-marked-powers}
\end{align}
Assume $M>1$; $M=1$ is deterministic. The exclusion is the same exact-count requirement as in F.
The binomial identity gives $\sum_j N_{w,j}=Z$ mode by mode, not only after the full Fourier sum.

Keep a frequency set $\mathcal L$ closed under conjugation and discard the other modes. The resulting
estimates are real but need not be nonnegative. They still obey the common-normalizer identity, but this
alone does not make them count probabilities. Absolute remainders are
\begin{align}
 |Z-\widetilde Z|&\leq
       \sum_{\ell\notin\mathcal L}|c_\ell||\chi_\ell|^K=:\delta_Z,\\
 |N_{w,j}-\widetilde N_{w,j}|&\leq
       \pi_j\sum_{\ell\notin\mathcal L}|c_\ell|
                 |\chi_{\ell,-w}|^{K-j}=:\delta_{w,j}.
 \label{eq:fourier-marked-remainder}
\end{align}
Tail numerators sum the relevant $j$ bounds; event/complement intervals or
\eqref{eq:ratio-error} then control the original count tails. One may also use the exact prior bounds
$0\leq N_{w,j}\leq\pi_j$, since the physical likelihood is at most one.

The distinction between powers is substantive. Unmarked evidence is damped by $K$, but a query with $j$
copies of a fixed word leaves only $K-j$ independent residual draws. At $j=K$ there is no such damping.
Its prior factor $p^K$ is small, but its posterior relevance still depends on the evidence. Accurate evidence
alone does not establish accurate marked inference.

\subsection{What a cheap spectral certificate would need}
Suppose the omitted modes obey $|\chi_\ell|\leq\rho<1$. The exclusion identity gives
\[
 |\chi_{\ell,-w}|\leq\rho_-:=
       \min\{1,(\rho+p)/(1-p)\}.
\]
With $A_{\rm om}=\sum_{\ell\notin\mathcal L}|c_\ell|$,
\begin{equation}
 \delta_Z\leq A_{\rm om}\rho^K,\qquad
 \delta_{w,j}\leq\pi_j A_{\rm om}\rho_-^{K-j}.
 \label{eq:fourier-gap-bound}
\end{equation}
Thus the proposed assumption has two parts: constructible control of the omitted likelihood spectrum and
of the marked prior characteristic functions. A bound on only one of them is incomplete.
If neither is accessible cheaply, writing \eqref{eq:fourier-marked-remainder} has not reduced the work.

There is a universal but potentially poor bound. Finite Parseval gives
$\sum_\ell|c_\ell|^2=N_g^{-1}\sum_tg_y(t)^2\leq1$, hence
\[
 A_{\rm om}\leq\sqrt{N_g-|\mathcal L|}.
\]
This bound needs no omitted coefficients, but exposes the exponential grid dimension. A favourable power
$\rho^K$ must overcome it. Sharper energy bounds or certified structured coefficient access are additional
algebraic problems, not a free property of a Gaussian observation.

Nor are all nonzero frequencies necessarily damped. Equality $|\chi_\ell|=1$ holds exactly when all the
unit phases $\zeta_{\ell,w}$ agree: equivalently,
$\theta_\ell^T(t_w-t_v)\in2\pi\mathbb Z$ for every $w,v$. These modes detect the affine lattice supporting
the signatures. Near equality also gives slow damping. Remove exact signature redundancies before
postulating a useful gap; retaining only frequencies near the origin is not justified by their names.
Finally, $g_y$ on this finite grid is a periodically extended array, which need not join smoothly across its
boundary. Its coefficients cannot simply be replaced by the continuous Gaussian Fourier formula.

\subsection{Count coefficient access as well as retained modes}
Let $R=|\mathcal L|$. Once the $c_\ell$ are available, constructing all $\chi_\ell$ and all marked masses
costs $O(RM(C_{\rm phase}+K))$, where $C_{\rm phase}$ includes waveform/signature access and phase
evaluation. Powers can be accumulated without dividing by $\chi_{\ell,-w}$, which may vanish.
Modes can be streamed with $O(MK)$ accumulated count storage.

The missing cost is often the largest one. Direct access to $R$ likelihood coefficients requires all
$N_g$ likelihood values and $O(RN_g)$ summation, or a full $O(N_g\log N_g)$ transform.
Neither is a scalable way to obtain a few modes when $d$ is large. A short final Fourier sum is therefore
only a conditional reduction until coefficient construction and the omitted-mode certificate are reduced too.
The exact distinct-source formulas instead use cardinality coefficients of
$\prod_w(1+u\zeta_{\ell,w})$ and appropriate count markers. The IID powers and gap claim above must not be
applied to them unchanged.

\subsection{An alternative exact integral avoids the finite-grid coefficient table}
There is a different spectral starting point that exposes another approximation opportunity.
For $\Omega\sim\mathcal N(0,\sigma^{-2}I_n)$, the Gaussian characteristic identity gives
\[
 e^{-\|y-s\|^2/(2\sigma^2)}
   =\E_\Omega e^{i\Omega^T(y-s)}.
\]
Put $\chi(\Omega)=M^{-1}\sum_w e^{-i\Omega^T\phi_w}$. Then, exactly for IID messages,
\[
 Z=\operatorname{Re}\E_\Omega[e^{i\Omega^Ty}\chi(\Omega)^K],
\]
and the marked expression in \eqref{eq:fourier-marked-powers} holds inside the same expectation with
$\zeta_w=e^{-i\Omega^T\phi_w}$. This uses the physical dimension $n$ and does not construct
$(2K+1)^d$ coefficients. It still requires the characteristic sum or its own justified approximation.

For any count event $A$, its integrand is
$\operatorname{Re}\{e^{i\Omega^Ty}\E_\pi[\ind_Ae^{-i\Omega^T\Phi a}]\}$ and lies in
$[-\pi(A),\pi(A)]$. Independent frequency samples therefore give, for a fixed family of $J$ marked queries
plus the evidence, simultaneous probability at least $1-\delta$ bounds
\begin{equation}
 |\widehat N_A-N_A|\leq\pi(A)
       \sqrt{\frac{2\log[2(J+1)/\delta]}{S}},\qquad
 |\widehat Z-Z|\leq\sqrt{\frac{2\log[2(J+1)/\delta]}{S}}.
 \label{eq:random-frequency-bound}
\end{equation}
These follow from the bounded-variable exponential inequality and a union bound. All frequencies can be
shared between queries. Ratios still require a positive evidence lower bound; the estimates can be signed.
For all word counts, direct work is $O(SM(C_{\rm phase}+K))$. No large frequency grid is needed, but
explicit alphabet access and cancellation may remain prohibitive. This construction is distinct from
sampling \emph{legal configurations} in the next entry, which preserves a positive approximate posterior.

\clearpage
\section{Entry H: positive weighted samples for count expectations}
\label{sec:weighted-sampling}
\subsection{Specify a drawable proposal, not a posterior-sampling oracle}
For the fixed observation $y$, let $q(a)>0$ on every legal configuration of positive target weight.
It must be possible both to draw independent configurations from $q$ and to evaluate $q(a)$.
Define the unnormalized importance weight
\[
 W(a)=\frac{\omega(a)}{q(a)},\qquad
 0<W(a)\leq W_{\max},\qquad \E_qW=Z.
\]
The upper bound must be justified over the legal support; the largest sampled weight is not such a bound.
One always specified proposal is the actual source prior $q=\pi$, for which $W$ is the Gaussian likelihood
and $W_{\max}=1$ is valid. Drawing the prior requires no knowledge of the transmitted messages.
It may be an extremely poor proposal for the posterior.

For independent draws $a^{(1)},\ldots,a^{(S)}$, form
\[
 \widehat Z=\frac1S\sum_sW_s,\qquad
 \widehat N_A=\frac1S\sum_sW_s\ind\{a^{(s)}\in A\},\qquad
 \widehat P=\frac{\sum_sW_s\delta_{a^{(s)}}}{\sum_sW_s}.
\]
All samples are legal and preserve the known total $K$ and the declared source exclusions. The count tails
of $\widehat P$ are nonnegative, monotone and jointly consistent. Self-normalized ratios are not exactly
unbiased at finite $S$; the following bounds control their error directly. Probabilities of these bounds
refer to the numerical sampling randomness conditional on $y$.

\subsection{A computable simultaneous count-tail certificate}
For a fixed family of $J$ count events, the variables $W$ and $W\ind_A$ lie in $[0,W_{\max}]$.
The bounded-variable exponential bound gives, with probability at least $1-\delta$, all of
\[
 |\widehat Z-Z|\leq b_S,\qquad
 |\widehat N_A-N_A|\leq b_S,\qquad
 b_S=W_{\max}\sqrt{\frac{\log[2(J+1)/\delta]}{2S}}.
\]
For completeness, the bound follows by applying
$\E e^{t(X-\E X)}\leq e^{t^2W_{\max}^2/8}$ to each bounded variable, optimizing $t$, and union-bounding
the two tails for $J+1$ quantities. If $\widehat Z>b_S$, the ratio calculation yields
\begin{equation}
 |\widehat P(A)-P_y(A)|
 \leq \min\!\left\{1,\frac{b_S[1+\widehat P(A)]}{\widehat Z-b_S}\right\}
 \leq e_S:=\min\!\left\{1,\frac{2b_S}{\widehat Z-b_S}\right\}.
 \label{eq:importance-computable}
\end{equation}
When the denominator cannot be certified positive, use the trivial bound one.
Taking $J=MK$ covers all message count tails, including words never drawn. Their empirical tails are zero,
but they have not been declared impossible. The Bayes action for $\widehat P$ consequently has conditional
excess loss at most $\min(1,2e_S)$ on the same simultaneous event.

This can work even when the sample-supported law misses most individual posterior configurations.
Uniformly accurate values of the required count events are a weaker requirement than small TV on the full
count-vector space. Applying a full-posterior support-mass bound here can therefore be unnecessarily severe.
It remains essential to cover the events of \emph{unseen words}, not just those with nonzero empirical scores.

\subsection{The informative parameter is importance-weight concentration}
The preceding computable certificate is conservative. A variance argument reveals the more relevant sample
scale. Put $\gamma=W_{\max}/Z\geq1$. For $H=\ind_A$ and $T=P_y(A)$,
\[
 Y=W(H-T),\quad \E_qY=0,\quad |Y|\leq W_{\max},\quad
 \E_qY^2\leq W_{\max}Z\,T(1-T).
\]
The last inequality follows by replacing one of the two $W$ factors by $W_{\max}$ and averaging the other
under $P_y$. Also $\operatorname{Var}_q(W)\leq W_{\max}Z$.
For independent centered variables bounded by $b$ with variance at most $v$, a useful exponential bound is
\[
 \left|\frac1S\sum_s X_s\right|
 \leq\sqrt{\frac{2vt}{S}}+\frac{2bt}{3S}
 \quad\text{with failure probability at most }2e^{-t}.
\]
One obtains it by bounding the centered moment-generating series with
$\exp[t_0^2v/(2(1-|t_0|b/3))]$ for $|t_0|b<3$, then optimizing the exponential tail bound.
Thus no assumed Gaussian distribution of sampling error is needed.

Use $t=\log[2(J+1)/\delta]$ for the denominator and the $J$ centered observables above.
If $S\geq16\gamma t$, the denominator bound gives $\widehat Z\geq Z/2$ on this event. Dividing the
bound on the sample mean of $Y$ by $\widehat Z$, and using $T(1-T)\leq1/4$, proves
\begin{equation}
 \max_A|\widehat P(A)-P_y(A)|
 \leq \min\!\left\{1,\sqrt{\frac{2\gamma t}{S}}+
                         \frac{4\gamma t}{3S}\right\}.
 \label{eq:importance-concentration}
\end{equation}
Hence a sufficient nontrivial sample scale is of order $\gamma\log(J/\delta)/e^2$, not necessarily
$\gamma^2/e^2$. This is a conditional rate statement: $Z$ or a useful deterministic lower bound need not
be available. Equation \eqref{eq:importance-computable} supplies a certificate from the sampled sums when
its own conditions hold. Arbitrarily rescaling all weights changes neither $\gamma$ nor the normalized
estimate. The difficulty is proposal mismatch, not a chosen common likelihood normalization.

\subsection{A complete sampled decoder need not enumerate the alphabet}
For the prior proposal, draw $K$ messages according to the source, construct their sum, compute one
likelihood, and record their sparse counts. In a sample, a word with count $c$ contributes its weight to
tails $1,\ldots,c$; all updates together number exactly $K$. Across $S$ samples there are at most $SK$
positive word/tail entries. Accumulating them, normalizing, and selecting the best $K$ occurrences defines
a complete approximate Bayes action without an $M$-entry array.

Let $C_q$ include drawing a legal configuration and evaluating its proposal probability, and let
$C_{\rm enc}$ construct one waveform. A direct cost bound is
\[
 O\!\left(S[C_q+K C_{\rm enc}+nK]+SK\log(2+SK)\right).
\]
Sparse tail storage is $O(SK)$ plus one signal/count record. For IID prior draws, word generation costs
$O(KB)$ per configuration. Distinct prior draws must actually sample without replacement and retain their
source law. The confidence factor $\log(MK)$ can be calculated from $B,K$ without listing all messages.
Thus this is a complete construction with no explicit full-alphabet enumeration. It is \emph{not} a
polynomial-complexity guarantee: for a poor proposal $\gamma$ and the required sample count can be enormous.
Weight arithmetic and certified numerical rounding remain additional costs.

\subsection{Use a positive approximate decoder as a proposal and correct it}
A previous approximation can supply a more informed proposal. Suppose its legal law is
$Q=\widetilde\omega/\widetilde Z$ and a certified ratio satisfies
$u\leq\omega/\widetilde\omega\leq v$. With $q=Q$,
\[
 W(a)=\widetilde Z\,\frac{\omega(a)}{\widetilde\omega(a)},\qquad
 W_{\max}=v\widetilde Z,\qquad
 \gamma\leq v/u.
\]
For a log-ratio range $\Omega$, this gives $\gamma\leq e^\Omega$. Sampling from the surrogate and using
the exact importance weights therefore corrects its model of the posterior, subject to the finite-sample
bound. The surrogate is a computational proposal, not an additional physical noise model.

The positive finite state graph in B and the cardinality-constrained forest in C/D admit exact backward
sampling once their forward tables exist. Those tables, their normalizer and evaluation of a sampled
configuration's approximate weight must all be charged. The construction does not bypass their earlier
explicit-$M$ preprocessing. The separate per-word moment estimates in F are not automatically a drawable,
evaluable joint proposal.

If the surrogate is supported only on a retained event $E$, it can correct inference only to $P_y(\cdot\mid E)$.
An omitted-mass bound is still needed for unconditional tails; a conditional tail error $e$ gives unconditional
error at most $\eta+(1-\eta)e$. Alternatively a proposal mixture with a positive prior component can
restore full support, but its density and weight bound must be evaluated for the actual mixture.
A per-word prior mixture and a whole-configuration prior mixture have different lower bounds; they are
not interchangeable shortcuts.

\subsection{Integrate a manageable block before sampling the rest}
There is also a variance reduction that uses exact local algebra. Write a legal configuration as $(u,v)$,
and let $q(u,v)=q_U(u)q(v\mid u)$. If the $v$ block can be summed exactly, define
\[
 Z(u)=\sum_v\omega(u,v),\qquad
 N_A(u)=\sum_v\omega(u,v)\ind_A,\qquad
 \overline W(u)=Z(u)/q_U(u),\quad
 \overline W_A(u)=N_A(u)/q_U(u).
\]
These are precisely the conditional expectations of $W$ and $W\ind_A$ under $q(v\mid u)$.
Sampling only $u$ therefore preserves their expectations and cannot increase their variances.
More directly, the centered numerator for a normalized query is the conditional expectation of
$W(\ind_A-T)$, so its variance also cannot increase. This does not prove that finite-sample ratio mean-square
error decreases in every case, nor that computing the conditional sums is cheap.

Local polar sums, small exact regions or an exact component contraction could supply such a block.
They must retain the exact source constraints and count marks, and their construction/query cost belongs
in $C_q$ and the per-sample work. Exact summation over one unknown user can still scan the full alphabet.
This is a concrete way to combine algebraic elimination with stochastic approximation, not an oracle.

\subsection{Adaptive stopping requires simultaneous bounds in time as well}
The displayed guarantees concern a fixed sample count and a fixed family of events. For prescribed rounds
$S_1<S_2<\cdots$, give round $r$ failure budget $\delta_r=\delta/[r(r+1)]$.
A union bound makes every round's event valid simultaneously with total failure probability at most $\delta$.
It is then safe to stop when a decision certificate first succeeds. With a fixed proposal, rounds may reuse
the existing draws. If a proposal is changed using earlier samples, the stated argument applies to fresh
independent draws conditional on that history, with a newly justified weight bound.
Reusing arbitrary adaptive draws or substituting correlated chain samples needs a separate martingale or
dependence argument. No mixing time or effective sample count has been assumed here.

\clearpage
\section{Entry I: spend accuracy on the prefixes that can change the decision}
\label{sec:prefix-decisions}
\subsection{A prefix count bounds every allocation inside that prefix}
Most previous deterministic constructions still compute $MK$ scores. This is not intrinsically required
by the Bayes objective. Partition the message alphabet into a frontier of disjoint prefix cells
$\mathcal C$ covering all messages. A cell may be one full word or a large set of unresolved words.
For a cell $C$, define
\[
 A_C=\sum_{w\in C}a_w,\qquad
 T_{C,j}=P_y(A_C\geq j),\qquad
 G_C(h)=\E_{P_y}\min(A_C,h)=\sum_{j=1}^hT_{C,j}.
\]
These are posterior counts of transmitted copies in the cell, not probabilities that every proposed word
is correct. Any actual output assigning $h$ copies to words inside $C$ satisfies, pointwise,
\begin{equation}
 \sum_{w\in C}\min(a_w,d_w)\leq\min(A_C,h),\qquad h=\sum_{w\in C}d_w.
 \label{eq:prefix-gain-bound}
\end{equation}
It cannot match more than either the true count or its output budget there. Thus $G_C(h)$ is an optimistic
gain bound even before discovering the words. For a singleton it is the exact gain function.
An unresolved prefix is never substituted for an actual decoded message.

Suppose simultaneous intervals $L_{C,j}\leq T_{C,j}\leq U_{C,j}$ are available.
Intersect them with $[0,1]$, then tighten by count-tail monotonicity:
\[
 L'_{C,j}=\max_{\ell\geq j}L_{C,\ell},\qquad
 U'_{C,j}=\min_{\ell\leq j}U_{C,\ell}.
\]
Both sequences are nonincreasing and still contain the true tails. This costs $O(K)$ per cell.
If a signed method needs approximate scores, their midpoints are monotone tails with error at most
half the largest interval width. This is a justified marginal construction; it does not assert a common
joint posterior. Inconsistent intervals instead indicate an invalid bound or arithmetic issue.

\subsection{A global decision certificate without all leaf scores}
For the frontier, form the optimistic gain
\[
 U_*=\max_{\substack{h_C\in\mathbb Z_+\\\sum_C h_C=K}}
       \sum_C\sum_{j=1}^{h_C}U'_{C,j}.
\]
The true optimum is no larger, by \eqref{eq:prefix-gain-bound}. Since the upper sequences are nonincreasing,
the same largest-next-tail allocation computes $U_*$. Let $d$ be any feasible actual word allocation
and let $L(d)=\sum_w\sum_{j=1}^{d_w}L'_{\{w\},j}$ be its lower gain bound.
Include every word used by $d$ as a resolved singleton, splitting its ancestors if necessary. Then
\begin{equation}
 0\leq R_T(d)-R_T(d^*)\leq
       \min\{1,[U_*-L(d)]/K\},\qquad
 R_T(d)\leq1-L(d)/K.
 \label{eq:prefix-regret}
\end{equation}
The second bound distinguishes absolute posterior risk from closeness to Bayes: if even Bayes is poor,
a small excess-loss certificate alone is not strong recovery.

This is not a union bound on all undiscovered words. One aggregate random count controls every allocation
inside a prefix, with the common output budget $K$ retained. If a prefix's first upper tail is already below
all relevant selected lower tails, none of its later tails can force an improving exchange.
More generally \eqref{eq:prefix-regret} controls the entire optimistic allocation at once.

A complete refinement calculation can now be stated:
\begin{enumerate}
\item Start from the root cell; its total count is exactly $K$. Obtain cell intervals from a specified
      provider and split prefixes to discover actual candidate words.
\item Maintain a feasible incumbent word allocation and its lower gain. Optimize the frontier upper gain
      using the monotone cell sequences.
\item Refine a cell or a query accuracy that contributes to the remaining gain gap. Splitting a cell keeps
      both children, so unvisited messages remain covered.
\item Stop when \eqref{eq:prefix-regret} meets the desired tolerance, or return the incumbent and the
      unresolved bound at the declared resource limit.
\end{enumerate}
Keeping the best incumbent lower bound and the smallest valid global upper bound seen so far makes the
reported gap nonincreasing. This does not require successive approximate providers to define the same law;
their intervals must cover the same exact target. A fixed inaccurate provider may never close the gap,
even after all leaves are visited. Further accuracy refinement or a finite exact fallback is then needed.

\subsection{Moments can answer prefix queries without enumerating their payloads}
Here is a concrete provider extending F. For a nonempty proper cell $C$, put $p_C=|C|/M$.
Let $X_C,X_{\bar C}$ be independent prior draws conditional on being inside and outside the cell.
Under IID messages,
\[
 N_{C,j}=\binom Kj p_C^j(1-p_C)^{K-j}
 \E\exp\!\left[-\frac{\|y-\sum_{i=1}^jX_{C,i}
                          -\sum_{i=1}^{K-j}X_{\bar C,i}\|^2}{2\sigma^2}\right].
\]
Approximate the two conditional sums jointly by their matched Gaussian. Its mean and covariance are
\[
 m_j=j\mu_C+(K-j)\mu_{\bar C},\qquad
 V_j=j\Sigma_C+(K-j)\Sigma_{\bar C}.
\]
The likelihood formula is \eqref{eq:gaussian-residual} with $(h\mu_w,h\Sigma_w)$ replaced by $(m_j,V_j)$.
The independent-increment proof of F, which need not use identically distributed increments, gives
\[
 \delta_{C,j}\leq
 \min\!\left\{1,\frac{j\,b_C+(K-j)b_{\bar C}}{2\sigma^3}\right\},\quad
 b_C=\E\|X_C-\mu_C\|^3+
     \E\|\mathcal N(0,\Sigma_C)\|^3.
\]
Use trace bounds as in F, multiply by the binomial prior masses, then apply the event/complement
normalization bounds. This supplies valid cell-tail intervals even when separate approximate cell laws are
inconsistent. Root and empty-cell counts are known exactly and require no Gaussian approximation.

The polar encoder gives an exact way to construct useful cell moments. Order prefix bits as pattern bits
followed by payload bits. Once the pattern is fixed, a payload prefix leaves a uniform affine set
$v=v_0+Rz$ over $\mathbb F_2$, where $z$ comprises its free bits. Regard generator rows as column vectors
$\ell_t$. Character averages are
\begin{align}
 \E(-1)^{\ell_t^Tv}
   &=(-1)^{\ell_t^Tv_0}\ind\{R^T\ell_t=0\},\\
 \E(-1)^{(\ell_t+\ell_u)^Tv}
   &=(-1)^{(\ell_t+\ell_u)^Tv_0}\ind\{R^T(\ell_t+\ell_u)=0\}.
 \label{eq:prefix-character-moments}
\end{align}
Indeed, a nonconstant character sums to zero over the free bits. Multiply these moments by the fixed
amplitudes and insertion matrix exactly as in F. A prefix of pattern bits instead averages the full-payload
moments over its allowed patterns; that pattern work remains charged.
With global mean $\mu$ and second moment $R_2$, complement moments follow from
\[
 \mu_{\bar C}=\frac{\mu-p_C\mu_C}{1-p_C},\qquad
 R_{2,\bar C}=\frac{R_2-p_CR_{2,C}}{1-p_C},\qquad
 \Sigma_{\bar C}=R_{2,\bar C}-\mu_{\bar C}\mu_{\bar C}^T.
\]
Thus fixing payload bits need not cause a sum over all remaining payloads.
For general cells the covariance changes by more than a rank-one update, so F's cheap leaf-query matrix
reuse does not automatically apply. Direct work is up to $O(Kn^3)$ per visited cell after its moments,
plus their construction and $O(K)$ scalar interval work. Conditional-group mixtures from F can refine a
poor Gaussian bound, with their stated sector cost. Distinct-source prefix queries need a
without-replacement derivation; the IID binomial formulas must not replace that source law.

\subsection{Fourier and sample providers fit the same query interface}
In G, replace a single-word phase by the unnormalized cell characteristic
\[
 s_C(\theta)=\frac1M\sum_{w\in C}\zeta_w(\theta),\qquad
 s_{\bar C}(\theta)=\chi(\theta)-s_C(\theta).
\]
The count-$j$ transform is exactly
$\binom Kj s_C(\theta)^j s_{\bar C}(\theta)^{K-j}$.
Finite mode or random-frequency bounds can therefore target prefixes directly. Computing $s_C$ is itself
a nonlinear character sum and is not automatically as cheap as the first two moments in
\eqref{eq:prefix-character-moments}. A short query list does not remove this coefficient-access bottleneck.

For H, accumulate $A_C$ from each sampled configuration. At each bit depth, the occupied prefixes have
counts summing to $K$, so all their count-tail updates require $O(K)$ work per sample and depth.
All sampled prefix statistics can thus be formed in $O(SKB)$ work and at most $O(SKB)$ sparse storage.
Unsampled cells have zero empirical tails but retain their confidence upper bounds.
A binary message tree has at most $2M-1$ cells. Setting $J=K(2M-1)$ in the simultaneous bound covers
\emph{all} possible prefix queries, even if the algorithm chooses which to inspect after seeing the samples.
The logarithm of that number is available from $B$; no tree enumeration is required for the union bound.
Adaptive sample counts use H's round-wise confidence allocation.

With $V_{\rm visit}$ queried cells and frontier size $F_{\rm front}$, maintaining monotone bounds costs
$O(KV_{\rm visit})$ besides the provider work. One optimistic allocation can use a heap in
$O(F_{\rm front}+K\log(2+F_{\rm front}))$ time and $O(KF_{\rm front})$ stored intervals.
The number of refinements and repeated provider work must be counted; neither is guaranteed small.
In the worst case the tree still visits all $M$ leaves or its intervals remain uninformative.
The positive result is narrower: full-message readout can be driven by certified gain bounds, and both
moment and sample providers have been derived without assuming an oracle for those bounds.

\clearpage
\section{Entry J: retain cycles through joint features and conditioning}
\label{sec:cyclic-features}
\subsection{A feature expansion supplies a variable that can be conditioned on}
Eliminating the triangle's middle message gives the joint factor $\sum_v\nu_v b(u,v)b(v,t)$.
This generally differs from the product of its two separate averages. Entry D offers a way to retain
such dependence with a different intermediate domain. On C's regional graph, write each chosen positive
edge approximation as
\[
 \widetilde b_e(a_r,a_q)=\sum_{\lambda_e=1}^{R_e}
              u_{e,\lambda_e}(a_r)v_{e,\lambda_e}(a_q).
\]
The feature index is an auxiliary summation variable. When every retained edge has this form, expanding
them gives the exact identity for the \emph{specified surrogate}
\begin{equation}
 \widetilde Z_K=[z^K]\sum_{\lambda}
       \prod_r\left[\sum_{a_r}f_r(a_r)z^{k_r}
                       \prod_{e\ni r}h_{r,e,\lambda_e}(a_r)\right].
 \label{eq:feature-network}
\end{equation}
Here $h$ denotes the appropriate endpoint feature; common source/likelihood constants are as in C.
All feature indices touching $r$ depend on the \emph{same} local count state $a_r$. Summing that state
creates a joint tensor on its incident indices. Replacing it by separate edge marginals would introduce
another approximation. This is the same obstruction as Report 7's union of overlap-parity requirements.

\subsection{Break cycles by summing a small feature cutset}
Choose edges $\mathcal F$ whose removal leaves a forest. Fix their feature indices first. Each fixed edge
then contributes only two unary factors, and the remaining forest is evaluated by C/D, still extracting
total degree $K$. Finally sum the \emph{unnormalized} sector answers. There are
\[
 N_{\mathcal F}=\prod_{e\in\mathcal F}R_e
\]
sectors. A single cycle needs only one cut edge: the multiplicative cost is its feature count, rather than
automatically a product of all physical state sizes. No closing edge has been discarded beyond the already
bounded feature approximation. Multiple cycles can make $N_{\mathcal F}$ prohibitive.

Let $C_{\rm forest}$ include the remaining forest's two passes and count readout in
\eqref{eq:feature-cost}. One direct bound, after local states and features have been constructed, is
\begin{equation}
 O\!\left(N_{\mathcal F}
       \left[C_{\rm forest}+\sum_r\deg_{\mathcal F}(r)D_r\right]\right).
 \label{eq:cycle-cost}
\end{equation}
The extra term forms sector-specific unaries. Sectors can be streamed with one forest workspace and
accumulated count numerators; no $N_{\mathcal F}$ copies of all tables are required. A zero-evidence sector
contributes zero and is never normalized. Nonzero sectors contribute their evidence times their conditional
count probabilities. Averaging their normalized answers uniformly would be incorrect.

For subsequent exact sampling from this surrogate, select sectors with their evidence weights and sample
the conditional forest. Store the sector weights or pay to recompute them; forest tables for an unrecorded
chosen sector also need rebuilding. Evaluate the marginal configuration density
$Q(a)\propto\prod_r f_r(a_r)\prod_e\widetilde b_e(a_r,a_q)\ind\{\sum k_r=K\}$.
Using the probability of the particular sampled auxiliary feature assignment as though it were $Q(a)$
would give wrong importance weights.

An alternative is to condition on regional states whose removal makes a forest. Its sector count is the
product of their $D_r$ values, with the remaining occupancy reduced by their known counts. The edges among
conditioned states and every source constraint remain in the sector weight. One can mix state conditioning
and feature conditioning, or use a general elimination schedule with joint feature indices. Compare the
complete costs of these schedules; a small isolated $R_e$ does not bound their largest joint intermediate.

\subsection{Use the exact polar interface before adding a feature approximation}
In a conditional slot block, a payload influences all its neighbours through the union of their overlap
parities. Group using that union, then form the required weighted sums or feature tables. This removes
exact redundancy before approximating anything. A full-message marker must distinguish its own weight
inside a group, exactly as in Report 7's marked triangle. Unknown patterns remain variables to sum or sample.
The selected-bank census already shows that a union can expose all payload bits; no small interface is
assumed here.

The full retained-support error is still the sum of local log-ratio ranges from C/D. Conditioning and exact
sector summation add no approximation error, although they add work. This is a concrete extension of the
positive contraction decoder to cyclic graphs. Its favourable regime requires both affordable local
states/features and an affordable joint schedule; it is not supplied by pairwise rank alone.

\section{Entry K: measure the correction where the reference places probability}
\label{sec:weighted-corrections}
\subsection{A second-moment certificate follows from the exact closure identity}
On the same legal support let $Q=\widetilde\omega/\widetilde Z$ and
$R=\omega/\widetilde\omega>0$. Report 7's exact reweighting identity is
\[
 P_y=\frac{QR}{\E_QR},\qquad
 \E_{P_y}H-\E_QH=\frac{\operatorname{Cov}_Q(H,R)}{\E_QR}.
\]
Define the scale-invariant correction variance
\[
 v_R=\frac{\E_QR^2}{(\E_QR)^2}-1.
\]
Cauchy--Schwarz immediately gives
\begin{equation}
 \TV(P_y,Q)\leq\min\{1,\tfrac12\sqrt{v_R}\},\qquad
 |\E_{P_y}H-\E_QH|\leq\sqrt{\operatorname{Var}_Q(H)\,v_R}.
 \label{eq:weighted-correction-bound}
\end{equation}
The first follows by writing TV as $\frac12\E_Q|R/\E_QR-1|$.
For a tail event the second contains $Q(A)(1-Q(A))$. For a competing-copy contrast
$H=\ind_A-\ind_B$, its variance is at most $Q(A\mathbin\triangle B)$: configurations where both events
agree do not contribute. Thus the required decision can be stable even when a uniform factor range is poor.

This identifies a different favourable assumption, not a free certificate. Computing $\E_QR$ is the
original evidence ratio, and computing $\E_QR^2$ can restore the difficult interactions. A usable bound needs
a cheaper lower bound on the first and an upper bound on the second, from positive envelopes, a compatible
conditional contraction, or separately justified sampling confidence. For example, bounds $m_->0$ and
$s_+\geq\E_QR^2$ give $v_R\leq s_+/m_-^2-1$ when $m_-\leq\E_QR$.
Their construction belongs in the cost. An empirical variance without a remainder or confidence bound
does not certify an unseen large correction.

If several modifications are made, use their complete product $R=\prod_e R_e$.
Zero covariance with each $R_e$ separately need not imply zero covariance with their product; Report 7
gives an explicit two-factor counterexample. Joint feature conditioning from J can evaluate small coupled
correction blocks without assuming those blocks independent of their boundary.

\subsection{A rare bad region is useful only if reweighting cannot amplify it too much}
Let $G$ be a region where $r_-\leq R\leq r_+$, with $r_->0$, and suppose
$Q(G^c)=q$ and $R\leq b$ on $G^c$. Positivity yields
\[
 P_y(G^c)\leq\eta_*:=\frac{qb}{(1-q)r_-+qb}.
\]
For $0<q<1$, set $m=\max(q,\eta_*)$ and $\Omega_G=\log(r_+/r_-)$. Comparing the masses of $G,G^c$
and then their conditional laws gives
\begin{equation}
 \TV(P_y,Q)\leq m+(1-m)\tanh(\Omega_G/4).
 \label{eq:good-region-correction}
\end{equation}
Indeed, the mismatch of region masses and arbitrary disagreement inside $G^c$ cost at most $m$;
the common mass inside $G$ is $1-m$ and uses the range bound. Endpoint cases follow by continuity.
Certified bounds on $q$, rather than its exact value, can be substituted monotonically.
This connects A's mass accounting with D's local accuracy without actually deleting the bad region.
A small $q$ alone is insufficient if $b/r_-$ is enormous. The observation-weighted rarity and possible
likelihood amplification must both be controlled.

\subsection{Known surrogate answers can be control variates for correction}
Suppose $h=\E_QH$ is already available from an exact contraction of $Q$. For any fixed scalar $c$,
\begin{equation}
 \E_{P_y}H=h+
       \frac{\E_Q[(R-c)(H-h)]}{\E_QR}.
 \label{eq:correction-control-variate}
\end{equation}
This follows because $\E_Q(H-h)=0$. Estimate only the displayed correction and denominator from independent
$Q$ samples. The correction sample mean is unbiased; their ratio generally is not. If $R\in[r_-,r_+]$,
choose $c=(r_-+r_+)/2$ and put $d=(r_+-r_-)/2$. For an indicator $H$,
\[
 \E_Q[(R-c)^2(H-h)^2]\leq d^2\operatorname{Var}_Q(H).
\]
This bounds its variance as well. A nearly constant correction can therefore be cheap even when the
uncorrected count statistic is variable. Since $|(R-c)(H-h)|\leq d$, the bounded-variable bounds from H
give finite numerator confidence; the denominator still needs its own lower bound. All queried corrections
and adaptive rounds need a common failure-budget account. The corrected scores may be signed or lose
monotonicity, so use valid intervals and I's tightening, not an assertion of a positive empirical law.

An exact $h$ is essential to the displayed identity. A moment estimate substituted for $h$ needs its own
error propagation. If $h$ is unavailable, ordinary importance sampling remains specified by H.
Control variates are most natural when C/D/J have already paid for the surrogate's count or prefix answers.

Report 7's alias conditionals offer another exact variance reduction. Restore the known source-prior
conditional inside a complete equal-signal class. The correction ratio is then constant within that class,
and $H$ may be replaced there by its prior-weighted conditional average. Conditional expectation
cannot increase its variance. This preserves irreducible message ambiguity; selecting one representative
message would not. Finding and summing the classes is still charged, and partial classes must be explicit
disjoint groups rather than an assumed census of all aliases.

\subsection{Near aliases: retain identity and charge likelihood variation}
The 29 October geometry extension separates a zero distance from a small distance. Consider a declared
finite group $\mathcal C$ of legal configurations, with their exact source masses $\pi(a)$, and replace
their conditional likelihoods by a common value. The resulting conditional law is
$Q_{\mathcal C}(a)=\pi(a)/\sum_{c\in\mathcal C}\pi(c)$. Put $s_a=\Phi a$ and choose a center $s_0$
with a certified radius $\|s_a-s_0\|\leq\rho$. For $r=y-s_0$,
\[
 \log g(y-s_a)-\log g(r)
 =\frac{r^T(s_a-s_0)-\|s_a-s_0\|^2/2}{\sigma^2}.
\]
Hence a valid within-group log-likelihood oscillation bound is
\begin{equation}
 \Omega_{\mathcal C}\leq\frac{2\|r\|\rho+\rho^2}{\sigma^2},\qquad
 \TV(P(\cdot\mid\mathcal C,y),Q_{\mathcal C})
 \leq\tanh(\Omega_{\mathcal C}/4).
 \label{eq:near-alias-conditional}
\end{equation}
This is the earlier likelihood-range bound applied to a geometric group, not a new approximation family.
Every conditional count tail has the same absolute-error bound. If all states share the same outside
$K-m$ copies, the conditional excess multiset-PUPE of the surrogate Bayes action is at most
$2(m/K)\tanh(\Omega_{\mathcal C}/4)$. For each fixed action only the remaining $m$ true copies vary;
the two action comparisons give the factor two. At $\rho=0$ the exact source conditional is recovered.

A small Euclidean radius alone is insufficient uniformly in the observation or as $\sigma^2$ decreases.
Preserving this conditional law does not certify the group's total posterior mass; changing its evidence
requires the separate normalization/support account. Nor may identities be replaced by one representative.
Report 11's affine translation orbit supplies exact membership and selected count-tail queries for one
zero-radius family in polynomial work. General group discovery, radius certification and allocation over
many overlapping compact groups remain charged. Report 8 provides the physical character-energy metric
and a counterexample to using resource balance as a surrogate for this radius.

\clearpage
\section{Entry L: residual-guided proposals and block updates}
\label{sec:guided-decoder}
\subsection{The residual recursion specifies the ideal next branch}
H left the choice of an informative, drawable proposal open. The exact report and F/I now supply one.
Use ordered slots $W=(w_1,\ldots,w_K)$ for this calculation. Under IID uniform messages,
\[
 \omega_{\rm slot}(W)=M^{-K}g\!\left(y-\sum_i\phi(w_i)\right),\qquad
 g(r)=e^{-\|r\|^2/(2\sigma^2)}.
\]
Its evidence is the same $Z$ as in \eqref{eq:target}. Summing the $K!/\prod_w a_w!$ orderings of each
count vector gives exactly the original factorial source weight. A proposal need not be slot-exchangeable;
importance-weighting ordered slots and then recording counts preserves that target.

Let $\xi$ be a prefix of the $KB$ slot bits, and let
$L(\xi)=\E_\pi[g(y-\sum_i\phi(w_i))\mid\xi]$.
If $\pi(b\mid\xi)$ is the prior probability of the next bit, the exact posterior branch is
\begin{equation}
 p(b\mid\xi,y)=
 \frac{\pi(b\mid\xi)L(\xi b)}{\sum_{c=0}^1\pi(c\mid\xi)L(\xi c)}.
 \label{eq:ideal-bit-proposal}
\end{equation}
This is Bayes' rule applied to a partition into two children. Under the IID bit representation,
$\pi(b\mid\xi)=1/2$. Exact branch sampling would produce the exact posterior and constant importance
weight $Z$. The hard operation is evaluating its continuation likelihoods.

Replace those two likelihoods by any specified strictly positive guide $\widetilde L$. Then
\[
 q_g(b\mid\xi)=\frac{\pi(b\mid\xi)\widetilde L(\xi b)}
                         {\sum_c\pi(c\mid\xi)\widetilde L(\xi c)},\qquad
 q_g(W)=\prod_{t=1}^{KB}q_g(b_t\mid\xi_{t-1}).
\]
There are two child queries per bit, rather than an explicit $M$-word normalization at each slot.
Both sampling and evaluating the full proposal density use the same recorded branch probabilities.
All messages remain legal outputs of the fixed encoder. Approximate guides influence where computation
is spent; exact importance weights or acceptance ratios below retain the physical posterior.

\subsection{Conditional polar moments give a concrete positive guide}
After several complete slots have been chosen, subtract their waveforms from $y$ to obtain $r$.
Suppose the current message lies in prefix cell $C$ and there are $h$ further whole messages. The exact
continuation is the likelihood expectation of
$X_C+X_1+\cdots+X_h$, with independent conditional-current and full-prior future draws. F gives
\begin{align}
 m_{C,h}&=\mu_C+h\mu,\qquad V_{C,h}=\Sigma_C+h\Sigma,\\
 \widetilde L(C;r,h)&=
 \frac{\exp[-\tfrac12(r-m_{C,h})^T(\sigma^2I+V_{C,h})^{-1}(r-m_{C,h})]}
      {\sqrt{\det(I+V_{C,h}/\sigma^2)}},\label{eq:guided-likelihood}\\
 |L-\widetilde L|&\leq
 \min\{1,(b_C+h b)/(2\sigma^3)\},\qquad
 b=\E\|X-\mu\|^3+\E\|\mathcal N(0,\Sigma)\|^3.
 \label{eq:guided-remainder}
\end{align}
The guide is strictly positive for $\sigma^2>0$. Prefix moments and their construction are given in I.
Once a message is complete, its conditional covariance is zero and its signal enters the residual exactly.
The remaining covariance describes unresolved prior draws inside this approximation; it is not a changed
physical noise variance. Retaining conditional groups from F supplies a more expensive multimodal guide.

This construction does not assume that a moment guide is accurate. Its explicit smoothing bound may be
vacuous. Nevertheless it specifies a legal, evaluable proposal; correction can use it even then.
Pattern-prefix moments may require averaging many patterns. Payload moments use polar characters, avoiding
enumeration of all payloads, but matrix work and pattern access remain charged.

For an accuracy-dependent proposal guarantee, suppose at every reachable depth-$t$ prefix the two ratios
$L(\xi b)/\widetilde L(\xi b)$ have log range at most $\Omega_t$. Dividing each ratio by its
$q_g$-weighted mean gives
\[
 e^{-\Omega_t}\leq\frac{p(b\mid\xi,y)}{q_g(b\mid\xi)}\leq e^{\Omega_t},\qquad
 \sup_W\frac{P_y^{\rm slot}(W)}{q_g(W)}\leq
                 \exp\!\left(\sum_{t=1}^{KB}\Omega_t\right).
\]
This is a direct importance-concentration bound. Uniform control over reachable prefixes is required;
checking only sampled paths does not establish it. An absolute likelihood remainder becomes a relative
bound only when its corresponding likelihood has a positive, informative lower bound.

\subsection{An independent corrected decoder is now fully specified}
To obtain a known weight bound without assuming the guide works, choose a fixed $0<\epsilon\leq1$ and use
the whole-configuration mixture
\begin{equation}
 q(W)=(1-\epsilon)q_g(W)+\epsilon\pi_{\rm slot}(W),\qquad
 W_{\rm imp}(W)=\frac{\omega_{\rm slot}(W)}{q(W)}\leq\epsilon^{-1}.
 \label{eq:guided-mixture}
\end{equation}
Choose the mixture component once per configuration, then draw its complete ordered tuple. Evaluate the
\emph{mixture} density, including $q_g(W)$ also for a prior draw. This is not a per-bit prior mixture.
The bound follows from $q\geq\epsilon\pi_{\rm slot}$ and $g\leq1$.

Draw independent tuples, accumulate weighted count tails or prefix counts, and allocate $K$ actual message
copies using H/I. H's confidence statements apply unchanged on this ordered probability space, with
$W_{\max}=\epsilon^{-1}$. No unknown count-proposal density or sum over its slot permutations is needed.
Confidence covers all words, including unseen ones; a resource-limited run returns its unresolved bound.
Guide error is not added as posterior bias after exact importance correction. It affects concentration
and therefore the computation needed to reach a useful bound.

Let $C_{\rm init}$ construct the global/conditional moment resources and let $C_L$ bound a child guide query,
including its moment access. A conservative full cost is
\begin{equation}
 O\!\left(C_{\rm init}+S[KB C_L+K C_{\rm enc}+nK]
                         +SK\log(2+SK)\right).
 \label{eq:guided-decoder-cost}
\end{equation}
Direct dense guide work is $O(n^3)$ per query after moments; any reuse must be established separately.
Prefix readout adds $O(SKB)$ sparse work/storage if used. Moment construction can visit $P$ patterns
without visiting $M=PV$ full messages. This is a complete decoder without explicit full-alphabet
enumeration, not a polynomial-in-$B,K$ guarantee: pattern work and the sufficient sample count can still
be enormous. With only the defensive mixture guarantee, $\gamma\leq1/(\epsilon Z)$.

The formulas for the guide's remainder above are IID. Distinct-source sampling must use legal
without-replacement branch probabilities and the matching ordered prior. A positive heuristic guide with
correct source/importance weights could still be used there, but the IID residual error proof would not
carry over. No such without-replacement closure is asserted here.

\subsection{Corrected block updates retain every outside user's effect}
Another use of a drawable full-support $q$ is independent-proposal Metropolis--Hastings. At state $W$, draw
$W'$ from $q$ and accept with probability
\[
 \alpha(W,W')=\min\left\{1,
        \frac{\omega_{\rm slot}(W')q(W)}{\omega_{\rm slot}(W)q(W')}\right\}.
\]
The accepted flow between two states is the smaller of their two opposing proposal flows, so detailed
balance holds for the exact ordered posterior. Its normalizer cancels. This is a standard exact-target
correction, not a claim that a finite chain is already at stationarity.

Local updates can use more of Report 7. Choose a fixed-size slot block $\mathcal B$ and hold all other
messages at their current values. Its conditional likelihood uses
\[
 r_{\mathcal B}=y-\sum_{i\notin\mathcal B}\phi(w_i),\qquad
 g\!\left(r_{\mathcal B}-\sum_{i\in\mathcal B}\phi(w_i)\right).
\]
For two or three slots, the exact K2/K3 operations therefore apply with this residual. The outside users
are neither known truths nor deleted interactions; they are current states of the joint sampler.
Distinct sampling additionally excludes every outside message and within-block repetition.

An exact conditional \emph{joint draw} is needed, not independent draws from block marginals. For K2,
draw the first word with weight $\nu_u(B\nu)_u$ and the second with conditional weight $\nu_vB_{uv}$.
For K3, its first-slot marginal is one third of its expected-count vector; after drawing that word, the
remaining two-slot law has unaries $\nu_vB_{uv}$ and their mutual factor $B$.
These give exact block draws after the relevant contractions. A K3 update still pays the report's repeated
pair-operator cost and explicit alphabet work. Small blocks have not made that work free.

Alternatively use the bitwise guide for a block and accept/reject with its exact conditional target and
forward/reverse proposal densities. The resulting local kernel also preserves the full posterior.
Choose blocks with state-independent probabilities, or include any state-dependent selection probabilities
in the acceptance ratio. Mixing fixed block kernels and global proposals preserves invariance.
Freeze the guide and kernel-selection settings during each chain run for the bounds below.
This is a rigorous role for two- or three-user calculations at large $K$: local moves inside a full joint
law, with no assumption that higher-user effects fade away.

\subsection{A conservative finite-chain bound closes the remaining guarantee}
Suppose a justified uniform bound gives $P_y^{\rm slot}(W)\leq\bar\gamma q(W)$ for every state.
For the independent-proposal kernel,
\[
 q(W')\alpha(W,W')\geq P_y^{\rm slot}(W')/\bar\gamma.
\]
This follows directly by considering which term realizes the minimum. The full transition, including
rejection, therefore contains a $1/\bar\gamma$ copy of the target distribution. This finite-state argument
is the usual independence-sampler minorization; the general result is given by Mengersen and Tweedie
\cite{mengersen-tweedie}.

Use that global kernel with fixed probability $\beta>0$ and any exact-target local block kernel otherwise.
Put $\kappa=\beta/\bar\gamma$. Every transition still contains a $\kappa$ target component. Subtracting
that component and iterating gives, from any starting state,
\[
 \TV(\mathcal L(W_t),P_y^{\rm slot})\leq(1-\kappa)^t\leq e^{-\kappa t}.
\]
For $S$ recorded states spaced $t$ transitions apart, with the first also after $t$ transitions, sequential
coupling bounds their joint distance from $S$ independent target draws by $S e^{-\kappa t}$.
Thus choose
\begin{equation}
 t\geq\kappa^{-1}\log(S/\delta_{\rm mix}),\qquad
 e=\sqrt{\frac{\log(2J/\delta_{\rm stat})}{2S}}.
 \label{eq:block-chain-certificate}
\end{equation}
With probability at least $1-\delta_{\rm mix}-\delta_{\rm stat}$, every one of $J$ empirical tail queries
has error at most $e$. The proof applies the IID bounded-variable bound on the coupled sample and charges
coupling failure separately. Use $J=MK$, or $K(2M-1)$ for adaptive prefix selection, and the usual Bayes
excess-loss bound $\min(1,2e)$. Adaptive stopping needs H's round-wise confidence allocation as well.

A finite bound on $\bar\gamma$ can come from a uniformly controlled guide, H's positive-surrogate ratio,
or $W_{\max}/Z_-$ with a valid evidence lower bound. Even one legal configuration supplies $Z_->0$ in
real arithmetic, but it can be far too small to give an affordable schedule. A confidence-based lower
bound needs its own failure budget. Unweighted chain records then cost $O(St)$ transitions plus their
sparse readout; exact local contractions and proposal evaluation are charged in each transition.
The conservative schedule does not prove a speed advantage over independent importance sampling.
Local blocks may help exploration, but a better mixing claim requires new evidence or analysis.

\clearpage
\section{Entry M: truncate count multiplicity with a controlled moment hierarchy}
\label{sec:factorial-moments}
\subsection{The unexcluded residual is useful when interpreted as a moment}
The exact report distinguished count probabilities from the unexcluded residual formula. That distinction
also supplies another approximation opportunity. Define
\[
 F_{w,r}=\E_\pi\!\left[\binom{a_w}{r}g(y-\Phi a)\right],\qquad
 \mu_{w,r}=F_{w,r}/Z=\E_{P_y}\binom{a_w}{r}.
\]
Under the IID prior, select the $r$ labeled slots forced to equal $w$. The other slots remain unrestricted,
including further copies of $w$. Hence
\begin{equation}
 F_{w,r}=\binom Kr M^{-r}
  \E\,g\!\left(y-r\phi(w)-\sum_{i=1}^{K-r}X_i\right).
 \label{eq:factorial-residual}
\end{equation}
This is a factorial moment, exactly because no word was excluded. It is not $P_y(a_w=r)$.
For a fixed $r$, all words now share the same residual source moments and covariance. F's closure can reuse
the factorization of $\sigma^2I+(K-r)\Sigma$ across all queried words. Its smoothing error is multiplied by
$\binom Kr M^{-r}$ before normalization. G's characteristic function similarly uses the full source power
$\chi^{K-r}$. The residual still contains $K-r$ users; fixing a small number of slots has not solved it by K2.

\subsection{Alternating finite sums give upper and lower count-tail bounds}
For $j\leq m\leq K$, define
\[
 S_{w,j}^{(m)}=\sum_{r=j}^m(-1)^{r-j}\binom{r-1}{j-1}\mu_{w,r}.
\]
Then
\begin{equation}
 \begin{cases}
 S_{w,j}^{(m)}\geq T_{w,j},&m-j\text{ even},\\
 S_{w,j}^{(m)}\leq T_{w,j},&m-j\text{ odd},
 \end{cases}\qquad
 |S_{w,j}^{(m)}-T_{w,j}|\leq\binom m{j-1}\mu_{w,m+1},
 \label{eq:factorial-tail-bounds}
\end{equation}
where $\mu_{w,K+1}=0$. This is a finite combinatorial statement, requiring no asymptotic collision model.
To verify it for a deterministic count $a\geq j$, use
\[
 1=j\binom aj\int_0^1 t^{j-1}(1-t)^{a-j}\,dt.
\]
Expand $(1-t)^{a-j}$ at zero through degree $m-j$. The integrated coefficients are
$(-1)^{r-j}\binom{r-1}{j-1}\binom ar$. Taylor's integral remainder has the sign of the next term and
magnitude no greater than that term, because the remaining power of $1-t$ lies in $[0,1]$.
For $a<j$ both sides vanish; if the polynomial ends before the truncation there is no remainder.
Averaging over the posterior proves \eqref{eq:factorial-tail-bounds}.

For example,
\[
 \mu_{w,1}-\mu_{w,2}\leq T_{w,1}\leq\mu_{w,1},\qquad
 \mu_{w,2}-2\mu_{w,3}\leq T_{w,2}\leq\mu_{w,2}.
\]
More generally, $T_{w,j}\leq\mu_{w,r}/\binom jr$ whenever $1\leq r\leq j$, by bounding
$\binom{a_w}{r}$ on $\{a_w\geq j\}$. Thus a short moment hierarchy also controls tails above its retained
order. Intersect these inequalities with $[0,1]$ and I's monotonicity constraints.
If the moments themselves have errors $\epsilon_{w,r}$, the error in the displayed partial sum is at most
$\sum_{r=j}^m\binom{r-1}{j-1}\epsilon_{w,r}$. Alternating coefficients can amplify numerical or approximation
errors; truncation and moment-estimation error are distinct. Full order recovers the exact tails, but is
not necessarily the most stable representation.

\newpage
\subsection{When posterior means alone have a defensible decision guarantee}
Assume $M\geq K$ and let $d_{\rm mean}$ select one copy of the $K$ largest exact posterior means.
Define
\[
 C_2=\sum_w\mu_{w,2}=\E_{P_y}\sum_w\binom{a_w}{2}.
\]
For any allocation $d$, the optimistic gain
$\overline G(d)=\sum_{w:d_w>0}\E_{P_y}a_w$ is at least its true matching gain $G(d)$.
It is maximized by $d_{\rm mean}$. For that one-copy output,
$\overline G(d_{\rm mean})-G(d_{\rm mean})
 =\sum_{w:d_{{\rm mean},w}=1}\E(a_w-1)_+\leq C_2$.
Consequently
\begin{equation}
 R_T(d_{\rm mean})-R_T(d^*)\leq\min\{1,C_2/K\}.
 \label{eq:mean-collision-regret}
\end{equation}
Uniform mean-estimation error $e$ adds at most $2e$ to this excess-loss bound.
This does not change the loss or source: it states a sufficient condition under which an easier decision
is close to the original optimum. Under distinct sampling $C_2=0$, recovering the exact mean/tail equivalence.
Under IID, a small prior collision rate is insufficient. The elementary bound
$C_2\leq\binom K2/(MZ)$ follows from $g\leq1$, but can be useless when the evidence is small.
An informative posterior collision bound must be derived or certified.

For moments through order $m$, a direct Gaussian evaluation after global moments costs
$O(mn^3+Mmn^2)$ for all word queries, plus waveform access, evidence bounds and count selection.
The $M$ dependence remains unless a discovery/query mechanism avoids the full scan. I's prefix bounds
can supply such discovery, while moment intervals refine only the resolved candidates that matter.
For sample-based H/L readout, exact empirical tails already cost only $K$ updates per sample; replacing
them by a moment hierarchy offers no automatic saving. This entry is most useful when low-order
analytical moment queries are cheaper than the excluded-source count calculation.

\clearpage
\section{Entry N: tilt the residual law before approximating it}
\label{sec:tilted-residual}
\subsection{Move the approximation to the part of the distribution being queried}
F replaces a residual sum by a Gaussian with its prior mean and covariance. A queried residual can be far
from that mean, where a small absolute error gives little control. Instead, first change the reference
distribution exactly. This uses the one-user partition operation from Report 7. Exponential tilting and
saddlepoint methods are classical \cite{daniels}; the identities and finite bounds below are derived for
our Gaussian-smoothed, marked residual queries, without importing an asymptotic accuracy claim.

Allow independent, not necessarily identical, residual waveforms $X_1,\ldots,X_h$. Each has a specified
finite prior: the full alphabet, an excluded-word alphabet, or a conditional prefix alphabet. Define
\[
 S=\sum_{i=1}^hX_i,\quad L(r)=\E g(r-S),\quad
 \psi_i(\theta)=\E e^{\theta^TX_i},\quad
 \kappa(\theta)=\sum_i\log\psi_i(\theta).
\]
The tilted prior of $X_i$ multiplies each original probability by
$e^{\theta^TX_i}/\psi_i(\theta)$. It is a computational reference, not a changed source model.
Writing $\E_\theta$ for the independent product of these tilted laws, completing the square gives
\begin{equation}
 L(r)=A_\theta(r)J_\theta(r),\qquad
 A_\theta(r)=e^{\kappa(\theta)-\theta^Tr+\sigma^2\|\theta\|^2/2},\qquad
 J_\theta(r)=\E_\theta g(r-\sigma^2\theta-S).
 \label{eq:tilted-residual-identity}
\end{equation}
Indeed the change of measure supplies $e^{\kappa-\theta^TS}$, and
$-\theta^TS-\|r-S\|^2/(2\sigma^2)
=-\theta^Tr+\sigma^2\|\theta\|^2/2-\|r-\sigma^2\theta-S\|^2/(2\sigma^2)$.
No likelihood or noise term has been added.

Let $\mu_{i,\theta}=\nabla\log\psi_i(\theta)$ and
$\Sigma_{i,\theta}=\nabla^2\log\psi_i(\theta)$, with
$m_\theta=\sum_i\mu_{i,\theta}$ and $V_\theta=\sum_i\Sigma_{i,\theta}$.
The function $\log A_\theta$ is strongly convex:
\begin{equation}
 \nabla\log A_\theta=m_\theta+\sigma^2\theta-r,\qquad
 \nabla^2\log A_\theta=V_\theta+\sigma^2I\succ0.
 \label{eq:tilt-stationary}
\end{equation}
Finite bounded alphabets and the quadratic term make it coercive, so its unique minimizer $\theta_*$
satisfies $m_{\theta_*}=r-\sigma^2\theta_*$. At this point the new residual query is at the new sum's mean.
An approximate optimizer is also allowed: \eqref{eq:tilted-residual-identity} holds at every $\theta$.
There is no optimizer error to add if the actual $\theta$, shifted residual and moments are used consistently.

\subsection{A relative closure certificate with the small common scale removed}
Replace the tilted sum by its matching Gaussian, giving
\begin{align}
 \widetilde J_\theta(r)
 &=c_\theta
  \exp\!\left[-\tfrac12(r-\sigma^2\theta-m_\theta)^T
        (\sigma^2I+V_\theta)^{-1}(r-\sigma^2\theta-m_\theta)\right],
 \quad c_\theta=\det(I+V_\theta/\sigma^2)^{-1/2},\\
 \widetilde L_\theta(r)&=A_\theta(r)\widetilde J_\theta(r).
 \label{eq:tilted-gaussian}
\end{align}
For $\xi_{i,\theta}=X_i-\mu_{i,\theta}$ under the tilted law and
$G_{i,\theta}\sim\mathcal N(0,\Sigma_{i,\theta})$, F's replacement proof gives
\[
 |J_\theta-\widetilde J_\theta|\leq
 \delta_\theta:=
 \min\!\left\{1,\frac1{2\sigma^3}\sum_i
       \big[\E_\theta\|\xi_{i,\theta}\|^3+\E\|G_{i,\theta}\|^3\big]\right\}.
\]
Its shaped version also applies, using the covariances of the unreplaced Gaussian increments.
Consequently, whenever $\delta_\theta<\widetilde J_\theta$,
\begin{equation}
 \frac{|L-\widetilde L_\theta|}{L}
 \leq\frac{\delta_\theta}{\widetilde J_\theta-\delta_\theta}.
 \label{eq:tilted-relative-error}
\end{equation}
At $\theta_*$ the denominator uses $c_{\theta_*}$. The potentially tiny factor $A_\theta$ has cancelled.
This is the precise possible benefit; it does not prove that $\delta_\theta/c_\theta$ is small.
The determinant can itself be small, particularly in many dimensions.

There is also a strictly positive lower bound even when the smoothing remainder is unhelpful.
Jensen's inequality for $e^{-x}$ gives
\begin{equation}
 A_\theta
 e^{-[\|r-\sigma^2\theta-m_\theta\|^2+\operatorname{tr}V_\theta]/(2\sigma^2)}
 \ \leq L(r)\ \leq A_\theta.
 \label{eq:tilted-jensen}
\end{equation}
Intersect these bounds with $[0,1]$ and
$A_\theta[\max(0,\widetilde J_\theta-\delta_\theta),
            \min(1,\widetilde J_\theta+\delta_\theta)]$.
At zero tilt this includes an ordinary prior-moment Jensen lower bound. At the optimum its residual
misfit vanishes. Such lower bounds can support F/G/H's normalization certificates, even when no Gaussian
approximation is accepted.

\subsection{The tilted bound refines the exact support bound}
Let $\mathcal C$ be the convex hull of the possible sums $S$ and
$H_{\mathcal C}(\theta)=\max_{s\in\mathcal C}\theta^Ts$.
Since $\kappa(\theta)=\log\E e^{\theta^TS}\leq H_{\mathcal C}(\theta)$,
\begin{equation}
 L(r)\leq A_{\theta_*}(r)
 \leq\inf_\theta e^{H_{\mathcal C}(\theta)-\theta^Tr+\sigma^2\|\theta\|^2/2}
 =e^{-\operatorname{dist}(r,\mathcal C)^2/(2\sigma^2)}.
 \label{eq:tilt-versus-support}
\end{equation}
The final equality is Report 7's support dual with $u=\sigma^2\theta$.
The log moment-generating function retains prior probabilities which the support maximum discards.
Thus exact optimization of the tilted upper bound is no weaker than the convex-hull likelihood envelope.
A finite trial $\theta$ still gives a valid bound, but need not improve the optimized geometric bound.
For an unfinished branch, multiply by its own prior mass before combining it with other branches.

\subsection{Where the gain is paid for, and which combinations survive}
For the full uniform polar alphabet, put $A_p=S_p\diag(\sqrt{e_p})$ and
$Z_{\rm pol}(\lambda)=\sum_v e^{\lambda^T(-1)^{Lv}/2}$. Then
\[
 \psi(\theta)=\frac1M\sum_{p,v}e^{\theta^T\phi(p,v)}
 =\frac1M\sum_p Z_{\rm pol}(2A_p^T\theta).
\]
Its derivatives give the tilted moments. At $\theta=0$, uniform-payload character orthogonality makes
F/I's moments inexpensive. At nonzero tilt that orthogonality shortcut generally disappears. Exact polar
partition sums, their derivatives, all unknown patterns and any exclusions must now be computed or bounded.
SC or a short candidate list is not an exact partition oracle. Approximate partition/moment operations
need their own valid bounds; a covariance guessed independently of a partition is not automatically its Hessian.

If one tilted partition/moment evaluation costs $C_{\rm tilt}$ and $q$ optimization evaluations are used,
a conservative query cost is $O(q(C_{\rm tilt}+n^3))$, plus third-moment bounds and source construction.
Direct summation can cost $O(Mn^2)$ per evaluation before exploiting structure. This can undo F's original
payload savings. The approach is attractive only if improved concentration or tighter query bounds repay it.

There is a useful reuse choice: for M's order-$r$ factorial queries, all words share the same unrestricted
$K-r$-user source. Fixing one $\theta$ for that order reuses its tilted moments and matrix factorization
across residuals $y-r\phi(w)$. Optimizing a separate tilt for every word may tighten the bound but loses that
reuse. F's conditional mixtures can instead tilt each retained occupancy sector, with the sector's original
prior weight restored. L can use a positive tilted likelihood as a guide; exact importance correction then
targets the original law, while guide quality controls concentration. I can use the certified intervals
directly. These are different uses of the same calculation, not interchangeable guarantees.

For a Gaussian increment law the closure is exact at every tilt. For discrete increments with small physical
noise, tilting need not remove separated modes or lattice effects; F's discrete counterexample is not ruled
out by a change of measure. Distinct sampling is dependent and does not admit the displayed product
$\kappa=\sum_i\log\psi_i$ unchanged. The exact identity can use a joint sum's log MGF, but computing it can
restore the original without-replacement problem.

\clearpage
\section{Entry O: positive Gaussian fields and controlled reductions}
\label{sec:positive-gaussian-field}
\subsection{Use fixed cardinality and unit energy to change the sign of a quadratic}
The continuous Fourier route uses complex phases. There is also a positive integral, obtained by using
two exact facts together: there are exactly $K$ slots and every waveform has norm one.
For an ordered tuple $W$, stack its waveforms into $x_W\in\mathbb R^{nK}$ and write $s_W=\sum_k\phi(w_k)$.
Define
\[
 C=\frac{(KI_K-\one\one^T)\otimes I_n}{\sigma^2}\succeq0,\qquad
 \Gamma=(\Gamma_1,\ldots,\Gamma_K)\sim\mathcal N(0,C).
\]
The covariance is singular: $\sum_k\Gamma_k=0$. One can draw it as
$\Gamma_k=\sqrt K(Z_k-\overline Z)/\sigma$ with independent $Z_k\sim\mathcal N(0,I_n)$.
The Gaussian exponential identity gives
\[
 \E e^{\Gamma^Tx_W}
 =e^{x_W^TCx_W/2}
 =e^{(K^2-\|s_W\|^2)/(2\sigma^2)},
\]
because $\|x_W\|^2=K$. Therefore, putting $t=y/\sigma^2$ and
$c_y=e^{-(\|y\|^2+K^2)/(2\sigma^2)}$,
\begin{equation}
 g(y-s_W)=c_y\,\E_\Gamma\prod_{k=1}^K
       e^{(t+\Gamma_k)^T\phi(w_k)}.
 \label{eq:positive-field-identity}
\end{equation}
This is a direct Gaussian decoupling identity, not a Gaussian approximation to interference.
All cross-user dependence has moved into a shared field integral.

With $\psi(u)=M^{-1}\sum_w e^{u^T\phi(w)}$, the IID evidence becomes
\begin{equation}
 Z=c_y\,\E_\Gamma h(\Gamma),\qquad
 h(\Gamma)=\prod_k\psi(t+\Gamma_k).
 \label{eq:positive-field-evidence}
\end{equation}
Conditional on the field under the joint posterior, the $K$ messages are independent, with
\[
 p_{k,w}(\Gamma)=\frac{e^{(t+\Gamma_k)^T\phi(w)}}{M\psi(t+\Gamma_k)}.
\]
The posterior field law is $h(\Gamma)\,d\gamma(\Gamma)/\E_\gamma h$, where $\gamma=\mathcal N(0,C)$.
It is \emph{not} the unweighted Gaussian law. Mixing the conditional products with the wrong field law
would change the posterior.

\subsection{Count markers remain exact through the integral}
For a full word $w$, conditional count probabilities are coefficients of
\begin{equation}
 Q_w(z\mid\Gamma)=\prod_{k=1}^K[1-p_{k,w}(\Gamma)+z\,p_{k,w}(\Gamma)].
 \label{eq:field-count-polynomial}
\end{equation}
This is a sum of independent, generally nonidentical Bernoulli indicators. Averaging its coefficients under
the posterior field law gives the exact $q_{w,j}$, hence the exact tails. For a prefix cell $D$, replace
$p_{k,w}$ by $\sum_{w\in D}p_{k,w}$. Total cardinality is exactly $K$ throughout, and repeated messages are
retained. A count polynomial costs $O(K^2)$ arithmetic per queried word/cell after its conditional masses.
One field evaluation still requires $K$ polar partition queries and their pattern sums. Direct all-word
readout retains $M$ dependence; prefix discovery can reduce queries only if its bounds are useful.

Under the distinct source the identity for a fixed tuple still holds, but the conditional product must be
restricted to distinct words and renormalized with that restriction. Its partition is a sum over injective
assignments, not the product in \eqref{eq:positive-field-evidence}. The IID factorization must not be reused
for that control. At $K=1$, $C=0$ and this representation reduces to the original one-user partition sum.

\newpage
\subsection{Positivity removes cancellation, not importance-weight variance}
Under the base Gaussian, $h$ is generally unbounded, so H's bounded-weight theorem cannot simply be applied.
There is nevertheless an exact second-moment identity. Expanding two copies of $h$ and integrating the
Gaussian gives
\begin{equation}
 \frac{\E_\gamma h^2}{(\E_\gamma h)^2}
 =\E_{W,W'\ {\rm IID}\ P_{\rm slot}(\cdot\mid y)}
       e^{x_W^TCx_{W'}}
 \leq e^{K^2/\sigma^2}.
 \label{eq:field-second-moment}
\end{equation}
For the equality, the separate factors
$e^{t^Ts_W+x_W^TCx_W/2}/M^K$ normalize to the exact slot posterior.
For the inequality use positive-semidefinite Cauchy--Schwarz and
$x_W^TCx_W=(K^2-\|s_W\|^2)/\sigma^2\leq K^2/\sigma^2$.
The bound is finite but can be enormous. It also bounds
$\E_\gamma[(hH)^2]/(\E_\gamma h)^2$ for any conditional count-tail observable $H\in[0,1]$.

This is a real obstruction, not just a loose worst-case warning. Take $K=2$, scalar waveforms $\{-1,1\}$
with uniform prior, and $y=0$. Then $\Gamma=(g,-g)$ with $g\sim\mathcal N(0,\sigma^{-2})$ and
$h=\cosh^2g$. Direct Gaussian moments give
\[
 \E h=\frac{1+e^{2/\sigma^2}}2,\qquad
 \E h^2=\frac{3+4e^{2/\sigma^2}+e^{8/\sigma^2}}8,\qquad
 \frac{\E h^2}{(\E h)^2}\sim\tfrac12 e^{4/\sigma^2}\quad(\sigma\downarrow0).
\]
Thus plain independent Gaussian-field sampling can be poor even for two users. A useful field proposal,
deterministic integration or conditional summation needs its own analysis. Replacing a complex integral
by a positive one has exchanged one obstacle for another.

\subsection{Truncate the field domain with a posterior-mass bound}
Choose $B_C\in\mathbb R^{nK\times d}$ with $B_CB_C^T=C$, $d=n(K-1)$, and write
$\Gamma=B_CU$, $U\sim\mathcal N(0,I_d)$. Put $a_W=B_C^Tx_W$.
Completing the square now shows that the exact posterior of $U$ is the mixture
\begin{equation}
 \mathcal L(U\mid y)=\sum_W P_{\rm slot}(W\mid y)\,
          \mathcal N(a_W,I_d),\qquad \|a_W\|\leq K/\sigma.
 \label{eq:field-gaussian-mixture}
\end{equation}
The mixture weights are unknown, but the uniform bound on their centers is known.
For $R>K/\sigma$ and $Z_d\sim\mathcal N(0,I_d)$ this proves
\begin{equation}
 \Pr(\|U\|>R\mid y)\leq
       \Pr(\|Z_d\|>R-K/\sigma)=:\eta_R.
 \label{eq:field-domain-tail}
\end{equation}
The right side is an ordinary Gaussian radial tail, with no small evidence denominator.
Conditioning the joint law on $\|U\|\leq R$ changes its message-count marginal by TV at most $\eta_R$.
Further approximation inside the ball combines by A's restriction-plus-conditional-error rule.
This is an analytic posterior-mass guarantee, not a claim that the ball has low dimension.
On the ball, $h\leq\exp(K\|y\|/\sigma^2+KR/\sigma)$, since
$\sum_k\|\Gamma_k\|\leq\sqrt K\|\Gamma\|\leq KR/\sigma$.
Bounded-domain sampling can therefore use a bounded-weight argument, but may still have very poor
normalization and concentration. At $K=1$ the field is absent and no truncation is needed.

\newpage
\subsection{Discard field directions with an explicit positive correction}
Choose an orthogonal projection $\Pi$ in the whitened field coordinates and retain only $\Pi U$.
Equivalently set $C_{\rm keep}=B_C\Pi B_C^T$ and
$C_{\rm omit}=C-C_{\rm keep}\succeq0$. The retained-field integral has tuple weight
\begin{equation}
 \widetilde\omega_{\rm slot}(W)
 =\omega_{\rm slot}(W)e^{-D(W)/2},\qquad
 D(W)=x_W^TC_{\rm omit}x_W\geq0.
 \label{eq:field-direction-drop}
\end{equation}
This follows by comparing the two Gaussian exponential moments. Keeping the conditional one-user
partition functions therefore yields a positive, fixed-$K$ approximation, even with far fewer field
coordinates. If $D(W)\in[D_-,D_+]$ on all legal tuples, the normalized count laws satisfy
\begin{equation}
 \TV(P_y,\widetilde P_y)\leq
       \tanh\!\left(\frac{D_+-D_-}{8}\right).
 \label{eq:field-direction-tv}
\end{equation}
Marginalizing slots to counts cannot increase TV. The correction is $e^{D/2}$, so K's
observation-weighted or good-region bounds may improve this worst-case range when their quantities can
actually be bounded. Small average discarded field energy alone is not a certificate.

Integrating the omitted directions \emph{exactly} would instead restore $e^{D/2}$ and generally recouple
the messages. That is why reducing the field dimension is an approximation. Choosing $\Pi$ requires
encoder/conditional geometry, not only the spectrum of $C$: all its nonzero eigenvalues are $K/\sigma^2$.
Report 7's exact full-bank observation span offers no global null directions here. More restricted
alphabets or decision-relevant regions can be different.

A common recentering $\phi(w)\mapsto\phi(w)-m$ does not by itself help this field reduction:
$C(\one\otimes m)=0$, and $B_C^T(\one\otimes m)=0$, so $a_W$ and $D(W)$ are unchanged.
This is a useful limit on combining reference changes with this particular representation.
A chosen projection or finite field quadrature can also break equal-signal conditional identities;
known alias groups must retain their source-weighted marks, or use the exact conditional restoration in K.
Positivity alone does not enforce those identities.

Field integration, its proposal density and the inner polar sums remain separate costs. Evaluating a finite
positive quadrature gives a mixture of normalized product laws only when its evidence weights and
conditional probabilities are normalized consistently. Independently approximated inner denominators can
destroy that interpretation. Exact sampling from the retained-field posterior is not supplied merely by
writing its integral, and correction does not make its unknown marginal tuple proposal free to evaluate.
The useful new assumption is small discarded quadratic variation, or an affordable accurate field
integration; neither has been established for the selected bank.

\clearpage
\section{Entry P: extract more information from the moment bounds already paid for}
\label{sec:moment-consistency}
\subsection{The count is an integer in a small known interval}
M's alternating formulas give valid bounds, but need not use all information in the retained moments.
For any word or prefix count $A\in\{0,\ldots,K\}$, suppose
\[
 \ell_r\leq\E_{P_y}\binom Ar\leq u_r,\qquad r=1,\ldots,m.
\]
For a required function $f(c)$, define the finite optimization problems
\begin{equation}
 \begin{split}
 \underline f,\overline f={}&\min,\max\ \sum_{c=0}^K f(c)p_c,\\
 &p_c\geq0,\quad \sum_cp_c=1,\quad
 \ell_r\leq\sum_c\binom crp_c\leq u_r\quad(1\leq r\leq m).
 \end{split}
 \label{eq:count-moment-lp}
\end{equation}
The true count law is feasible, so these bound its expectation. They are the sharp bounds given
\emph{these constraints}, not necessarily sharp bounds for the full physical posterior.
Take $f(c)=\ind\{c\geq j\}$ for a tail, or $f(c)=\min(c,h)$ for I's prefix gain.
This is a finite version of the classical optimization approach to moment inequalities
\cite{bertsimas-popescu}; here the known integer support makes it an ordinary linear program.

There are $K+1$ variables and $O(m)$ constraints, with $O(Km)$ coefficient data.
Solving it has a polynomial finite-LP cost at a stated numerical precision; that is additional work,
not an $O(Km)$ solve claim. Full order with exact moments identifies the count law because the
binomial-coefficient matrix is triangular. Lower order can remain ambiguous.
If valid intervals have empty intersection, an input bound, arithmetic error or source contract is wrong;
silently projecting to a convenient moment vector would invalidate the certificate.

For an algebraic example, let $K=3$, $\E A=1$ and $\E\binom A2=1/5$.
Writing $p_3=t$ gives
\[
 p_2=1/5-3t,\quad p_1=3/5+3t,\quad p_0=1/5-t,\quad 0\leq t\leq1/15.
\]
Hence $P(A\geq1)\in[4/5,13/15]$, with both endpoints attained by legal count laws.
The first two alternating bounds alone give $[4/5,1]$.
The stronger upper bound uses the bounded integer support, without another residual likelihood query.

\subsection{Dual polynomials explain both the bound and the next useful query}
For exact moments $\mu_r$ with $\mu_0=1$, any polynomial in the binomial basis
$p(c)=\sum_{r=0}^m\lambda_r\binom cr$ satisfying $p(c)\geq f(c)$ for all integers $0\leq c\leq K$ gives
\[
 \E f(A)\leq\lambda_0+\sum_{r=1}^m\lambda_r\mu_r.
\]
The best such upper polynomial is the finite LP dual. With moment intervals, the same feasible polynomial
gives the certified bound
\begin{equation}
 \E f(A)\leq\lambda_0+
       \sum_{\lambda_r\geq0}\lambda_r u_r+
       \sum_{\lambda_r<0}\lambda_r\ell_r.
 \label{eq:moment-dual-interval}
\end{equation}
A polynomial below $f$ gives the lower counterpart. Only the $K+1$ integer values must be checked;
nonnegativity between them is unnecessary. M's alternating and single-moment bounds are particular
feasible polynomials, so optimization over the same information cannot weaken their intersection.

The coefficients show which moment uncertainties contribute to this certificate: reducing an interval
width by $\Delta_r$ can change this fixed dual bound by at most $|\lambda_r|\Delta_r$.
This supplies an algebraic reason to refine a particular query rather than raise every moment order.
It predicts sensitivity of the certificate, not an assured improvement in decoding performance.
Floating-point optimization must be followed by valid dual feasibility and rounding bounds if a
machine-certified interval is claimed.

\subsection{Keep the shared evidence instead of dividing every interval separately}
Sometimes F/G/N provide unnormalized moment bounds
$L_r\leq F_r=\E_\pi[\binom Ar g(y-\Phi a)]\leq U_r$ and
$0<Z_-\leq Z\leq Z_+$. Introduce $\tau=1/Z$ and the normalized count probabilities $p_c$.
Then impose the linear constraints
\begin{equation}
 L_r\tau\leq\sum_c\binom crp_c\leq U_r\tau,\qquad
 Z_-\tau\leq1\leq Z_+\tau,\quad p_c\geq0,\quad\sum_cp_c=1.
 \label{eq:moment-shared-evidence}
\end{equation}
All queried words/cells can share the same $\tau$. This keeps the common denominator coupled to their
moments, instead of treating separate ratio endpoints as if their extremes could occur independently.
N's positive Jensen bound is one possible provider of $Z_-$. Bounds from different representations may
be intersected when each bounds the same exact numerator/evidence with the correct source normalization.
Shared approximate formulas without proven bounds are not additional exact constraints.

\subsection{Enforce consistency between prefix counts}
Let a finite binary prefix tree partition the full alphabet, and write $A_D$ for a cell's total count.
The root count is $K$. At every split $D=D_0\mathbin{\dot\cup}D_1$, introduce a nonnegative table
$q_D(c,d)$ for its two child counts, with $c,d\geq0$, $c+d\leq K$. Require
\begin{equation}
 \begin{split}
 \sum_d q_D(c,d)&=p_{D_0}(c),\qquad
 \sum_c q_D(c,d)=p_{D_1}(d),\\
 \sum_{c+d=a}q_D(c,d)&=p_D(a),\qquad p_{\rm root}(a)=\ind\{a=K\}.
 \end{split}
 \label{eq:prefix-count-consistency}
\end{equation}
Add available moment or tail intervals to the corresponding marginals.
There are $O(N_{\rm tree}K^2)$ table entries and consistency constraints for a constructed tree,
not one variable per message multiset. Under distinct sampling impose the extra capacities
$A_D\leq\min(K,|D|)$; singleton counts are then $0/1$.

These local consistency equations really do admit a global legal distribution over leaf counts.
To prove it, start at the deterministic root count and, at each node with count $a$, draw child counts
with probability $q_D(c,d)/p_D(a)$ on $c+d=a$. Zero-probability parent counts need no draw.
Continuing down the tree gives every prescribed node/split marginal and total count $K$.
The constructed law is a feasibility witness, not an assertion that the physical posterior has this
conditional independence. Likelihood information still comes from the supplied intervals; with few such
constraints the outer set of possible posteriors can be very large.

There is an exact source for selected mixed-moment constraints. For disjoint cells $D_0,D_1$ with
prior masses $\rho_0,\rho_1$, IID sampling gives, for $r+s\leq K$,
\begin{equation}
 \E_\pi\!\left[\binom{A_{D_0}}r\binom{A_{D_1}}s g(y-\Phi a)\right]
 =\frac{K!}{r!s!(K-r-s)!}\rho_0^r\rho_1^s
 \E g\!\left(y-\sum_{i=1}^rX_{0,i}-\sum_{i=1}^sX_{1,i}
                   -\sum_{i=1}^{K-r-s}X_i\right).
 \label{eq:mixed-prefix-factorial}
\end{equation}
Here the first two groups are independent draws conditional on their cells; the remaining draws use
the full prior, including those cells. Selecting the disjoint sets of marked slots proves the coefficient.
These are mixed factorial moments, not probabilities of exactly $r,s$ occurrences.
F or N can bound this nonidentical-increment residual using I's conditional moment information, and the
result constrains $\sum_{c,d}\binom cr\binom ds q_D(c,d)$.
One can purchase mixed moments only at ambiguous splits, with their construction cost stated.
The IID formula is not a without-replacement formula.

For fixed competing allocations supported on resolved full-word leaves, their expected gain contrast is
linear in the leaf count probabilities. Optimizing it over this consistency set can certify a comparison
even when individual tails remain broad. This does not make the simultaneous optimization over every
possible allocation and posterior a single LP. I's optimistic prefix-gain upper bound remains available,
now with tighter certified query intervals. Unresolved prefixes still cannot be output as recovered words.

\clearpage
\section{Entry Q: make the inner polar queries finite and certifiable}
\label{sec:polar-inner-closure}
\subsection{A conditional character algebra computes more than uniform moments}
N and O require nonuniform polar sums, not just F's uniform first two moments.
Fix a pattern and a payload prefix $v=v_0+Rz$, with $z$ uniform in $\mathrm{GF}(2)^b$.
Group equal restricted character rows to write a linear score as
\[
 s(z)=\theta^T\phi(p,v_0+Rz)
     =a_0+\sum_{g\in\mathcal G_*}a_g\chi_g(z),\qquad
 \chi_g(z)=(-1)^{g^Tz},\quad g\ne0.
\]
The coefficients include the known prefix signs and physical amplitudes. They obey
$\chi_g\chi_h=\chi_{g+h}$ and $\E\chi_g=\ind\{g=0\}$.
Consequently every polynomial in $s$ can be averaged by multiplying character polynomials, adding their
binary keys and reading the zero coefficient. Inserting $\chi_h$ reads coefficient $h$ instead.
The same calculation therefore gives weighted symbol means and second moments using the original
character rows and their pairwise sums. It retains all polar/CRC constraints.

This is finite arithmetic without enumerating every payload at a fixed polynomial degree, but not free
arithmetic. If $D=|\mathcal G_*|$, expansion through degree $d$ has at most
$\binom{D+d}{d}$ nonconstant-character monomials before grouping; a dense character table has $2^b$
keys. Either construction can be used, with polynomial coefficient/key work and the moment readout
charged. A sparse dictionary can reach the dense table, and high degree can cost more than direct payload
enumeration. Sum over every allowed pattern with its original conditional prior mass.

\subsection{A complementary subset expansion can exploit independent characters}
There is an exact alternative to expanding powers of the score:
\[
 e^{s(z)}=e^{a_0}\prod_{g\in\mathcal G_*}\cosh(a_g)
                         \prod_{g\in\mathcal G_*}[1+\tau_g\chi_g(z)],
 \qquad \tau_g=\tanh(a_g).
\]
For an inserted character $h$, its expectation is the common prefactor times
\[
 S_h=\sum_{\substack{A\subseteq\mathcal G_*\\\sum_{g\in A}g=h}}\prod_{g\in A}\tau_g.
\]
The partition uses $h=0$. Retaining subsets of size at most $d$ has absolute remainder bounded by
$\sum_{j>d}e_j(|\tau|)$, where
$\prod_g(1+z|\tau_g|)=\sum_j e_j(|\tau|)z^j$.
These elementary-symmetric coefficients cost $O(D^2)$ scalar work, while the retained subset terms cost
up to $\sum_{j\leq d}\binom Dj$ key/weight operations. Multiply by the common prefactor for the
unnormalized remainder. Negative $\tau_g$ can cause cancellation, so use numerator/denominator intervals.
The truncated sums do not automatically describe a positive tilted law.

If the keys are linearly independent, only the empty subset sums to zero, giving
$\psi=e^{a_0}\prod_g\cosh(a_g)$ exactly. Each attainable $h$ has one subset representation.
This can help after conditioning a prefix. Independence requires at most $b$ nonzero character keys,
so it does not provide a general shortcut for the full bank.
This expansion and the positive polynomial below exploit different regimes; neither is uniformly better.

\subsection{A positive polynomial gives one coherent approximate tilted source}
Choose certified score bounds $a\leq s(z)\leq b_+$; for example
$a=a_0-\sum_g|a_g|$ and $b_+=a_0+\sum_g|a_g|$. Put $L_s=b_+-a$ and
\[
 T_d(t)=\sum_{j=0}^d\frac{t^j}{j!},\qquad
 q_d(s)=e^aT_d(s-a),\qquad
 \rho_d(L_s)=e^{-L_s}T_d(L_s)>0.
\]
For $0\leq t\leq L_s$, $e^{-t}T_d(t)$ decreases, because its derivative is
$-e^{-t}t^d/d!$. Thus
\begin{equation}
 \rho_d(L_s)e^s\leq q_d(s)\leq e^s,\qquad
 \psi_d:=\E q_d(s)\leq\psi:=\E e^s\leq\psi_d/\rho_d(L_s).
 \label{eq:polar-positive-series}
\end{equation}
The polynomial is strictly positive on every legal word. Normalize its weights to obtain a genuine
tilted-source approximation $Q_d$. Since $q_d\leq e^s$, the exact tilted law is a mixture with component
weight $\psi_d/\psi$ equal to $Q_d$. Therefore
\[
 \TV(P_\theta,Q_d)\leq1-\psi_d/\psi\leq1-\rho_d(L_s).
\]
A normalized-weight range also gives the alternative bound
$\tanh[-\log\rho_d(L_s)/4]$; take whichever certified bound is smaller.
The series converges as $d$ increases for every finite score range. Its useful degree can grow with
tilt magnitude and inverse noise; fixed low degree has no uniform accuracy guarantee.

Calculate $Q_d$'s moments by inserting the waveform or its products in the \emph{same polynomial sum}.
Differentiating the truncated partition does not generally return these moments, because
$T_d'=T_{d-1}$, not $T_d$. This distinction prevents an inconsistent mean/covariance closure.
Third absolute moments can use F's unit-norm trace bounds under this actual positive law.

With different polynomial bounds on different patterns, sum their lower partitions and upper partitions
separately. If the totals are $\psi_-\leq\psi\leq\psi_+$ with $Q_d$ normalized by $\psi_-$, then
$\TV(P_\theta,Q_d)\leq1-\psi_-/\psi_+$. This can improve the worst-pattern bound.
Word exclusions use the allowed-source sum with its correct prior normalization; algebraic subtraction
of a forbidden term requires controlled rounding if cancellation occurs.
When computing child masses of a fixed $Q_d$, retain its same polynomial weights.
Choosing a fresh polynomial independently in each child instead defines a different sequential proposal,
which can still be used through L with its actual branch density.

\subsection{Propagate the inner approximation through the outer residual closure}
For N's independent tilted sources, suppose source $i$ has
$\psi_{i,-}\leq\psi_i\leq\psi_{i,+}$ and a coherent approximation $Q_i$ with TV error at most $\eta_i$.
A product coupling bounds the difference of the exact and approximate sum laws by
\[
 \eta_{\rm in}=1-\prod_i(1-\eta_i).
\]
Because $0\leq g\leq1$, the same quantity bounds their smoothed likelihood difference.
Use the exact means/covariances of the $Q_i$ in F's Gaussian formula at $r-\sigma^2\theta$, obtaining
$\widetilde J_Q$ and replacement bound $\delta_Q$. Define
\[
 A_\pm=e^{-\theta^Tr+\sigma^2\|\theta\|^2/2}\prod_i\psi_{i,\pm}.
\]
Then a complete nested bound is
\begin{equation}
 A_-\max(0,\widetilde J_Q-\delta_Q-\eta_{\rm in})
 \leq L(r)\leq
 A_+\min(1,\widetilde J_Q+\delta_Q+\eta_{\rm in}).
 \label{eq:nested-tilted-bound}
\end{equation}
Intersect with $[0,1]$. A positive lower bound is also available from Jensen under the \emph{original}
source, using its known uniform moments; do not apply N's tilted Jensen formula to guessed exact moments.
This closes the inner/outer error interface. The parameters are now series degree, polynomial/key cost,
inner TV, outer smoothing error and the resulting likelihood interval, rather than an unspecified tilted oracle.

\subsection{A convergent residual hierarchy need not assume Gaussian interference}
The same positive series also closes a limit of F: making a Gaussian closure more numerically accurate
does not remove its non-Gaussian modeling error. Instead approximate the likelihood function itself.
For $S$ the sum of $h$ unit waveforms, put
\[
 a_r=-\frac{(\|r\|+h)^2}{2\sigma^2},\qquad
 \ell_r(S)=-\frac{\|r-S\|^2}{2\sigma^2},\qquad
 g_d(r-S)=e^{a_r}T_d(\ell_r(S)-a_r).
\]
Since $a_r\leq\ell_r(S)\leq0$, \eqref{eq:polar-positive-series} gives
\[
 \rho_d(-a_r)g(r-S)\leq g_d(r-S)\leq g(r-S).
\]
Now $\ell_r(S)-a_r$ is quadratic in $S$, so $\E g_d(r-S)$ requires moments only through degree $2d$.
Independent-sum moments follow by multiplying truncated multivariate moment-generating polynomials;
each waveform moment is a character average with the correct prefix/excluded source.
There are $\binom{n+2d}{2d}$ possible monomials through that degree, before any symmetry or exact projection.
Constructing the one-user tensors, their source-specific versions and their polynomial convolutions is
charged. Dimension and degree can make this much worse than enumeration.

For a single coherent full-posterior approximation, use the same
$a=-(\|y\|+K)^2/(2\sigma^2)$ and degree for all configurations.
Its positive weights have log-ratio range at most $-\log\rho_d(-a)$, and the original count-tail error
is bounded through Section 4. Increasing $d$ drives this bound to zero without a Gaussian assumption.
Separate residual queries can use tighter local $a_r$ as interval providers, but need not be marginals
of that one polynomial-weight posterior. This is a finite-accuracy hierarchy, not a useful low-degree
claim at low noise. It gives a mathematically controlled alternative when F's Gaussian assumption fails.

\subsection{The same character calculation supplies continuous Fourier queries}
G's continuous Fourier route also has a missing one-user sum:
$\chi(\Omega)=\E e^{-i\Omega^T\phi(w)}$.
For the real score $s=\Omega^T\phi(w)$, Taylor's integral remainder gives
\[
 \left|e^{-is}-\sum_{j=0}^d\frac{(-is)^j}{j!}\right|
 \leq\frac{|s|^{d+1}}{(d+1)!}.
\]
Since $\|\phi(w)\|=1$, $|s|\leq\|\Omega\|$, or use a sharper conditional bound.
Character arithmetic computes the polynomial expectation. Projecting the resulting complex number onto
the closed unit disk cannot increase its distance from the true characteristic function.
For disk-valued $\chi,\widetilde\chi$,
$|\chi^h-\widetilde\chi^h|\leq h|\chi-\widetilde\chi|$ by telescoping the power.

For a cell $D$ of prior mass $\rho$, an exact-count frequency term has source factor
$\binom Kj\rho^j(1-\rho)^{K-j}\chi_D^j\chi_{D^c}^{K-j}$.
If both conditional characteristics have uniform error at most $\epsilon_\chi$, its error is at most
$\binom Kj\rho^j(1-\rho)^{K-j}K\epsilon_\chi$. The evidence factor has error at most $K\epsilon_\chi$.
Add these deterministic errors to G's independent-frequency confidence bounds before normalization.
Choosing the degree from the observed $\|\Omega\|$ can enforce the same tolerance at every frequency;
the degree and runtime are then random and potentially large. Frequency sampling remains independent.
This avoids a finite-grid coefficient table and can avoid payload enumeration at modest degree, while
retaining pattern access, cancellation and the small-evidence obstacle. It supplies queries, not a
positive approximate posterior. Empty or full cells are handled deterministically.

\clearpage
\section{Entry R: close the field geometry and integration questions}
\label{sec:field-closure}
\subsection{Use a physical projection that respects exchangeable slots}
O allowed a general field projection. A structured choice retains a physical subspace $V$ identically
for every slot, with orthogonal projection $P_V$:
\[
 C_{\rm keep}=\frac{(KI-\one\one^T)\otimes P_V}{\sigma^2},\qquad
 C_{\rm omit}=\frac{(KI-\one\one^T)\otimes(I-P_V)}{\sigma^2}.
\]
Its field dimension is $(K-1)\dim V$. Writing $z_w=(I-P_V)\phi(w)$, the discarded correction is exactly
\begin{equation}
 D(W)=\frac{K\sum_k\|z_{w_k}\|^2-\|\sum_kz_{w_k}\|^2}{\sigma^2}
     =\frac1{\sigma^2}\sum_{i<j}\|z_{w_i}-z_{w_j}\|^2.
 \label{eq:field-dispersion}
\end{equation}
It is slot-symmetric and hence a function of counts. If a certified omitted-space diameter is $\Delta_V$,
then $0\leq D\leq\binom K2\Delta_V^2/\sigma^2=:D_{\max}$, usable in O's TV bound.
The radius bound $\max_w\|(I-P_V)(\phi(w)-m)\|\leq r_V$ also gives
$D\leq K^2r_V^2/\sigma^2$; use the smaller available bound.
The center changes this radius certificate, not the exact value of $D$.

With prior covariance $\Sigma$,
\[
 \E_\pi D=\frac{K(K-1)}{\sigma^2}
                  \operatorname{tr}((I-P_V)\Sigma).
\]
The top covariance eigenvectors minimize this prior average among subspaces of a fixed dimension.
F's polar moments construct $\Sigma$ without payload enumeration; eigendecomposition costs $O(n^3)$.
This gives a concrete candidate subspace, not a posterior certificate. For example,
$\E_{P_y}D\leq\E_\pi D/Z_-$ uses a valid evidence lower bound but can be vacuous.
Worst-case radii can be bounded without enumerating payloads: if
$\phi(p,v)=C_{p,0}+\sum_{g\ne0}C_{p,g}\chi_g(v)$, then
\[
 r_V\leq\max_p\left[
 \|(I-P_V)(C_{p,0}-m)\|+
                  \sum_{g\ne0}\|(I-P_V)C_{p,g}\|\right].
\]
All patterns and projected character vectors are charged. These explicit bounds can be loose; that
limits the certificate, rather than authorizing a small-error claim from the covariance spectrum alone.

\subsection{Weighted dispersion can be queried through one-user moments}
Let $Q$ be the normalized retained-field law on tuples, and let $R=e^{D/2}$, so $P=QR/\E_QR$.
If $0\leq D\leq D_{\max}$ and $d_Q\geq\E_QD$, convexity on that interval gives
\begin{equation}
 x:=\E_Q(R-1)\leq
 a_D:=\frac{e^{D_{\max}/2}-1}{D_{\max}}d_Q,\qquad
 \TV(P,Q)\leq\frac{x}{1+x}\leq\frac{a_D}{1+a_D}.
 \label{eq:field-weighted-dispersion}
\end{equation}
Use the continuous coefficient limit $1/2$ at $D_{\max}=0$.
For the TV inequality, $R\geq1$ and
$\E|R-\E R|\leq2\E(R-1)$.
The exponential coefficient exposes why an average dispersion alone may be insufficient.

Conditional on a retained field, the slot laws are independent. Let their waveform means and covariances
be $\mu_k,\Sigma_k$. The conditional expectation of \eqref{eq:field-dispersion} is
\[
 \frac1{\sigma^2}\sum_{i<j}
 \left[\operatorname{tr}((I-P_V)(\Sigma_i+\Sigma_j))
              +\|(I-P_V)(\mu_i-\mu_j)\|^2\right].
\]
It needs one-user moments and $O(Kn^2+Kn)$ dense aggregate arithmetic, using the sum-of-squares identity,
rather than message-pair enumeration. Integrate or bound it under the evidence-weighted field law to
obtain $d_Q$; sampling requires its own confidence allowance. Q's inner moment errors must be enclosed.
Thus the weighted correction has a specified query, while outer integration remains charged.

\subsection{A bounded Gaussian proposal in a stated curvature regime}
For exact inner partitions define $\mu(u)=\nabla\log\psi(u)$ and
$\Sigma(u)=\nabla^2\log\psi(u)$. Suppose a known matrix $\overline\Sigma$ bounds
$0\preceq\Sigma(u)\preceq\overline\Sigma$ for every field argument being used.
Unit waveforms always permit $\overline\Sigma=I$; smaller global bounds require proof.
Write the retained field as $\Gamma=B U$, $BB^T=C_{\rm keep}$ and $U\sim\mathcal N(0,I_d)$, and put
\[
 A=B^T(I_K\otimes\overline\Sigma)B,\qquad h_0=\psi(y/\sigma^2)^K.
\]
Taylor's integral formula for $\log\psi$ gives
\[
 0\leq\log\psi(t+\Gamma_k)-\log\psi(t)-\mu(t)^T\Gamma_k
       \leq\tfrac12\Gamma_k^T\overline\Sigma\Gamma_k.
\]
The linear terms cancel because $\sum_k\Gamma_k=0$. Therefore
\begin{equation}
 h_0\leq h(BU)\leq h_0e^{U^TAU/2}.
 \label{eq:field-quadratic-envelope}
\end{equation}
If $A\prec I$, let $D_A=\det(I-A)^{-1/2}$ and use the explicit proposal
$q_A=\mathcal N(0,(I-A)^{-1})$. Its weight for the field integral is
\begin{equation}
 W(U)=\frac{\varphi_d(U)h(BU)}{q_A(U)}
     =D_Ah(BU)e^{-U^TAU/2}\leq D_Ah_0,\qquad
 h_0\leq\E_{\varphi_d}h\leq D_Ah_0.
 \label{eq:field-bounded-proposal}
\end{equation}
Here $\varphi_d$ is the standard Gaussian density. Thus H has a certified concentration parameter
$\gamma\leq D_A$, without a small physical-evidence denominator. Conditional count polynomials are bounded
observables in $[0,1]$, so the same bounded-variable proof applies to their weighted expectations.
They need not be sampled as categorical messages. Full-word discovery/readout is still necessary.

For $\overline\Sigma=cI$ and a rank-$r$ physical projection,
$A=(Kc/\sigma^2)I_{r(K-1)}$ in suitable whitened coordinates. The sufficient regime and cost factor are
\[
 Kc<\sigma^2,\qquad D_A=(1-Kc/\sigma^2)^{-r(K-1)/2}.
\]
Even inside that regime, the determinant can be large. This is a finite conditional guarantee, not an
assertion that the selected bank satisfies a useful regime. Constructing a general $A$ and factoring
$I-A$ can cost $O(d^3)$; the displayed isotropic/Kronecker case avoids a generic factorization.
Every field draw also needs its inner polar partitions and marked queries.

\subsection{Outside that regime the obstruction is visible in the algebra}
The exact log field density has Hessian
$-I+B^T\diag(\Sigma(t+\Gamma_1),\ldots,\Sigma(t+\Gamma_K))B$.
For the full field, zero is stationary by slot symmetry. Its curvature eigenvalues are
$-1+K\lambda_j(\Sigma(t))/\sigma^2$, each repeated $K-1$ times.
If one is positive, zero is not a local maximum: a single Gaussian expansion there cannot be justified
by silently clipping a negative precision eigenvalue. Failure of $A\prec I$ is only failure of this
global envelope; the exact field posterior remains well defined.

O's finite-radius restriction gives a separate universally finite route: its omitted posterior mass is
bounded, its integrand is bounded on the ball, and independent draws from the Gaussian conditioned on
the ball can use H. Rejection sampling has known acceptance probability
$\Pr(\|U\|\leq R)$; attempts must be charged. These alternatives close the integration definition and
error accounting without promising useful dimension or sample cost.

Inner approximations are another distinct error. If Q's positive word weights have ratio in $[\rho,1]$
to the exact exponentials uniformly over the retained field domain, their $K$-fold joint product has
ratio in $[\rho^K,1]$. The joint/count TV bound is at most $\tanh(K\log(1/\rho)/4)$.
Inside a bounded field ball, a finite degree can enforce this uniformly. On the full field space, degree
may instead grow with the field norm. A pointwise approximation at sampled fields alone does not prove
a uniform error bound. Add field restriction, direction, inner-query and numerical sampling errors
according to their actual conditional laws, not as interchangeable variance terms.

\clearpage
\section{Entry S: choose a reference by a declared variational objective}
\label{sec:variational-reference}
\subsection{The common center has a principled, but limited, optimum}
Report 7 allowed any common center $m$ in its exact pair/reference reparameterization.
One way to choose it is to optimize the IID product family $Q_\nu(W)=\prod_{k=1}^K\nu(w_k)$, with the
same distribution in each slot. This family retains a multinomial count law and total $K$; it does
approximate the interactions. Variational lower bounds are classical \cite{wainwright-jordan}.

For $m=\E_\nu\phi(w)$ and $u(w)=1/M$, the evidence lower bound is
\begin{align}
 \mathcal F(\nu)
 &:=\E_{Q_\nu}\log\frac{\omega_{\rm slot}(W)}{Q_\nu(W)}\\
 &=-\frac{\|y\|^2-2K y^Tm+K+K(K-1)\|m\|^2}{2\sigma^2}
            -K D_{\rm KL}(\nu\|u)\leq\log Z.
 \label{eq:variational-product}
\end{align}
The quadratic follows by separating the $K$ unit self-energies and the $K(K-1)$ independent cross terms.
The bound follows from Jensen applied to $\log\E_Q(\omega/Q)$.
It is strictly concave in $\nu$: the negative quadratic is concave and negative relative entropy is
strictly concave. Its unique maximizer satisfies
\begin{equation}
 \nu_*(w)=\frac{e^{\theta_*^T\phi(w)}}{M\psi(\theta_*)},\qquad
 \sigma^2\theta_*+(K-1)\mu(\theta_*)=y.
 \label{eq:variational-reference-optimum}
\end{equation}
The tilt is the unique minimizer of
$(K-1)\log\psi(\theta)+\sigma^2\|\theta\|^2/2-y^T\theta$, whose Hessian is
$(K-1)\Sigma(\theta)+\sigma^2I\succ0$.
With $m_*=\mu(\theta_*)$, this is exactly Report 7's unary reference
$\nu^{(m_*)}(w)\propto e^{[y-(K-1)m_*]^T\phi(w)/\sigma^2}$.
It gives a declared criterion for centering before a connected expansion or choosing a product proposal.
It does not minimize a connected-tree remainder, forward KL, or Bayes decision loss.
N's full-$K$ residual tilt instead has coefficient $K$; these two optimizations must not be conflated.

This optimization also has a finite accuracy route. Since $\Sigma(\theta)\preceq I$,
the tilt objective has Hessian between $\sigma^2I$ and $(\sigma^2+K-1)I$.
Gradient descent with step $1/(\sigma^2+K-1)$ contracts distance to $\theta_*$ by at most
$1-\sigma^2/(\sigma^2+K-1)$ per exact-gradient step, as follows by integrating the Hessian along
the line to the optimum. The equation gives
$\|\theta_*\|\leq(\|y\|+K-1)/\sigma^2$, so starting at zero has an explicit initial bound.
Q provides finite inner sums and certified gradient errors; a uniform gradient error $\delta$
adds at most $\delta/\sigma^2$ to the limiting distance bound.
Every iteration's partition/moment cost is charged. Any finite trial tilt still defines a valid
reference, so exact convergence is unnecessary: use that reference's actual objective and remainder.

For the centered interaction $b^{(m)}(u,v)=e^{-(\phi(u)-m)^T(\phi(v)-m)/\sigma^2}$,
a certified radius $\|\phi(w)-m\|\leq r_m$ gives the positive floor
$b_0=e^{-r_m^2/\sigma^2}$. Then $b^{(m)}/b_0-1\geq0$ and E's auxiliary-component construction applies
with the centered unary reference under the IID source. The exact centered pair operator is still a diagonal rescaling of
Report 7's original operator; the conservative floor itself needs no pair-table enumeration.
Its improved radius may or may not compensate for the changed reference and feature magnitudes.
Freeze $m$, the reference and the floor for the target $K$ throughout the connected coefficient recurrence;
choosing a new center for each intermediate block size would change that identity.

Strict concavity concerns identical slot marginals. Allowing different $\nu_k$ gives a different,
generally nonconcave problem; no best-product claim for that larger family is made.
The distinct source also needs a different constrained family. The tilt equation above is not a
without-replacement variational formula.

\newpage
\subsection{An evidence sandwich can certify the variational approximation}
For any evaluable positive legal law $Q$,
\[
 D_{\rm KL}(Q\|P_{\rm slot})=\log Z-\mathcal F(Q).
\]
If a separate calculation gives $Z\leq Z_+$ and a certified lower bound
$\mathcal F(Q)\geq F_-$, then
\begin{equation}
 \TV(Q_{\rm counts},P_y)
 \leq\min\!\left\{1,\sqrt{\frac{\log Z_+-F_-}{2}}\right\}.
 \label{eq:variational-certificate}
\end{equation}
This is the finite probability inequality $\TV^2\leq D_{\rm KL}/2$, followed by marginalization.
Together with the decision bound it controls excess multiset risk. Under the identical-slot product,
both laws are uniform over orderings given counts, so their slot/count KL divergences are equal.
A negative computed gap signals inconsistent bounds or rounding, not exceptional accuracy.
For the inequality, coarse-grain onto the event attaining TV. Relative entropy can only decrease under
this grouping, and binary relative entropy as a function of its first probability has second derivative
$1/[q(1-q)]\geq4$ and value/first derivative zero at the second probability. Taylor's bound gives
$D_{\rm KL}\geq2\TV^2$.

For an exponential reference,
$D_{\rm KL}(\nu_\theta\|u)=\theta^T\mu(\theta)-\log\psi(\theta)$.
Thus its objective uses the same partition and moments as N/Q. Certified intervals must bound the entire
objective, including its quadratic, if these quantities are approximate.
The upper evidence can come from N, A/I, a positive contraction, or a signed computation with a valid
upper remainder. Merely increasing a lower bound gives no approximation guarantee.

For $K=2$, scalar $\phi=\pm1$ and $y=0$, the optimum in this symmetric family is uniform for every
$\sigma$. As $\sigma\downarrow0$, the exact posterior puts probability tending to one on one copy of
each word, whereas $Q_\nu$ assigns that count vector probability $1/2$.
The approximation therefore remains wrong despite exact optimization.
Indeed $\mathcal F=-1/\sigma^2$ and $\log Z\to-\log2$, so the KL gap diverges.
Optimizing a reference and certifying its adequacy are separate tasks. Using that reference inside an
exactly corrected sampler or the full connected expansion remains valid even when its standalone
posterior approximation is poor.

\clearpage
\section{Entry T: signed compression with explicit error bounds}
\label{sec:signed-compression}
\subsection{A low-rank query provider need not be a positive posterior}
D/J retained positive features so that their contractions define probabilities. Signed factors can
sometimes use a smaller rank; Report 7's exact rank obstruction does not forbid accurate low-rank
approximation. Classical tensor decompositions provide such representations \cite{oseledets}.
The issue here is to transport the approximation error through source constraints and count queries.

At a finite cut let $G$ be the exact likelihood table, $\widetilde G$ its rank-$r$ approximation and
$a,b$ the actual incoming vectors, including their known prior/likelihood weights. A certified
operator-norm residual $\|G-\widetilde G\|_2\leq\epsilon_r$ gives
\begin{equation}
 |a^T(G-\widetilde G)b|\leq\epsilon_r\|a\|_2\|b\|_2.
 \label{eq:signed-cut-error}
\end{equation}
With a full SVD, $\epsilon_r$ is the next singular value. The vectors must be those of the actual
contraction or have valid norm bounds; exact outside messages are not free oracles.
Preserve occupancy and any query counter in the cut indices. A small error for the unmarked evidence
alone does not bound every separately marked contraction.

This can reduce applying a known $s$-by-$t$ table to $O(r(s+t))$ work, but its construction and residual
certificate may still require the whole table. A numerical rank chosen from a few sampled entries
without an omitted-entry/norm guarantee is not a certified approximation.
Signed intermediate values require G's numerator/denominator rules; clipping negative answers does
not repair a missing remainder.

\subsection{An explicit positive envelope handles many approximate factors}
For a legal configuration $x$, write its weight as $\pi(x)\prod_{e=1}^E f_e(x)$ with $f_e\geq0$.
Keep all hard source, CRC and cardinality constraints exact. Suppose
$|f_e-\widetilde f_e|\leq\epsilon_e(x)$ and
$A_e(x)\geq|\widetilde f_e(x)|$, with nonnegative, constructible envelopes.
Expanding a telescoping product proves
\begin{equation}
 \left|\prod_e f_e-\prod_e\widetilde f_e\right|
 \leq\sum_{j=1}^E\epsilon_j
          \prod_{e<j}(A_e+\epsilon_e)\prod_{e>j}A_e
 =:\mathcal E(x).
 \label{eq:signed-factor-envelope}
\end{equation}
The proof replaces exact factors before the changed factor by
$f_e\leq|\widetilde f_e|+\epsilon_e\leq A_e+\epsilon_e$, and later approximate factors by $A_e$.
Every summand is nonnegative. Contract it with the original source and count markers to obtain
\[
 B_Z=\sum_x\pi(x)\mathcal E(x),\qquad
 B_H=\sum_x\pi(x)H(x)\mathcal E(x),\quad 0\leq H\leq1.
\]
Then $|Z-\widetilde Z|\leq B_Z$ and $|N_H-\widetilde N_H|\leq B_H$.
For a signed decision contrast use $|H|$ in its remainder.
Event/complement normalization or P's shared-evidence constraints complete the query interface.

If $\widetilde f_e=\sum_{\ell=1}^r u_\ell v_\ell$, the bound
$A_e=\sum_\ell|u_\ell||v_\ell|$ preserves a positive feature representation, though it may be much
larger than $|\widetilde f_e|$. A residual envelope can use
$\epsilon_e=\epsilon_r$ entrywise from an operator-norm bound, or a tighter structured bound.
One explicit finite method performs the $E$ positive contractions in
\eqref{eq:signed-factor-envelope}; their joint separator/feature costs remain charged.
The envelope can destroy an apparent speedup or be too loose, particularly when signed factors cancel.
This is the boundary of the route, not a reason to assume positivity or a small downstream error.

Zero residual recovers exact contraction. Replacing all factors by coarse low ranks is useful only when
both the main contraction and its marked error envelope are affordable. This completes a signed
alternative to D/J without requiring it to serve as a drawable proposal.

\clearpage
\section{Complete the guide and moment information arguments}
\label{sec:guide-moment-completion}
\subsection{A robust branch proposal follows from likelihood intervals}
At a legal binary prefix in L, let the child likelihoods have certified intervals
$L_b\in[\ell_b,u_b]$ and prior branch probabilities $\pi_b$.
Bayes' rule gives an interval $p_0\in[l,u]$ with
\[
 l=\frac{\pi_0\ell_0}{\pi_0\ell_0+\pi_1u_1},\qquad
 u=\frac{\pi_0u_0}{\pi_0u_0+\pi_1\ell_1}.
\]
Use $[0,1]$ when denominator information is absent, and handle truly impossible branches by the source
constraint. The minimax proposal over this interval is
\begin{equation}
 q_0=\frac{u}{1+u-l},\qquad
 q_1=\frac{1-l}{1+u-l},\qquad
 \sup_{p_0\in[l,u]}\max_b\frac{p_b}{q_b}=1+u-l.
 \label{eq:robust-bit-guide}
\end{equation}
To prove it, the worst two ratios are $u/q_0$ and $(1-l)/(1-q_0)$.
Their maximum is minimized when they are equal. Degenerate deterministic intervals select the legal
branch with probability one. With positive likelihoods on both legal children, valid bounds give
positive proposal probabilities.

This converts interval quality into a precise proposal-concentration resource: width $w=u-l$ costs a
factor $1+w$, at most two. If widths are uniformly at most $w_t$ over every reachable depth-$t$ prefix,
\[
 \sup_W\frac{P_{\rm slot}(W\mid y)}{q(W)}
 \leq\prod_{t=1}^{KB}(1+w_t).
\]
The proposed tuple density is the product of its recorded branch probabilities.
Bounds for visited prefixes alone do not establish the displayed global product.
Unexamined prefixes can always be assigned width one, giving a valid but potentially exhaustive-scale
bound $2^{KB}$. A constructed prefix tree can propagate maxima of the remaining products with those
unexamined bounds. L's whole-configuration prior mixture supplies its separate universal weight bound.

F/N/Q can refine the intervals that make this product large. This is a concrete alternative to plugging
in a single approximate likelihood. Its cost is two bounded continuation queries per bit, including
all inner polynomial, pattern and matrix work. These are conditioned-on-history slot queries;
I/P's unconditional prefix-count queries use related primitives but are not numerically the same events.
Independent guided draws, H's confidence accounting, I's decision bounds and corrected local blocks now
have no unspecified branch probability. There is still no useful uniform concentration guarantee for
the bank without establishing informative interval widths or another proposal bound.

\subsection{A finite moment obstruction says when further queries are unavoidable}
P's LP returns the best bound obtainable from its supplied constraints. It cannot extract information
that those constraints do not contain. For $m<K$, consider counts $c=0,\ldots,m+1$ and the two laws
\[
 p_c^{\rm even}=2^{-m}\binom{m+1}{c}\ind\{c\text{ even}\},\qquad
 p_c^{\rm odd}=2^{-m}\binom{m+1}{c}\ind\{c\text{ odd}\}.
\]
Both normalize to one. The $(m+1)$st finite difference annihilates every polynomial of degree at most
$m$, so their first $m$ factorial moments agree. Nevertheless their tail differences are
\begin{equation}
 \left|P_{\rm even}(A\geq j)-P_{\rm odd}(A\geq j)\right|
 =2^{-m}\binom m{j-1},\qquad 1\leq j\leq m+1.
 \label{eq:moment-information-obstruction}
\end{equation}
The partial alternating binomial identity proves the formula. Thus central count tails can remain
substantially ambiguous even with exact low-order moments.
Set a second leaf count to $K-A$. Both constructions then obey known cardinality and agree on every
mixed moment of total degree at most $m$. A two-leaf prefix-consistency tree alone cannot distinguish them.
Mixing both laws with the same small full-support distribution preserves moment equality and scales
the tail difference, so zero probabilities are not essential to the example.

This is an information obstruction for the supplied moment constraints, not a claim that both laws are
posteriors of this fixed polar bank at some observation. Additional physical likelihood or alias
constraints can narrow the feasible set. Full order recovers each leaf's exact count law; more cheaply,
only enough information to resolve competing allocation gains is required.
If the LP remains wide after numerical moment errors have been reduced, the next step is a higher or
mixed moment, another likelihood/contrast bound, or explicit refinement. Repeatedly optimizing the same
low-order constraints cannot resolve the gap.

\subsection{Distinct sampling has a finite control, but not the IID closure proof}
The primary constructions remain IID. A useful explicit control for a uniform alphabet of size $m$
is that $h$ draws without replacement equal IID draws conditioned on having no repeat. The conditioning
event has prior probability $(m)_h/m^h$, so the two ordered source laws have TV
\[
 \eta_{\rm rep}=1-\frac{(m)_h}{m^h}
 \leq\min\!\left\{1,\binom h2/m\right\},\qquad h\leq m.
\]
Their waveform-sum laws have TV at most this quantity. Consequently their normalized residual likelihood
expectations differ by at most $\eta_{\rm rep}$.
One may add this absolute error to a certified IID residual calculation on the same allowed alphabet,
then use the correct without-replacement source masses in the surrounding count query.
For example, with $K\leq M$, let $\overline D_h^{(-w)}(r)$ be the normalized likelihood expectation for
$h$ distinct draws from the alphabet excluding $w$. The exact word-count numerators are
\[
 N_{w,1}=\frac KM\,\overline D_{K-1}^{(-w)}(y-\phi(w)),\qquad
 N_{w,0}=\left(1-\frac KM\right)\overline D_K^{(-w)}(y).
\]
At $K=M$, the zero-count term is zero and its infeasible residual need not be evaluated.
After fixing distinct outside messages, remove them first. For prefix-dependent nonidentical laws,
the collision probability must be recomputed; the uniform formula is not applicable unchanged.

At the full-posterior level the distinct target is the IID target conditioned on no repeat.
The relevant error is then the \emph{posterior} repeat probability, not $\eta_{\rm rep}$.
Small prior collision probability alone does not justify swapping decoder targets.
This supplies a bounded optional control while keeping the exact distinct recursions and legal sampling
as the general alternatives. No complete without-replacement Gaussian theorem is assumed.

\clearpage
\section{What the combinations establish}
\label{sec:continuing-combinations}
The full exact report suggests several distinct ways to combine the approximation calculations.
The new connections below have explicit error objects; none selects a receiver specification.

\begin{enumerate}
\item \textbf{Residual closure, support geometry and guided or deterministic queries.}
      N uses the exact one-user partition to move F's approximation toward a queried residual and to
      strengthen Report 7's geometric completion bound with prior probabilities.
      The resulting positive estimate can guide L, while its certified intervals can feed I/P.
      A guide need not be certified to define an evaluable full-support proposal; an interval provider does.
      The unresolved tradeoff is tilted-partition access against better concentration and smaller relative errors.
\item \textbf{Moment truncation, prefix discovery and shared normalization.}
      M's inexpensive unrestricted residual moments can be processed by P before buying higher orders.
      I's count hierarchy makes information from different cells mutually constraining; selected mixed
      moments retain some cross-cell dependence. N/F/G can provide the same unnormalized constraints.
      This improves the use of computed information without assuming that different approximation errors
      are independent or that their approximate marginals already form a posterior.
\item \textbf{Exact Gaussian decoupling, positive compression and weighted corrections.}
      O provides another positive representation whose one-user operation is precisely N's partition.
      Its field-domain omission is controlled by posterior mass; its discarded directions are controlled
      by a quadratic correction. A and K therefore apply to different, explicitly identified errors.
      Accurate inner polar sums, outer integration and full-message readout are all still required.
      Common recentering leaves these field directions unchanged, and unweighted field draws can have
      exponentially bad variance. These negative conclusions rule out two easy but unjustified shortcuts.
\item \textbf{Preserve exact local structure before using any of these replacements.}
      Pair/triangle parity quotients and J's joint feature conditioning still reduce local construction
      where their boundaries are affordable. Compatible future projection still precedes B's quantization.
      None of the new formulas turns a low-dimensional interface into a small state table, or turns an
      exact K2/K3 conditional into an approximation that higher users have disappeared.
      Known alias conditionals can improve positive approximations, while count markers preserve their
      irreducible uncertainty. A covariance, a field mode or a moment bound does not identify an alias class.
\end{enumerate}

Connected expansions and marked Fourier sums remain useful alternatives when their computable remainder
is smaller than a moment-closure remainder. N can remove a small common likelihood scale before such a
residual evaluation, but then the Fourier characteristics or connected reference must be recomputed for
the \emph{tilted} source. Unit-modulus Fourier modes do not disappear by assertion, and a signed remainder
does not become a positive proposal. Likewise, exact recentering can improve a connected bound while
leaving O's field correction invariant: the usefulness of a reference change is representation-specific.

The final pass supplies finite answers to the preceding open construction questions.
Q provides character-polynomial tilted and Fourier queries with nested error bounds.
R provides a physical direction rule, dispersion certificates and bounded field integration alternatives.
S connects the common center to a declared variational optimum and an evidence-sandwich certificate.
T permits signed low-rank query calculations with positive propagated remainders.
Section~\ref{sec:guide-moment-completion} derives interval-based guided branches and a moment-information
obstruction. These constructions can now be assessed by their complete cost and error, rather than by
assuming that their inner operations are cheap. None guarantees that its sufficient regime is useful
at the intended load.

\clearpage
\section{Compose positive, nonuniform error envelopes}
\label{sec:envelopes}
The worst-case bounds above are easy to state, but may be controlled by unlikely configurations. There is a
more selective calculation when the approximate contraction can also carry nonnegative error bounds.
On a retained support $E$, suppose
\[
 l(a)\widetilde\omega(a)\leq\omega(a)\leq u(a)\widetilde\omega(a),
 \qquad 0\leq l(a)\leq u(a).
\]
For a count-tail event $A$, calculate lower and upper sums on both it and its complement \emph{within $E$}:
\[
 L_A=\sum_{a\in A\cap E}l(a)\widetilde\omega(a),\quad
 U_A=\sum_{a\in A\cap E}u(a)\widetilde\omega(a),
\]
and similarly $L_{\bar A},U_{\bar A}$. Positivity and monotonicity of $x/(x+y)$ prove
\begin{equation}
 \frac{L_A}{L_A+U_{\bar A}}
 \leq P_y(A\mid E)\leq
 \frac{U_A}{U_A+L_{\bar A}}.
 \label{eq:event-envelopes}
\end{equation}
Use $[0,1]$ when a denominator cannot be bounded positively. Compared with separately bounding evidence,
this preserves that the numerator and its complement form the same total. If omitted-event masses have
upper bounds $U_{D,A}$ and $U_{D,\bar A}$, unconditional bounds follow by adding $U_{D,\bar A}$ to the
left denominator and $U_{D,A}$ to the right numerator and denominator. Setting each to the total $U_D$
is valid but can be loose.

For an exchange contrast $H=\ind_A-\ind_B$, its disjoint positive and negative events are
$A\setminus B$ and $B\setminus A$. If
\[
 L_{A\setminus B}>U_{B\setminus A}+U_{D,B\setminus A},
\]
then $\E_{P_y}H>0$ regardless of normalization and of omitted states favourable to the comparison.
This combines approximate contraction with an exact decision certificate; it does not require accurate
probabilities for every configuration. Multiple exchanges must still certify the final feasible allocation.

For compatible positive modifications, local ratio bounds multiply. Equivalently, certified absolute
log errors add along a path, so $l(a)=e^{-E(a)}$, $u(a)=e^{E(a)}$ are valid. State-specific merging errors
and occupancy-specific omitted-edge errors can enter this way. If the bounds themselves factor on the same
small graph, positive summation propagates them with comparable contraction cost. If obtaining the sums
requires enumeration of all paths or occupancy sectors, this certificate has not saved that work. For
example, $e^{\sum\gamma_{rq}k_rk_q}$ may restore dense occupancy coupling after the physical edge was dropped.
An affordable looser bound may be preferable, but its informativeness must be assessed algebraically.

If only a uniform retained-support log error $E_*$ and omitted-mass bound $\bar\eta$ are available, the
simple composite guarantee remains
\begin{equation}
 R(\widetilde d)-R(d^*)\leq
 \min\!\left\{1,\,2\left[\bar\eta+(1-\bar\eta)\tanh(E_*/2)\right]\right\}.
 \label{eq:composite-error}
\end{equation}
The retained law and every modification must refer to the same original legal configurations; rebuilding
unrelated approximate posteriors at successive cuts does not meet this requirement. These are real-arithmetic
statements. Floating-point rounding would need a separate enclosure if a later implementation claims certification.

Known equal-signal classes provide one further exact correction: keep the approximate mass of each class
but restore the source-prior conditional distribution inside it. Report 7 proves that this cannot increase
TV error. It preserves the irreducible ambiguity instead of assigning an arbitrary message identity.
Finding complete classes may be expensive; known disjoint classes can be corrected without assuming a full census.

\clearpage
\section{Existing illustrative decoder compositions}
\label{sec:decoder-synthesis-stage}
The previous pass assembled the following finite constructions to check that proposal/query calculations,
count readout and error accounting can connect. They are retained as illustrative compositions, not a
commitment to an architecture. Sections~\ref{sec:tilted-residual}--\ref{sec:continuing-combinations}
provide further ingredients and certificates that can be substituted at their stated interfaces.
They share the same encoder, source, physical likelihood and Bayes loss; none has demonstrated an affordable
accuracy regime on the selected bank.

\subsection{Path 1: residual-guided independent importance decoding}
\begin{enumerate}
\item Construct the global polar moments and the conditional-prefix moment operations in F/I. Choose a
      specified positive guide, initially \eqref{eq:guided-likelihood}, and a whole-configuration prior
      mixture weight $\epsilon$. No transmitted pattern or message is supplied.
\item Generate independent ordered tuples with L's bitwise sampler. Evaluate the actual mixture density
      and the original likelihood; use \eqref{eq:guided-mixture} for every importance weight.
\item Accumulate the sparse full-word count tails from H. Optionally also accumulate prefix counts for I.
      Greedily allocate $K$ copies by the empirical tails, preserving multiplicity.
\item At prescribed sample rounds, use H's simultaneous confidence intervals for all possible queries.
      Stop at the desired excess-loss/gain certificate or the declared sample budget. Return the actual
      allocation, numerical failure budget and unresolved bound.
\end{enumerate}
Its direct cost is \eqref{eq:guided-decoder-cost}; no full-message alphabet array is required.
The central favourable condition is useful proposal concentration, alongside affordable moment access.
The guide's Gaussian assumption affects that condition, but is not an uncorrected bias in the target.
Conditional mixtures, exact summation of a small block, or better continuation calculations can improve
the proposal without changing the decoder's readout or target. A poor guide can make the required $S$
prohibitive, even though every individual sample is easy to generate.

An optional variant uses L's global Metropolis proposals mixed with exact or accepted/rejected local
block moves. Retain unweighted count tails from recorded states and use
\eqref{eq:block-chain-certificate}, including all intervening transitions and block contraction costs.
This reuses K2/K3 structure inside a full-$K$ law. Its extra dependence analysis and computation are
justified only if they improve exploration enough; the current conservative bound establishes correctness,
not superiority over the independent version.

\subsection{Path 2: deterministic prefix queries with an explicit gain gap}
\begin{enumerate}
\item Start with a frontier covering all messages and an actual incumbent allocation, for example $K$
      specified distinct words when $M\geq K$. Split its ancestor paths so those words are singleton cells.
\item Query F/I's conditional-moment likelihood intervals for the frontier and incumbent. Tighten them by
      probability and tail monotonicity. Compute I's global optimistic gain $U_*$ and incumbent lower gain.
\item Allocate incumbent copies by the largest resolved-leaf lower tails. Refine an unresolved cell used
      by the optimistic allocation, prioritizing its contribution to the remaining gap. Alternatively
      tighten a limiting query, using conditional groups, a valid G/E remainder, M's moment hierarchy,
      or the finite exact calculation when its cost is acceptable.
\item Keep both children of every split and intersect valid repeated intervals. Stop when the gain gap
      divided by $K$ meets the tolerance, or return the incumbent and current bound at the resource limit.
\end{enumerate}
Taking the best incumbent lower gain and smallest valid global upper gain makes the reported gap
nonincreasing. A concrete budgeted rule can break equal priorities lexicographically and cap each provider's
work; no favourable termination claim depends on that choice. Fully refining the tree and using exact
queries is a finite endpoint if a certified answer is required irrespective of cost.

The base query cost is up to $O(Kn^3)$ per visited cell after moments, plus the frontier/heap work in I.
Every repeated query and conditional-mixture sector counts. Useful scaling requires both informative
likelihood intervals relative to evidence and a small set of competitive prefixes. When the bounds remain
wide, the algorithm returns a weak certificate or incurs extensive refinement. It never reports an
unresolved prefix as a decoded word. M can simplify resolved-word multiplicity queries when posterior
collisions are controlled; it does not automatically solve discovery.

\subsection{Path 3: a positive structured contraction, with optional correction}
\begin{enumerate}
\item Use the whole alphabet, or a candidate alphabet whose omitted posterior mass is bounded by A.
      Partition it into disjoint regions and construct legal local count states, source weights and signals.
      Any occupancy caps also require omitted-mass accounting.
\item Choose exact reference changes first. Retain selected residual interactions; approximate appropriate
      factors with D's positive features. Choose a forest, J's feature/state conditioning schedule, or a
      compatible B state graph, paying for the actual constructed joint states.
\item Contract with the total-$K$ constraint and full-message marks intact. Compute count tails from the
      same normalized positive law, then greedily allocate actual copies.
\item Bound the retained-support distortion with local ranges, positive envelopes or K's justified
      weighted correction bounds. Combine it with omitted mass. Refine factors/states when affordable,
      or return the current action and bound.
\end{enumerate}
For the C/D/J version, if $E_{\rm cut}$ bounds the absolute log error of omitted edges, a valid range is
\[
 \Omega=2E_{\rm cut}+\sum_{\text{feature edges }e}
                  [2\epsilon_e-\log p_{L_e}(\Lambda_e)].
\]
Additional B errors add their ranges only if their pathwise bounds hold for this same final surrogate.
Conditioning in J adds work but no further error.
The direct excess-loss bound is
\[
 \min\{1,\,2[\bar\eta+(1-\bar\eta)\tanh(\Omega/4)]\}.
\]
Forests use \eqref{eq:feature-cost}; cyclic schedules use \eqref{eq:cycle-cost}; the B alternative uses its
generated-state/transition accounting. Candidate discovery, local-state construction and certificates
remain included. An arbitrary restricted configuration event can introduce extra coupling and must be
retained explicitly; it cannot be replaced by independent region alphabets.

Once this law $Q$ can be sampled and evaluated, H can correct it by exact importance weights, or K can
estimate corrections to already-known query values. On retained support $\gamma\leq e^\Omega$ is available.
This offers a deliberate tradeoff: reduce the contraction effort until correction cost outweighs the
saving. Source-prior mixing can restore support, at the cost of evaluating the actual mixture.
Unlike Path 1, this path generally has explicit alphabet/local-state preprocessing. Its potential advantage
is a structured, jointly consistent posterior and a controllable correction; affordable construction is
its principal unresolved condition.

\subsection{Combination rules that prevent an error or cost from being hidden}
These paths permit combinations, but their guarantees are not interchangeable.
\begin{enumerate}
\item \textbf{Restriction is separate.} If only $P_y(\cdot\mid E)$ is approximated and its tail error is $e$,
      unconditional error is at most $\bar\eta+(1-\bar\eta)e$. Conditional intervals $[L,U]$ become valid
      unconditional intervals $[(1-\bar\eta)L,\,(1-\bar\eta)U+\bar\eta]$.
\item \textbf{A corrected proposal does not retain its old deterministic bias.} Exact importance or
      Metropolis correction targets the original supported law. Surrogate error then controls weights or
      mixing, while the reported decision error uses finite-sampling bounds. Adding both the full surrogate
      bias and the corrected sampling error would be unnecessarily pessimistic.
\item \textbf{Direct query approximations need truthful intervals, not a fictitious joint law.} F/G/M can
      supply separate intervals to I. Intersect intervals for the same exact query; random providers need
      simultaneous confidence across queries and rounds. Independently derived providers need not be
      probabilistically independent for a union bound.
\item \textbf{All likelihood dependence is used through one target.} Conditioning a block changes its
      residual; exact reweighting uses the original likelihood once in its target weight. Neither operation
      adds another observation or treats an inference residual as an independent measurement.
\item \textbf{A useful certificate can still be expensive.} Evaluating correction moments, rare bad-region
      mass, feature cutsets, alias classes or proposal densities must not secretly solve the original
      exponential problem. A loose explicit bound is preferable to an uncharged oracle.
\end{enumerate}
If a randomized guarantee says the conditional excess risk is at most $\rho$ except with probability
$\delta$, its expectation over numerical randomness is at most $\rho+\delta$, capped at one.
This does not establish a favourable source/noise-averaged PUPE value; absolute posterior risk and later
controlled performance evaluation remain distinct.

\subsection{What the combinations have and have not settled}
The exact resources now have concrete roles. Conditional polar characters build guides and query moments;
joint overlap parities reduce local construction before approximation; K2/K3 can supply exact conditional
block draws; compatible future projections precede state quantization; support directions and adverse prior
masses certify discovery; alias conditionals preserve irreducible ambiguity; unexcluded residuals provide
factorial moments with bounded reconstruction. The same encoder therefore supports several complementary
computational reductions, with count identity retained throughout.

S supplies one specified variational product family and an evidence-gap certificate; C's zero-edge case
already includes a cardinality-constrained product over regional states. Other local message projections
must still specify what joint factor is replaced and how cardinality/count queries survive. A bitwise or loopy
message calculation could provide a guide in L, but its name alone supplies neither a count-law guarantee
nor a mixing bound. E's connected coefficients and G's signed sums remain alternative query/evidence
providers where their access cost and remainder are favourable. They need not be forced into a positive
joint proposal. This places further ideas at their actual computational role without an assumption atlas.

\clearpage
\section{Completion audit before implementation}
\label{sec:completion-audit}
\subsection{What is complete, and what that statement means}
The concrete theoretical prerequisites for the retained paths are now in place. Each has an exact
origin, a finite replacement or query calculation, a source/count contract, an error statement and an
explicit account of the potentially expensive work. Where an approximation has an irreducible limitation,
the report gives a counterexample or a condition that would have to be established. \emph{Complete} here
means that an implementation need not invent a missing statistical target or invoke an unspecified
posterior/partition oracle. It does not mean that the bounds are tight, that the required regime holds
for this encoder, or that the calculations are affordable.

The following is an audit of the existing derivations, not a new assumption menu. Some limits recover
exact inference; others give arbitrary accuracy with an explicit tolerance or failure probability.
A Gaussian closure by itself does not become exact merely by evaluating its formula more accurately.

\begin{center}\small
\begin{tabularx}{\textwidth}{@{}p{.12\textwidth}YY@{}}\toprule
Entries & Closed calculation and refinement limit & Remaining usefulness or cost question\\\midrule
A, I & Omitted mass and optimistic prefix gains bound unresolved decisions; full refinement restores the exact alphabet. &
How many prefixes and likelihood envelopes are needed before a decision separates?\\
B & Compatible state quantization has a positive sequential construction and occupied-transition error bound; zero rounding restores exact states. &
State/key construction, alphabet traversal and covering dimension can remain exponential.\\
C, D, J & Regional counts, positive features and joint cyclic cutsets preserve cardinality and markers; restoring interactions/features gives exact contraction. &
Local state counts, feature construction and joint separator products determine cost.\\
E, S & A fixed centered reference, positive edge floor and connected-component restriction have explicit remainders; order $K$ restores the expansion. &
Center optimization need not reduce the connected remainder; the remainder itself may be loose.\\
F, N, Q & Excluded/conditional residuals have smoothing, tilt and inner-series bounds. Q adds a convergent likelihood-polynomial hierarchy when Gaussian closure fails. &
Moment tensors, degree, pattern sums and small normalized likelihood intervals can dominate.\\
G, Q, T & Fourier and signed low-rank queries carry explicit inner, sampling or contraction remainders; full modes/ranks or increasing accuracy recover their exact queries. &
Coefficient access, cancellation, marked envelopes and evidence normalization can remove the gain.\\
H, K, L & Evaluable full-support proposals, exact correction, finite confidence and conservative block-chain bounds are specified. Robust bit intervals define another proposal. &
Importance concentration, certificate cost, mixing and all-prefix guarantees can be prohibitive.\\
M, P & Factorial moments, common evidence and prefix consistency give finite tail/gain programs; full-order leaf moments determine the required tails. &
Low-order ambiguity can survive exact moments; new physical or higher-order information must then be purchased.\\
O, R & Positive fields have controlled direction/domain reduction and bounded integration routes; Q supplies inner sums. Restoring directions and integration accuracy recovers the target. &
Field dimension, determinant/variance, inner-query degree and full-word discovery remain charged.\\
S & A strictly concave identical-slot objective defines a unique reference; an evidence sandwich certifies its posterior error when informative. &
The best member of this family can still be inaccurate; exact correction or richer interactions may be needed.\\
\bottomrule
\end{tabularx}
\end{center}

The distinct-message control remains separately conditioned. The finite collision bound supplies one
optional comparison of residual source laws, while the exact distinct recursions remain available.
No IID product, Gaussian-field factorization or variational formula is silently reused as a distinct law.

\subsection{One final error account for a composition}
Keep three operations separate: restricting support, approximating the retained target, and numerically
evaluating its queries. Suppose a retained event $E$ has omitted posterior mass at most $\eta$,
the normalized positive approximation to $P_y(\,\cdot\mid E)$ has log-weight range error at most $\Omega$,
and every returned conditional tail, clipped to $[0,1]$, has numerical error at most $e_{\rm num}$
relative to that approximate law. On the event where the numerical bounds hold, a valid total tail error is
\begin{equation}
 e_{\rm total}\leq
 \eta+(1-\eta)\min\{1,\tanh(\Omega/4)+e_{\rm num}\}.
 \label{eq:completion-error-ledger}
\end{equation}
This follows by comparing conditional tails first, then writing the exact posterior as the mixture
on $E$ and $E^c$. Compatible log-weight errors can be added before applying the hyperbolic tangent.
For an action maximizing the estimated gains as in Section~\ref{sec:decision-foundation}, the excess
multiset risk is at most $\min(1,2e_{\rm total})$; a sharper direct decision-contrast or prefix-gain
certificate may replace it. Numerical failure probabilities are combined separately.

If the approximate law is used only as a proposal and exact importance correction targets
$P_y(\,\cdot\mid E)$, its $\Omega$ is not also a posterior-bias term. The remaining error is
$\eta+(1-\eta)e_{\rm num}$, with the proposal affecting concentration and cost instead.
Conversely, approximate physical likelihoods used in the correction weights, or approximate inner
partitions used as target field weights, do introduce target error and need their own bounds.
Correcting with an unavailable or incorrectly marginalized proposal density is not exact correction.

Interval providers need not describe one common approximate posterior. They may instead constrain the
same exact numerators and shared evidence in P. Intersecting their valid constraints is allowed without
independence assumptions; turning unrelated approximations into artificial joint marginals is not.
For weighted or signed corrections, the construction of their moments/envelopes is part of the cost.
The same observation is used once in the physical target throughout these changes of representation.

\subsection{The next questions are specific enough to implement}
This closes the broad preparatory pass for the paths above. It is not a theorem that every useful
approximation has been found. Further algebra should now answer a specific missing bound or observed
bottleneck, rather than be a prerequisite to beginning any receiver.

The next stage can specify and implement a small exact-reference comparison of selected compositions.
It should distinguish three things: actual query/decision error, the looseness or failure probability
of its certificate, and the total cost including construction and discovery. A good estimate with a
vacuous bound, a strong bound obtained by enumeration, and a useful inexpensive certified estimate are
different outcomes. The guided and deterministic-prefix paths remain natural first comparisons, while
positive contractions and fields remain alternatives; this report has not selected an empirical winner.

No receiver implementation, new mathematical test, simulation, tuning or PUPE ranking accompanies this
pass. No useful polynomial-in-$B,K$ guarantee is established for the selected bank. Report 7 retains the
exact reference; \path{docs/POLAR_APPROXIMATION_LOG.md} records the deductions and cautions, and
\path{docs/CURRENT_STATE.md} owns the active stage.

\clearpage
\begin{thebibliography}{9}
\bibitem{mengersen-tweedie} K. L. Mengersen and R. L. Tweedie,
\emph{Rates of Convergence of the Hastings and Metropolis Algorithms},
The Annals of Statistics 24(1), 101--121, 1996.
\href{https://www2.stat.duke.edu/~scs/Courses/Stat376/Papers/ConvergeRates/GeneralState/MengersenTweedieAnnStat1996.pdf}
{Paper}. Cited for the classical independence-sampler convergence principle. The finite-state minorization,
count-query coupling bound and decoder cost accounting are derived here for the fixed model; no general
approximation technique is claimed as novel.
\bibitem{daniels} H. E. Daniels,
\emph{Saddlepoint Approximations in Statistics},
The Annals of Mathematical Statistics 25(4), 631--650, 1954.
\href{https://people.eecs.berkeley.edu/~jordan/sail/readings/archive/daniels-1954.pdf}{Paper}.
Cited for the classical tilting/saddlepoint context. Entry N derives its own finite Gaussian-smoothed
identity and error bounds; it assumes no asymptotic saddlepoint accuracy theorem.
\bibitem{bertsimas-popescu} D. Bertsimas and I. Popescu,
\emph{Optimal Inequalities in Probability Theory: A Convex Optimization Approach},
SIAM Journal on Optimization 15(3), 780--804, 2005.
\href{https://www.mit.edu/~dbertsim/papers/MomentProblems/Optimal-inequalities-in-probability-theory-A-convex-optimization-approach-SIAM15.pdf}{Paper}.
Cited for optimization over moment information. Entry P uses the elementary finite-support LP and
derives its count/prefix constraints directly; no continuous-support semidefinite result is needed.
\bibitem{wainwright-jordan} M. J. Wainwright and M. I. Jordan,
\emph{Graphical Models, Exponential Families, and Variational Inference},
Foundations and Trends in Machine Learning 1(1--2), 1--305, 2008.
\href{https://www.cs.princeton.edu/courses/archive/fall11/cos597C/reading/WainwrightJordan2008.pdf}{Paper}.
Cited for the variational/exponential-family context. Entry S derives the identical-slot optimum,
its limitations and its evidence-gap certificate for this particular source and loss.
\bibitem{oseledets} I. V. Oseledets,
\emph{Tensor-Train Decomposition}, SIAM Journal on Scientific Computing 33(5), 2295--2317, 2011.
\href{https://epubs.siam.org/doi/10.1137/090752286}{Paper}.
Cited as context for tensor compression. Entry T derives finite marked-contraction error transport;
no low-rank existence, cheap construction or posterior guarantee is imported from a generic decomposition.
\end{thebibliography}
\end{document}
